\documentclass[11pt]{aghdpl}
% \documentclass[en,11pt]{aghdpl}  % praca w języku angielskim

% Lista wszystkich języków stanowiących języki pozycji bibliograficznych użytych w pracy.
% (Zgodnie z zasadami tworzenia bibliografii każda pozycja powinna zostać utworzona zgodnie z zasadami języka, w którym dana publikacja została napisana.)
\usepackage[english,polish]{babel}

% Użyj polskiego łamania wyrazów (zamiast domyślnego angielskiego).
\usepackage{polski}

\usepackage[utf8]{inputenc}

% dodatkowe pakiety

\usepackage{float}
\usepackage{mathtools}
\usepackage{amsfonts}
\usepackage{amsmath}
\usepackage{amsthm}
\usepackage{siunitx}
%\usepackage{caption}

% \usepackage{tabularx}

% Dodane biblioteki
\usepackage{changepage} % To adjust the width of the column for the title part and figures/tables (adjustwidth environment)
\usepackage{hyperref}

% --- < bibliografia > ---

\usepackage[
style=numeric,
sorting=none,
%
% Zastosuj styl wpisu bibliograficznego właściwy językowi publikacji.
language=autobib,
autolang=other,
% Zapisuj datę dostępu do strony WWW w formacie RRRR-MM-DD.
urldate=edtf,
seconds=true,
% Nie dodawaj numerów stron, na których występuje cytowanie.
backref=false,
% Podawaj ISBN.
isbn=true,
% Nie podawaj URL-i, o ile nie jest to konieczne.
url=false,
%
% Ustawienia związane z polskimi normami dla bibliografii.
maxbibnames=3,
% Jeżeli używamy BibTeXa:
backend=bibtex
]{biblatex}

\usepackage{csquotes}
% Ponieważ `csquotes` nie posiada polskiego stylu, można skorzystać z mocno zbliżonego stylu chorwackiego.
\DeclareQuoteAlias{croatian}{polish}

\addbibresource{bibliografia.bib}

% Nie wyświetlaj wybranych pól.
%\AtEveryBibitem{\clearfield{note}}


% ------------------------
% --- < listingi > ---

% Użyj czcionki kroju Courier.
\usepackage{courier}

\usepackage{listings}
\lstloadlanguages{TeX}

\lstset{
	literate={ą}{{\k{a}}}1
           {ć}{{\'c}}1
           {ę}{{\k{e}}}1
           {ó}{{\'o}}1
           {ń}{{\'n}}1
           {ł}{{\l{}}}1
           {ś}{{\'s}}1
           {ź}{{\'z}}1
           {ż}{{\.z}}1
           {Ą}{{\k{A}}}1
           {Ć}{{\'C}}1
           {Ę}{{\k{E}}}1
           {Ó}{{\'O}}1
           {Ń}{{\'N}}1
           {Ł}{{\L{}}}1
           {Ś}{{\'S}}1
           {Ź}{{\'Z}}1
           {Ż}{{\.Z}}1,
	basicstyle=\footnotesize\ttfamily,
}

% ------------------------

\AtBeginDocument{
	\renewcommand{\tablename}{Tabela}
	\renewcommand{\figurename}{Rys.}
}

% ------------------------
% --- < tabele > ---

\usepackage{array}
\usepackage{tabularx}
\usepackage{multirow}
\usepackage{booktabs}
\usepackage{makecell}
\usepackage[flushleft]{threeparttable}

% defines the X column to use m (\parbox[c]) instead of p (`parbox[t]`)
\newcolumntype{C}[1]{>{\hsize=#1\hsize\centering\arraybackslash}X}


%---------------------------------------------------------------------------

\author{Roman Nowak}
\shortauthor{Roman Nowak}

%\titlePL{Przygotowanie bardzo długiej i pasjonującej pracy dyplomowej w~systemie~\LaTeX}
%\titleEN{Preparation of a very long and fascinating bachelor or master thesis in \LaTeX}

\titlePL{Estymacja stanu robota mobilnego z wykorzystaniem danych zdarzeniowych}
\titleEN{State estimation of a mobile robot using event data}


\shorttitlePL{Estymacja stanu robota} % skrócona wersja tytułu jeśli jest bardzo długi
\shorttitleEN{Event-based state estimation}


% Dopuszczalne wartości[1,2]:
% * "Projekt dyplomowy" - na koniec studiów I stopnia
% * "Praca dyplomowa" - na koniec studiów II stopnia
% [1] Zasady dyplomowania w roku akademickim 2020/2021 (Decyzja Dziekana WEAIiIB nr 16/2020 z dnia 9 grudnia 2020 roku)
% [2] Załącznik nr 1a) do Decyzji nr 16/2020 Dziekana Wydziału EAIiIB z dnia 09 grudnia 2020 r.
\thesistype{Praca dyplomowa}
%\thesistype{Master of Science Thesis}

\supervisor{dr inż. Tomasz Kryjak}
%\supervisor{Paweł Kłeczek, PhD}

\degreeprogramme{Automatyka i Robotyka}
%\degreeprogramme{Automatics and Robotics}

\date{2026}

%\department{Katedra Informatyki Stosowanej}
%\department{Department of Applied Computer Science}

\faculty{Wydział Elektrotechniki, Automatyki, Informatyki i Inżynierii Biomedycznej}
%\faculty{Faculty of Electrical Engineering, Automatics, Computer Science and Biomedical Engineering}

%\acknowledgements{Serdecznie dziękuję \dots tu ciąg dalszych podziękowań np. dla promotora, żony, sąsiada itp.}

\setlength{\cftsecnumwidth}{10mm}

%---------------------------------------------------------------------------
\setcounter{secnumdepth}{4}
\brokenpenalty=10000\relax

% rubber: bibtex.path ./


\newenvironment{abstractpage}
{\cleardoublepage\vspace*{\fill}\thispagestyle{empty}}
{\vfill\cleardoublepage}
\renewenvironment{abstract}[1]
{\bigskip\selectlanguage{#1}%
    \begin{center}\bfseries\abstractname\end{center}}
{\par\bigskip}
\begin{document}

    \titlepages
    % Ponowne zdefiniowanie stylu `plain`, aby usunąć numer strony z pierwszej strony spisu treści i poszczególnych rozdziałów.
    \fancypagestyle{plain}
    {
    	% Usuń nagłówek i stopkę
    	\fancyhf{}
    	% Usuń linie.
    	\renewcommand{\headrulewidth}{0pt}
    	\renewcommand{\footrulewidth}{0pt}
    }

    \begin{abstractpage}
\begin{abstract}{polish}

Kamery zdarzeniowe rejestrują lokalne zmiany jasności z~wysoką rozdzielczością czasową i~dużą rozpiętością tonalną, dzięki czemu są szczególnie przydatne podczas szybkiego ruchu oraz w~trudnych warunkach oświetleniowych. W~niniejszej pracy zaprezentowano system jednoczesnej lokalizacji i~mapowania (SLAM), wyznaczający trajektorię na podstawie danych z~pary kamer zdarzeniowych w~układzie stereo. Zdarzenia są synchronizowane i~agregowane do wykładniczo zanikających reprezentacji obrazowych. Punkty wykrywane metodą GFTT są śledzone metodą Lucas--Kanade, a~ich głębokość jest wyznaczana przez triangulację stereo. Poza jest estymowane przez PnP z~algorytmem RANSAC.  System korzysta z~danych z~żyroskopu, aby kompensować rozmycie wynikające z~ruchu obrotowego oraz wyznaczać zmianę orientacji. W~celu ograniczania dryfu, poza jest estymowana względem punktów rzadkiej mapy tworzonej z~klatek kluczowych, a~rozpoznawanie miejsc umożliwia domykanie pętli i~optymalizację grafu trajektorii. System oceniono na danych ze zbiorów EvSLAM~\cite{evslam2025challenge} i~M3ED~\cite{m3ed2025challenge}, o~rozdzielczościach odpowiednio $640\times480$ i~$1280\times720$ pikseli. Dla siedmiu sekwencji EvSLAM uzyskano sumę ATE równą $3{,}654\,\mathrm{m}$ i~sumę AUC równą $4{,}996$, wobec $0{,}817\,\mathrm{m}$ i~$6{,}286$ dla zwycięskiego rozwiązania. Dla sekwencji M3ED z~dostępną trajektorią referencyjną wyniki różnią się w zależności od charakterystyki sekwencji -- dla dwóch badanych wyniosło odpowiednio $0{,}162\,\mathrm{m}$ i~$25{,}305\,\mathrm{m}$.

\end{abstract}

\noindent\textbf{Słowa kluczowe:} kamery zdarzeniowe, stereowizja, jednoczesna lokalizacja i~mapowanie, odometria wizyjna, estymacja pozy

\end{abstractpage}

% \newpage

\begin{abstractpage}
\begin{abstract}{english}

Event cameras register local brightness changes with high temporal resolution and a~wide dynamic range, making them particularly useful during rapid motion and under challenging illumination. This thesis presents a~simultaneous localization and mapping (SLAM) system that estimates a~trajectory from a~stereo pair of event cameras. Events are synchronized and aggregated into exponentially decaying image representations. Points detected with GFTT are tracked using the Lucas--Kanade method, and their depth is recovered through stereo triangulation. The resulting 3D--2D correspondences are used for PnP estimation with RANSAC. The system uses gyroscope measurements to~compensate for the smearing of~structures caused by rotational motion and to determine orientation changes. To reduce drift, the pose is estimated relative to a~sparse map constructed from keyframes, while place recognition enables loop closure and pose graph optimization. The system was evaluated on data from the EvSLAM~\cite{evslam2025challenge} and M3ED~\cite{m3ed2025challenge} datasets, with resolutions of~$640\times480$ and $1280\times720$ pixels, respectively. Across seven EvSLAM sequences, the system obtained a~total ATE of~$3.654\,\mathrm{m}$ and a~total AUC of~$4.996$, compared with $0.817\,\mathrm{m}$ and $6.286$ for the winning entry. For the M3ED sequences with available ground truth, the results varied depending on the sequence characteristics: the ATE values for the two evaluated sequences were $0.162\,\mathrm{m}$ and $25.305\,\mathrm{m}$, respectively.

\end{abstract}

\noindent\textbf{Keywords:} event cameras, stereo vision, simultaneous localization and mapping, SLAM, visual odometry, VO, pose estimation

\end{abstractpage}

    \setcounter{tocdepth}{2}
    \tableofcontents
    \clearpage
    
    \chapter{Wprowadzenie}
\label{cha:wprowadzenie}

Estymacja stanu jest jednym z~podstawowych zadań systemu percepcji robota mobilnego. Informacja o~pozycji, orientacji i~prędkości jest potrzebna między innymi do planowania ruchu, nawigacji, stabilizacji platformy oraz tworzenia mapy otoczenia. Pozwala ona określić, gdzie znajduje się robot, w~jakim kierunku jest zwrócony i~jak zmienia się jego położenie. Bez dostatecznie dokładnego oszacowania tych wielkości nie jest możliwe wiarygodne odnoszenie obserwacji czujników do otoczenia ani bezpieczne wykonywanie zaplanowanego ruchu.

Położenie robota może być wyznaczane przez zewnętrzne systemy lokalizacji, takie jak globalne systemy nawigacji satelitarnej (GNSS, ang. \textit{Global Navigation Satellite System}), lub infrastrukturę rozmieszczoną w środowisku. Rozwiązania te nie zawsze są jednak dostępne. Sygnał GNSS może zanikać wewnątrz budynków, pomiędzy wysokimi obiektami lub pod ziemią, a jego precyzja jest ograniczona, natomiast specjalna infrastruktura pozwala na ruch robota tylko w ściśle określonym obszarze. W~takich warunkach ruch można wyznaczać na podstawie obserwacji z~czujników pokładowych. Odometria pozwala śledzić kolejne zmiany położenia, a~system SLAM dodatkowo tworzy mapę i~wykorzystuje ją do lokalizacji~\cite{cadena2016slam}. Takie rozwiązanie może uzupełniać zewnętrzne źródło pozycji, czasowo je zastępować lub stanowić samodzielny mechanizm estymacji położenia i prędkości.

Kamery są atrakcyjnym źródłem danych dla pokładowego systemu lokalizacji, ponieważ dostarczają informacji o~strukturze obserwowanej sceny bez konieczności modyfikowania otoczenia. Ruch kamery można wyznaczać na podstawie zmian położenia tych samych elementów sceny w~kolejnych obserwacjach. Jakość takiej estymacji zależy jednak od czytelności obrazu. Szybki ruch może powodować jego rozmycie, natomiast duże różnice oświetlenia utrudniają jednoczesną obserwację jasnych i~ciemnych fragmentów sceny.

Kamery zdarzeniowe rejestrują zmiany jasności niezależnie dla poszczególnych pikseli. Wysoka rozdzielczość czasowa, małe opóźnienie i~duża rozpiętość tonalna czynią je interesującym źródłem danych podczas szybkiego ruchu oraz przy trudnym oświetleniu~\cite{gallego2022event}. Ich wyjściem jest jednak asynchroniczny strumień zdarzeń, a~nie gotowe obrazy. Opracowanie systemu lokalizacji wymaga zatem dobrania reprezentacji danych i~metod estymacji ruchu odpowiednich dla tego formatu.

Zastosowanie pary kamer pozwala wyznaczać głębokość obserwowanych punktów i~skalę metryczną ruchu na podstawie geometrii stereo~\cite{hartley2004multiple}. Informacja o~prędkości kątowej z~żyroskopu dostarcza natomiast niezależnej obserwacji ruchu obrotowego i~może wspierać przetwarzanie podczas szybkich zmian orientacji~\cite{kneip2011gyro}. Połączenie tych czujników umożliwia wykorzystanie zalet kamer zdarzeniowych, a~jednocześnie dostarcza informacji potrzebnych do wyznaczania przestrzennego ruchu robota.

Opracowany system realizuje zdanie estymacji stanu przez obliczanie trajektorii położenia i prędkości. Może stanowić źródło informacji o~ruchu dla dalszych elementów systemu robota, takich jak nawigacja lub planowanie trajektorii, szczególnie w~warunkach ograniczonej dostępności zewnętrznych metod lokalizacji.

\section{Cel i zakres projektu}
\label{sec:cel_i_zakres}

Głównym celem projektu jest zaprojektowanie, zaimplementowanie i~przygotowanie do ewaluacji systemu wyznaczającego ruch robota na podstawie dwóch kamer zdarzeniowych, zdolnego do przetworzenia danych z publicznie dostępnych zbiorów danych, udostępnionych w ramach międzynarodowych konkursów badawczych: CVPR 2025 Workshop on Event-based Vision -- SLAM Challenge \cite{evslam2025challenge} (dalej opisywanym jako evSLAM) oraz IROS 2025 EvSLAM Challenge -- A Challenge on Event-based State Estimation for High-Speed Maneuvers \cite{m3ed2025challenge} (dalej opisywanym jako m3ed Challenge). Powstałe rozwiązanie realizuje odometrię wizyjną (ang. \textit{visual odometry}, VO), rozszerzoną o mechanizmy rzadkiego systemu jednoczesnej lokalizacji i~mapowania (ang. \textit{simultaneous localization and mapping}, SLAM).

Jako pierwszy etap pracy zrealizowany został przegląd literatury, ze szczególnym uwzględnieniem kamer zdarzeniowych, estymacji pozycji, metod VO oraz SLAM dostosowanych do danych asynchronicznych. Pozwoliło to na wybranie sposobu rozwiązania problemu. Drugim etapem była implementacja algorytmu.

Wymagane funkcjonalności systemu:
\begin{itemize}
    \item odczyt, kalibracja i~synchronizacja strumieni zdarzeń;
    \item utworzenie obrazowych reprezentacji zdarzeń oraz rektyfikację pary stereo;
    \item detekcja i~śledzenie cech, wyznaczanie głębokości oraz estymacja pozy kamery;
    \item opcjonalne wykorzystanie danych z IMU (ang. \textit{inertial measurment unit}) do wspomagania działania systemu;
    \item utrzymywanie mapy otoczenia, rozpoznawanie miejsc, domykanie pętli i~optymalizacja grafu pozy;
    \item zapis trajektorii, prędkości wyznaczonej z~kolejnych póz oraz mapy w~formatach umożliwiających ocenę na obu zbiorach.
\end{itemize}

System nie wykorzystuje kamery klatkowej, mimo że taki czujnik jest dostępny w ramach przyjętych zbiorów danych. Celem pracy jest skupienie się na kamerach zdarzeniowych i próba wykorzystania ich możliwości.

Trzecim i ostatnim etapem jest zebranie wyników oraz ich ewaluacja. Wydajność rozwiązania zostala oceniona przy użyciu standardowych metryk takich jak: ATE (ang. absolute trajectory error), RMSE (ang. root mean squared error) oraz RPE (ang. relative pose error). Wyniki zostały przedstawione na wykresach umożliwiających wizualną ocenę jakości działania systemu.

\section{Zawartość pracy}
\label{sec:zawartosc_pracy}

W~rozdziale \ref{cha:podstawy_teoretyczne} wyjaśniono podstawowe pojęcia używane w~tekście pracy i~metody wykorzystywane w implementacji systemu. W~rozdziale~\ref{cha:przeglad_literatury} omówiono istniejące rozwiązania VO, VIO i~SLAM. Rozdział~\ref{cha:projekt_systemu} opisuje założenia i~architekturę opracowanego systemu, natomiast rozdział ~\ref{cha:implementacja} jego realizację programową. Metodykę eksperymentów określono w~rozdziale~\ref{cha:metodyka_badan}, ich wyniki przedstawiono w~rozdziale~\ref{cha:wyniki}, a~wnioski i~ograniczenia rozwiązania zebrano w~rozdziale~\ref{cha:podsumowanie}.

    \chapter{Podstawy teoretyczne}
\label{cha:podstawy_teoretyczne}

W~rozdziale przedstawiono pojęcia potrzebne do opisu estymacji ruchu na podstawie danych zdarzeniowych. Wyjaśniono zasadę działania kamery zdarzeniowej, podstawowe reprezentacje jej danych, zapis pozy oraz różnice między VO, SLAM, i~VIO oraz przedstawiono sposoby wykorzystywania danych wizyjnych do estymacji ruchu.  Następnie omówiono geometrię stereo, estymację pozy i~grafową reprezentację trajektorii.

\section{Kamery zdarzeniowe}
\label{sec:kamery_zdarzeniowe}

Kamera klatkowa rejestruje jednocześnie natężenie światła dla wszystkich pikseli, a~kolejne obrazy powstają w~chwilach narzuconych przez częstotliwość pracy czujnika. Kamera zdarzeniowa działa inaczej. Każdy jej piksel reaguje niezależnie na zmianę jasności, dlatego wyjściem nie jest sekwencja pełnych ramek obrazów, lecz asynchroniczny strumień danych nazywanych zdarzeniami. Różnicę między tymi sposobami rejestracji przedstawiono na rysunku~\ref{fig:event_camera_comparison}.

\begin{figure}[!ht]
    \centering
    \includegraphics[width=0.78\linewidth]{theory_frame_camera_vs_event_camera.png}
    \caption{Porównanie wyjścia kamery klatkowej i~zdarzeniowej podczas obserwacji obracającej się tarczy. Źródło:~\cite{rebecq2018esim}.}
    \label{fig:event_camera_comparison}
\end{figure}

Taki sposób pomiaru zapewnia małe opóźnienie i~dużą rozdzielczość czasową. Niezależna praca pikseli pozwala również uzyskać dużą rozpiętość tonalną i~może ograniczać rozmycie ruchu typowe dla skończonego czasu ekspozycji kamery klatkowej. Właściwości te są przydatne podczas szybkiego ruchu robota oraz przy dużych zmianach oświetlenia~\cite{gallego2022event}. Pojedyncze zdarzenie nie zawiera jednak bezwzględnej wartości jasności, a~liczba pomiarów zależy od kontrastu sceny i~ruchu. W~scenie statycznej strumień może być bardzo rzadki, natomiast zmiany oświetlenia i~szum czujnika mogą powodować zdarzenia niezwiązane z~ruchem obserwowanych krawędzi. Działanie systemu opartego wyłącznie o kamery zdarzeniowe, jest więc uzależnione od występowania na scenie wystarczającej liczby zmian jasności.

\subsection{Model generowania zdarzeń}
\label{subsec:model_zdarzen}

Niech $I(\mathbf{x},t)$ oznacza natężenie światła mierzone w~pikselu $\mathbf{x}=(x,y)^T$ w~chwili $t$. Logarytm jasności $L$ definiuje równanie
\begin{equation}
    L(\mathbf{x},t)=\log I(\mathbf{x},t).
    \label{eq:log_intensity}
\end{equation}
Zdarzenie jest generowane, gdy zmiana $L$ względem wartości zapamiętanej podczas poprzedniego zdarzenia w~tym samym pikselu osiągnie próg kontrastu $C$~\cite{gallego2022event,rebecq2018esim}:
\begin{equation}
    L(\mathbf{x}_k,t_k)-L(\mathbf{x}_k,t_k-\Delta t_k)=p_k C.
    \label{eq:event_generation}
\end{equation}
Indeks $k$ oznacza numer zdarzenia, $t_k$ jego znacznik czasu, $\Delta t_k$ czas od poprzedniego zdarzenia w~pikselu $\mathbf{x}_k$, a~$p_k\in\{-1,1\}$ polaryzację. Wartość dodatnia oznacza wzrost, a~ujemna spadek jasności. Pojedynczy pomiar ma postać
\begin{equation}
    e_k=(x_k,y_k,t_k,p_k),
    \label{eq:event_tuple}
\end{equation}
gdzie $x_k$ i~$y_k$ są współrzędnymi piksela. Znacznik $t_k$ opisuje chwilę wystąpienia zmiany, a~nie czas wspólny dla całej matrycy. Zdarzenia pochodzące z~różnych pikseli nie tworzą więc naturalnych klatek.

Równanie~\eqref{eq:event_generation} jest modelem idealnym. W~rzeczywistym czujniku dodatni i~ujemny próg mogą się różnić, a~ich wartości nie są identyczne dla wszystkich pikseli. Na pomiar wpływają również inne niedoskonałości (\cite{gallego2022event}):

\begin{itemize}
    \item \textbf{szum śrutowy} (ang. \textit{shot noise}) -- wynika z~losowego
    charakteru liczby fotonów docierających do fotodiody. Powoduje fluktuacje
    mierzonego sygnału, które mogą prowadzić do wygenerowania zdarzenia mimo
    braku odpowiadającej mu zmiany jasności sceny;

    \item \textbf{prądy upływu} (ang. \textit{leakage currents}) -- niewielkie
    prądy pasożytnicze w~układzie piksela mogą powodować zmianę jego stanu
    wewnętrznego nawet przy stałym oświetleniu. Prowadzi to do powstawania
    fałszywych zdarzeń tła, określanych jako \textit{leak events};

    \item \textbf{czas martwy piksela} (ang. \textit{refractory period}) --
    po wygenerowaniu zdarzenia piksel przez krótki czas nie może wygenerować
    kolejnego. Przy szybko zmieniającym się sygnale może to prowadzić do
    pominięcia części zdarzeń;

    \item \textbf{ograniczone pasmo fotoreceptora} -- układ piksela ma
    skończoną szybkość odpowiedzi na zmianę natężenia światła. Szybkie zmiany
    mogą być przez to tłumione lub rejestrowane z~opóźnieniem, szczególnie
    przy niskim poziomie oświetlenia.
\end{itemize}

Wymienione efekty powodują zarówno powstawanie izolowanych zdarzeń tła,
jak i~pomijanie lub opóźnianie zdarzeń odpowiadających rzeczywistym zmianom
jasności.


\subsection{Reprezentacje zdarzeń}
\label{subsec:reprezentacje_zdarzen}

Surowy strumień zachowuje dokładne znaczniki czasu, ale nie ma regularnej struktury obrazu. Wiele metod grupuje zatem zdarzenia z~półotwartego okna $\mathcal{W}=[t_0,t_1)$ w~reprezentację o~stałym rozmiarze. Wybór długości okna jest kompromisem: dłuższe okno zawiera więcej aktywnych krawędzi, lecz zwiększa ich rozmycie podczas ruchu, a~krótsze może dostarczyć zbyt mało danych.

Do często stosowanych reprezentacji zdarzeń w~postaci obrazowej należą:

\begin{itemize}
    \item \textbf{Klatka binarna} (ang. \textit{binary frame}) (Rys. \ref{fig:event_binary_frame}) -- piksel informuje jedynie o~tym, czy w~danym położeniu w~rozpatrywanym oknie czasowym wystąpiło zdarzenie. W~najprostszej postaci piksel przyjmuje wartość aktywną niezależnie od liczby zdarzeń i~ich polaryzacji~\cite{blachut2023eventframe}.

    \begin{figure}[!ht]
        \centering
        \includegraphics[width=0.55\linewidth]{theory_event_binary_frame.png}
        \caption{Przykład klatki binarnej utworzonej z~danych zdarzeniowych. Źródło: \cite{blachut2023eventframe}.}
        \label{fig:event_binary_frame}
    \end{figure}

    \item \textbf{Klatka zdarzeniowa uwzględniająca polaryzację} (ang. \textit{event frame}) (Rys. \ref{fig:event_frame}), dla uproszczenia niżej nazywana standardową -- wartość piksela odpowiada polaryzacji zdarzenia zarejestrowanego w~danym położeniu:
    \begin{equation}
        F(\mathbf{x})=p_{\mathrm{last}}(\mathbf{x}),
        \label{eq:standard_event_frame}
    \end{equation}
    gdzie $p_{\mathrm{last}}(\mathbf{x})\in\{-1,1\}$ oznacza polaryzację ostatniego zdarzenia w~pikselu $\mathbf{x}$. Piksel, w~którym nie wystąpiło żadne zdarzenie, otrzymuje wartość neutralną. Reprezentacja zachowuje informację o~polaryzacji, ale nie o~dokładnym czasie wystąpienia zdarzenia~\cite{blachut2023eventframe}.

    \begin{figure}[!ht]
        \centering
        \includegraphics[width=0.55\linewidth]{theory_event_frame.png}
        \caption{Przykład klatki zdarzeniowej uwzględniającej polaryzację. Źródło: \cite{blachut2023eventframe}.}
        \label{fig:event_frame}
    \end{figure}

    \item \textbf{Wykładniczo zanikająca klatka zdarzeniowa} (ang. \textit{exponentially decaying time surface}) (Rys. \ref{fig:event_time_surface}) -- oprócz polaryzacji uwzględnia czas, jaki upłynął od ostatniego zdarzenia w~danym pikselu do chwili odniesienia. Wartość piksela można zapisać jako
    \begin{equation}
        F(\mathbf{x},t_r)=
        p_{\mathrm{last}}(\mathbf{x})
        \exp\left(
        -\frac{t_r-t_{\mathrm{last}}(\mathbf{x})}{\tau}
        \right),
        \label{eq:exponential_event_frame}
    \end{equation}
    gdzie $t_{\mathrm{last}}(\mathbf{x})$ jest czasem ostatniego zdarzenia w~pikselu, a~$\tau>0$ określa szybkość zaniku. Dzięki temu nowsze zdarzenia mają większy wpływ na wynikowy obraz, natomiast znaczenie starszych stopniowo maleje~\cite{blachut2023eventframe}.

    \begin{figure}[!ht]
        \centering
        \includegraphics[width=0.55\linewidth]{theory_event_exponential_frame.png}
        \caption{Przykład wykładniczo zanikającej klatki zdarzeniowej. Źródło: \cite{blachut2023eventframe}.}
        \label{fig:event_time_surface}
    \end{figure}

    \item \textbf{Siatka wokselowa} (ang. \textit{voxel grid}) -- oś czasu
    okna dzieli się na $B$ przedziałów, tworząc tensor
    $V\in\mathbb{R}^{B\times H\times W}$, gdzie $H$ i~$W$ oznaczają wysokość
    i~szerokość obrazu. Czas każdego zdarzenia jest przeskalowywany do zakresu
    odpowiadającego osi czasowej tensora:
    \begin{equation}
        \tilde{t}_k=(B-1)\frac{t_k-t_0}{t_1-t_0}.
        \label{eq:voxel_time}
    \end{equation}
    Zdarzenie jest następnie przypisywane do najbliższych przedziałów czasowych,
    a~jego wkład może być rozdzielony pomiędzy dwa sąsiednie przedziały za pomocą
    interpolacji liniowej. Dzięki temu siatka wokselowa zachowuje więcej informacji
    o~czasie występowania zdarzeń niż pojedyncza klatka, kosztem większego zużycia
    pamięci i~konieczności przetwarzania danych wielokanałowych
    ~\cite{zhu2019unsupervised}.
\end{itemize}

\section{Metody estymacji stanu}
\label{sec:estymacja_stanu}


\subsection{Odometria wizyjna, SLAM i VIO}
\label{subsec:vo_slam_vio}

\textbf{Odometria wizyjna} (ang. \textit{Visual Odometry}, VO) jest problemem estymacji ruchu kamery na podstawie sekwencji obserwacji wizyjnych~\cite{nister2004visualodometry}. Wyznaczane są względne przemieszczenia pomiędzy kolejnymi lub bliskimi w~czasie położeniami kamery, a~następnie składane w~trajektorię. W~podstawowym ujęciu VO nie wymaga tworzenia trwałej mapy ani rozpoznawania powrotów do wcześniej odwiedzonych miejsc. System może wykorzystywać lokalnie utrzymywane punkty sceny, jednak służą one przede wszystkim do wyznaczania bieżącego ruchu. Ponieważ każde oszacowanie jest obarczone błędem, jego kolejne składanie prowadzi do narastającego dryfu położenia i~orientacji.

\textbf{SLAM} (ang. \textit{Simultaneous Localization and Mapping}) polega na jednoczesnej estymacji stanu robota oraz mapy nieznanego otoczenia na podstawie wykonywanych pomiarów~\cite{cadena2016slam}. Mapa nie musi być gęstym modelem całej sceny. Może zawierać na przykład wybrane punkty przestrzeni i~pozy, z~których zostały zaobserwowane. System określa położenie czujnika względem tej reprezentacji, a~jednocześnie uzupełnia ją na podstawie kolejnych obserwacji.

W~porównaniu z~VO system SLAM utrzymuje zależności pomiędzy pomiarami wykonanymi w~różnych chwilach. Ponowne obserwacje tych samych elementów otoczenia wprowadzają dodatkowe ograniczenia pomiędzy pozami i~pozwalają poprawiać spójność trajektorii oraz mapy. Estymacja nie opiera się więc wyłącznie na składaniu kolejnych przemieszczeń, chociaż bez ograniczeń obejmujących dłuższe fragmenty trasy również może podlegać narastającemu dryfowi.

Jednym z~mechanizmów stosowanych w~systemach SLAM jest \textbf{domknięcie pętli} (ang. \textit{loop closure}), czyli wykrycie, że robot ponownie obserwuje wcześniej odwiedzony obszar. Pozwala ono powiązać odległe w~czasie fragmenty trajektorii i~skorygować część skumulowanego dryfu. Samo podobieństwo obserwacji nie jest wystarczające: przed wykorzystaniem takiego powiązania należy potwierdzić jego zgodność geometryczną.

\textbf{Odometria wizyjno-inercyjna} (ang. \textit{Visual-Inertial Odometry}, VIO) łączy obserwacje wizyjne z~pomiarami czujnika IMU~\cite{forster2015preintegration}. Kamera dostarcza informacji o~ruchu względem obserwowanych elementów otoczenia, natomiast żyroskop i~akcelerometr mierzą odpowiednio prędkość kątową oraz przyspieszenie właściwe. Pomiary inercyjne są zwykle dostępne z~większą częstotliwością i~pozwalają opisywać ruch pomiędzy kolejnymi obserwacjami wizyjnymi. Ich całkowanie powoduje jednak narastanie błędów związanych między innymi z~szumem i~błędami systematycznymi czujników. Ograniczenia wynikające z~obserwacji kamer pozwalają zmniejszać ten dryf, dlatego w~pełnym systemie VIO wspólnie estymuje się zazwyczaj pozę, prędkość oraz błędy systematyczne IMU.

VO i~VIO dotyczą przede wszystkim bieżącej estymacji ruchu, natomiast SLAM obejmuje również utrzymywanie mapy i~wykorzystywanie wcześniejszych obserwacji do poprawiania trajektorii. Pojęcia te nie są rozłączne: VIO określa sposób łączenia danych wizyjnych i~inercyjnych, który może być wykorzystywany zarówno w~odometrii, jak i~w~systemie SLAM.

\subsection{Podejścia do estymacji i wykorzystywane dane}
\label{subsec:podejscia_estymacja}

\textbf{Wizyjny SLAM} (ang. \textit{visual SLAM}, vSLAM) jest odmianą SLAM, w~której podstawowym źródłem informacji o~otoczeniu są czujniki wizyjne. System może wykorzystywać pojedynczą kamerę, układ stereo albo kamerę uzupełnioną innymi czujnikami. Kamera monokularna nie dostarcza bezpośrednio metrycznej skali translacji, natomiast znana geometria układu stereo pozwala wyznaczać głębokość punktów~\cite{hartley2004multiple}. Pomiary inercyjne mogą dodatkowo dostarczać informacji o~ruchu pomiędzy kolejnymi obserwacjami obrazu~\cite{forster2015preintegration}.

Chociaż tradycyjne algorytmy wizyjnego SLAM zaprojektowane zostały do przetwarzania klasycznych klatek obrazu, metody w nich wykorzystywane można z powodzeniem zastosować również do danych zdarzeniowych. Należy jednak pamiętać o ich odmiennej charakterystyce. Głębokie sieci neuronowe wymagają w tym celu dostosowania oraz dostrojenia wag. W celu zastosowania pozostałych metod, wykorzystywanych w tradycyjnym wizyjnym SLAM, można przekształcić strumień z kamery zdarzeniowej do reprezentacji obrazowej (~\ref{subsec:reprezentacje_zdarzen}), dzięki czemu otrzymujemy oparte na krawędziach obrazy w skali szarości, co w przybliżeniu odpowiada klatkom ze standardowej kamery.

Stosowane podejścia można podzielić na:~\cite{huang2026eventbased}:
\begin{itemize}
    \item \textbf{metody bezpośrednie} -- wyznaczają ruch bez tworzenia jawnych korespondencji punktowych, wyrównując zdarzenia z~obrazem odniesienia albo porównując ich regularne reprezentacje;
    \item \textbf{metody pośrednie, oparte na punktach charakterystycznych} -- wykrywają wybrane punkty, śledzą je lub dopasowują pomiędzy obserwacjami, a~następnie estymują ruch kamery i~strukturę sceny na podstawie jawnych korespondencji. Chociaż jest to metoda początkowo zaprojektowana jako sposób przetwarzania obrazów z tradycyjnej kamery, to zastosowanie reprezentacji danych zdarzeniowych jako klatki pozwala na użycie tej metody również dla czujników DVS. Powstały również algorytmy dedykowane do znajdowania punktów charakterystycznych w strumieniu danych z kamery zdarzeniowej \cite{vasco2016eharris}.
    \item \textbf{metody oparte na kompensacji ruchu} -- przekształcają zdarzenia zgodnie z~rozpatrywanym modelem ruchu i~dobierają jego parametry tak, aby zdarzenia po wyrównaniu tworzyły możliwie spójną reprezentację;
    \item \textbf{metody wykorzystujące uczenie głębokie} -- używają modeli przewidujących na podstawie zdarzeń, między innymi ruch kamery, przepływ optyczny lub głębokość sceny.
\end{itemize}

Dane zdarzeniowe można przetwarzać bezpośrednio jako asynchroniczny strumień, albo najpierw grupować je w~regularne reprezentacje obrazowe opisane w~podsekcji~\ref{subsec:reprezentacje_zdarzen}. Pierwsze podejście wykorzystuje położenia i~dokładne znaczniki czasu pojedynczych zdarzeń. Drugie umożliwia zastosowanie wybranych metod opracowanych dla obrazów, ponieważ wynik ma regularną siatkę pikseli. Wartości takiej reprezentacji opisują jednak aktywność zdarzeń~\cite{gallego2022event,huang2026eventbased}.

Dalsza część rozdziału koncentruje się na podejściu pośrednim, wykorzystującym klatki utworzone ze zdarzeń. Na regularnej reprezentacji można stosować klasyczne detektory punktów, takie jak GFTT i~FAST, a~następnie śledzić punkty metodą Lucas--Kanade albo porównywać je za pomocą deskryptorów~\cite{shi1994good,rosten2006fast,lucas1981iterative}.~W pracy~\cite{rebecq2017eventvio} oraz~w~\cite{vidal2018ultimate}, wykorzystano wirtualne klatki zdarzeniowe wraz z~detektorem FAST i~śledzeniem Lucas--Kanade. Korespondencje punktów pomiędzy kamerami pozwalają wyznaczyć ich położenia 3D, które wraz z~bieżącymi obserwacjami 2D służą do estymacji pozy kamery.

\section{Punkty charakterystyczne i ich śledzenie}
\label{sec:punkty_charakterystyczne}
\label{subsec:sledzenie_cech}

Metody, bazujące na charakterystycznych elementach obrazu, wybierają punkty o~lokalnym otoczeniu pozwalającym na ich ponowne odnalezienie, a~następnie tworzą korespondencje pomiędzy obserwacjami. Agregowanie strumienia zdarzeń do postaci ramek zdarzeniowych pozwala z zastosowanie w tego typu systemach algorytmów przeznaczonych do przetwarzania obrazu, co jest zaletą takiego podejścia.

Z~obrazu wybierane są punkty charakterystyczne, których położenie można wiarygodnie wyznaczyć i~porównywać pomiędzy kolejnymi obserwacjami. \textbf{Detektor punktów charakterystycznych} wskazuje najczęściej narożniki lub inne obszary o~wyraźnej zmianie wartości obrazu w~więcej niż jednym kierunku. Punkty znajdujące się na obszarach jednorodnych lub prostych krawędziach są mniej użyteczne, ponieważ ich położenie jest trudniejsze do jednoznacznego określenia.

\textbf{Detektor Shi--Tomasiego} -- GFTT (ang. \textit{Good Features to Track}) ocenia lokalne zmiany wartości obrazu na podstawie macierzy utworzonej z~gradientów. Punkt jest uznawany za odpowiedni do śledzenia, gdy obie wartości własne tej macierzy są dostatecznie duże; w~praktyce kryterium wykorzystuje mniejszą z~nich. Pozwala to odrzucić zarówno obszary jednorodne, jak i~punkty leżące na prostych krawędziach~\cite{shi1994good}.

\textbf{Detektor FAST} (ang. \textit{Features from Accelerated Segment Test}) wykrywa narożniki na podstawie wartości pikseli leżących na okręgu wokół analizowanego punktu. Punkt zostaje uznany za narożnik, jeżeli odpowiednio długi ciąg sąsiednich pikseli okręgu jest wyraźnie jaśniejszy lub ciemniejszy od jego środka. Metoda została zaprojektowana z~myślą o~bardzo szybkim wykrywaniu punktów charakterystycznych~\cite{rosten2006fast}.

Samo położenie punktu nie wystarcza jednak do jego ponownego rozpoznania po większej zmianie obrazu. W~tym celu stosuje się \textbf{deskryptory punktów charakterystycznych}, czyli reprezentacje numeryczne opisujące lokalne otoczenie punktu. Umożliwiają one porównywanie punktów wykrytych w~różnych obrazach i~wyznaczanie ich korespondencji.

Przykładem często stosowanej metody jest \textbf{ORB} (ang. \textit{Oriented FAST and Rotated BRIEF}), która łączy detekcję punktów FAST z~binarnym deskryptorem opartym na BRIEF (ang. \textit{Binary Robust Independent Elementary Features}). Orientacja punktu jest uwzględniana podczas konstrukcji deskryptora, co zwiększa jego odporność na obrót obrazu~\cite{rublee2011orb}. Ponieważ deskryptor ma postać binarną, podobieństwo dwóch punktów może być szybko oceniane za pomocą odległości Hamminga. Uzyskane dopasowania mogą jednak zawierać błędne korespondencje i~zwykle wymagają późniejszej weryfikacji geometrycznej.

Jeżeli punkt jest obserwowany w~kolejnych, blisko położonych w~czasie obrazach, zamiast ponownie wyszukiwać go na podstawie deskryptora można bezpośrednio \textbf{śledzić jego położenie}. Jako metoda śledzenia może być wykorzystany między innymi przepływ optyczny Lucas--Kanade~\cite{lucas1981iterative}. Metoda zakłada, że wartości obserwowanego fragmentu obrazu zmieniają się niewiele pomiędzy kolejnymi klatkami oraz że ruch w~niewielkim otoczeniu punktu jest w~przybliżeniu jednakowy. Na tej podstawie wyznaczane jest przesunięcie najlepiej dopasowujące lokalny fragment pierwszego obrazu do drugiego.

Podstawowa metoda Lucas--Kanade dobrze radzi sobie z~niewielkimi przemieszczeniami. W~wariancie piramidalnym obrazy są analizowane na kilku poziomach rozdzielczości: duże przemieszczenie jest najpierw oszacowywane na obrazie o~mniejszej rozdzielczości, a~następnie kolejno doprecyzowywane na dokładniejszych poziomach piramidy. Pozwala to stosować metodę również przy większym ruchu pomiędzy kolejnymi obrazami. W~przypadku klatek zdarzeniowych śledzony jest lokalny wzorzec aktywności zdarzeń, a~nie jasność zarejestrowana przez kamerę klatkową~\cite{rebecq2017eventvio,vidal2018ultimate}.


\section{Podstawy geometryczne estymacji ruchu}
\label{sec:geometria_stereo}

Estymacja ruchu wymaga jednoznacznego zapisu położenia i~orientacji kamery oraz modelu wiążącego punkty przestrzeni z~ich obserwacjami na obrazie. W~tej sekcji wprowadzono stosowaną konwencję transformacji, a~następnie przedstawiono geometrię skalibrowanego układu stereo.

\subsection{Rotacja i translacja}
\label{subsec:translacja_i_rotacja}

\textbf{Położenie (translacja)} określa pozycję początku jednego układu współrzędnych względem drugiego i~jest opisywane wektorem $\mathbf{t}\in\mathbb{R}^3$. \textbf{Orientacja (rotacja)} opisuje wzajemne ustawienie osi tych układów i~może być reprezentowana przez macierz obrotu $\mathbf{R}\in SO(3)$. Grupa $SO(3)$ obejmuje obroty w~przestrzeni trójwymiarowej reprezentowane przez macierze ortogonalne o~wyznaczniku równym~1~\cite{blanco2021se3}. Połączenie położenia i~orientacji nazywane jest \textbf{pozą} (ang. \textit{pose}).

Pozę jednego układu współrzędnych względem drugiego można reprezentować za pomocą jednorodnej macierzy transformacji $\mathbf{T}\in SE(3)$. Grupa $SE(3)$ obejmuje sztywne transformacje w~przestrzeni trójwymiarowej, złożone z~obrotu i~translacji, które zachowują kształt oraz rozmiar obiektów~\cite{blanco2021se3}. Macierz transformacji ma postać
\begin{equation}
    \mathbf{T}
    =
    \begin{bmatrix}
        \mathbf{R} & \mathbf{t}\\
        \mathbf{0}^{T} & 1
    \end{bmatrix},
    \label{eq:se3_transform}
\end{equation}
gdzie $\mathbf{R}$ opisuje rotację pomiędzy układami, $\mathbf{t}$ ich względne przesunięcie, a~symbol $\mathbf{0}^{T}$ oznacza trójelementowy wiersz zerowy.

Zapis $\mathbf{T}_{A B}$ oznacza transformację współrzędnych z~układu $B$ do układu $A$. Dla punktu $\bar{\mathbf{X}}_B$ zapisanego we współrzędnych jednorodnych zachodzi
\begin{equation}
    \bar{\mathbf{X}}_A
    =
    \mathbf{T}_{A B}\bar{\mathbf{X}}_B.
    \label{eq:frame_transform}
\end{equation}
Macierz $\mathbf{R}_{A B}$ i~wektor $\mathbf{t}_{A B}$ są odpowiednio częścią obrotową i~translacyjną transformacji $\mathbf{T}_{A B}$. Transformacje można składać zgodnie z~kolejnością układów, na przykład $\mathbf{T}_{A C}=\mathbf{T}_{A B}\mathbf{T}_{B C}$. Transformację w~przeciwnym kierunku otrzymuje się przez odwrócenie macierzy: $\mathbf{T}_{B A}=\mathbf{T}_{A B}^{-1}$.

W~dalszej części pracy poza kamery ma postać $\mathbf{T}_{W C}$, a~więc przekształca punkt z~układu kamery $C$ do układu świata $W$. Uporządkowany w~czasie zbiór takich póz tworzy \textbf{trajektorię}.

\subsection{Model kamery i kalibracja stereo}
\label{subsec:model_kamery}

Pojedyncza obserwacja punktu obrazu określa kierunek, w~którym znajduje się punkt sceny, lecz nie pozwala wyznaczyć jego odległości od kamery. Obserwacja tego samego punktu przez drugą kamerę dostarcza dodatkowego ograniczenia geometrycznego. Jeżeli względne położenie kamer jest znane z~kalibracji, położenie punktu w~przestrzeni może zostać wyznaczone przez triangulację, a~znana odległość pomiędzy kamerami pozwala określić jego głębokość w~skali metrycznej~\cite{fusiello2000rectification}.

W~modelu kamery otworkowej (ang. \textit{pinhole camera}) punkt
$\mathbf{X}_C=(X,Y,Z)^T$, zapisany w~układzie kamery, jest rzutowany
na punkt obrazu $(u,v)^T$ zgodnie z~zależnością~\cite{fusiello2000rectification}
\begin{equation}
    s
    \begin{bmatrix}
        u \\
        v \\
        1
    \end{bmatrix}
    =
    \mathbf{K}
    \begin{bmatrix}
        X \\
        Y \\
        Z
    \end{bmatrix},
    \qquad
    \mathbf{K}
    =
    \begin{bmatrix}
        f_x & 0   & c_x \\
        0   & f_y & c_y \\
        0   & 0   & 1
    \end{bmatrix}.
    \label{eq:pinhole_projection}
\end{equation}
gdzie $s$ jest współczynnikiem skali, $f_x$ i~$f_y$ oznaczają ogniskowe
wyrażone w~pikselach, a~$(c_x,c_y)$ jest punktem głównym, czyli miejscem
przecięcia osi optycznej z~płaszczyzną obrazu. Macierz $\mathbf{K}$
zawiera więc \textbf{parametry wewnętrzne kamery}.

Punkt obrazu można zapisać we \textbf{współrzędnych jednorodnych} jako
$\bar{\mathbf{u}}=(u,v,1)^T$. W~takim zapisie wektory różniące się niezerowym współczynnikiem skali reprezentują ten sam punkt. Pozwala to zapisać rzutowanie punktu 3D z~dowolnego układu odniesienia w~postaci
\begin{equation}
\bar{\mathbf{u}}
\sim
\mathbf{P}\bar{\mathbf{X}},
\qquad
\mathbf{P}
=
\mathbf{K}
\begin{bmatrix}
\mathbf{R} & \mathbf{t}
\end{bmatrix},
\label{eq:camera_projection_matrix}
\end{equation}
gdzie $\bar{\mathbf{X}}=(X,Y,Z,1)^T$ jest jednorodnym zapisem punktu 3D, a~symbol $\sim$ oznacza równość z~dokładnością do niezerowego współczynnika skali. Macierz $\mathbf{P}\in\mathbb{R}^{3\times4}$ nazywana jest \textbf{macierzą projekcji} kamery. Łączy ona parametry wewnętrzne $\mathbf{K}$ z~rotacją $\mathbf{R}$ i~translacją $\mathbf{t}$ określającymi położenie kamery względem przyjętego układu odniesienia~\cite{fusiello2000rectification} i~pozwala na odwzorowanie punktu przestrzeni trójwymiarowej na odpowiadający mu punkt obrazu.

Rzeczywista kamera odbiega od idealnego modelu otworkowego między innymi wskutek dystorsji obiektywu. Parametry modelu dystorsji są dlatego wyznaczane podczas kalibracji, a~w~dalszym opisie geometrii zakłada się punkty po jej korekcji.

\textbf{Kalibracja stereo} obejmuje wyznaczenie parametrów wewnętrznych i~modeli dystorsji obu kamer oraz ich względnego położenia i~orientacji. Przyjmując lewą kamerę jako układ odniesienia, transformację z~układu lewej kamery $C_L$ do prawej $C_R$ można zapisać jako
\begin{equation}
\mathbf{T}_{C_R C_L}
=
\begin{bmatrix}
\mathbf{R} & \mathbf{t}\\
\mathbf{0}^T & 1
\end{bmatrix},
\label{eq:stereo_extrinsics}
\end{equation}
co oznacza, że współrzędne punktu w~obu układach spełniają
\begin{equation}
\mathbf{X}_R
=
\mathbf{R}\mathbf{X}_L+\mathbf{t}.
\label{eq:stereo_relative_pose}
\end{equation}
Odległość pomiędzy środkami optycznymi kamer
\begin{equation}
B=\|\mathbf{t}\|
\label{eq:stereo_baseline}
\end{equation}
nazywana jest \textbf{bazą układu stereo}.

\textbf{Rektyfikacja stereo} przekształca obrazy obu kamer na podstawie ich kalibracji tak, aby odpowiadające sobie punkty znajdowały się w~tych samych wierszach (liniach epipolarnych). Można ją interpretować jako zastąpienie rzeczywistych kamer parą kamer wirtualnych o~równoległych płaszczyznach obrazowania i~osi poziomej zgodnej z~bazą stereo. Wyszukiwanie korespondencji zostaje dzięki temu ograniczone z~obszaru dwuwymiarowego do pojedynczego wiersza obrazu.~\cite{fusiello2000rectification}.

\subsection{Geometria epipolarna}
\label{subsec:geometria_epipolarna}

Punkt obrazu wyznacza promień wychodzący ze środka optycznego kamery, na którym może znajdować się obserwowany punkt sceny. Rzut tego promienia na płaszczyznę obrazu drugiej kamery tworzy \textbf{linię epipolarną}. Odpowiadający punkt w~drugim obrazie nie może więc znajdować się w~dowolnym miejscu, lecz musi leżeć na tej linii~\cite{fusiello2000rectification}.

Przyjmując układ lewej kamery jako układ odniesienia, macierze projekcji obu kamer mają postać
\begin{equation}
\mathbf{P}_L
=
\mathbf{K}_L
\begin{bmatrix}
\mathbf{I} & \mathbf{0}
\end{bmatrix},
\qquad
\mathbf{P}_R
=
\mathbf{K}_R
\begin{bmatrix}
\mathbf{R} & \mathbf{t}
\end{bmatrix},
\label{eq:stereo_projection_matrices}
\end{equation}
gdzie $\mathbf{I}$ jest macierzą jednostkową, natomiast $\mathbf{R}$ i~$\mathbf{t}$ są parametrami względnej pozy kamer z~równania~\eqref{eq:stereo_relative_pose}.

Wpływ parametrów wewnętrznych kamer można usunąć, przechodząc z~jednorodnych współrzędnych pikselowych $\bar{\mathbf{u}}_L$ i~$\bar{\mathbf{u}}_R$ do \textbf{znormalizowanych współrzędnych obrazu}
\begin{equation}
\mathbf{q}_L
=
\mathbf{K}_L^{-1}\bar{\mathbf{u}}_L,
\qquad
\mathbf{q}_R
=
\mathbf{K}_R^{-1}\bar{\mathbf{u}}_R.
\label{eq:normalized_image_coordinates}
\end{equation}

Dla skalibrowanej pary kamer geometrię odpowiadających sobie punktów opisuje macierz $\mathbf{E}$ (ang. \textit{essential matrix})
\begin{equation}
    \mathbf{E}
    =
    [\mathbf{t}]_{\times}\mathbf{R},
    \label{eq:essential_matrix}
\end{equation}
gdzie $[\mathbf{t}]_{\times}$ oznacza macierz operatora iloczynu
wektorowego, zdefiniowaną jako
\begin{equation}
    [\mathbf{t}]_{\times}
    =
    \begin{bmatrix}
        0    & -t_z & t_y \\
        t_z  & 0    & -t_x \\
        -t_y & t_x  & 0
    \end{bmatrix}.
    \label{eq:cross_product_matrix}
\end{equation}
Dla dowolnego wektora $\mathbf{a}\in\mathbb{R}^3$ zachodzi więc
\begin{equation}
    [\mathbf{t}]_{\times}\mathbf{a}
    =
    \mathbf{t}\times\mathbf{a}.
    \label{eq:cross_product_operator}
\end{equation}

Odpowiadające sobie punkty w~obu obrazach spełniają wówczas \textbf{ograniczenie epipolarne}
\begin{equation}
\mathbf{q}_R^T
\mathbf{E}
\mathbf{q}_L
=
0.
\label{eq:essential_constraint}
\end{equation}
Oznacza ono, że punkt odpowiadający obserwacji z~jednej kamery może znajdować się w~drugim obrazie jedynie na wyznaczonej przez geometrię kamer linii epipolarnej.

Jeżeli punkty są zapisane bezpośrednio we współrzędnych pikselowych,
tę samą zależność opisuje macierz $\mathbf{F}$
(ang. \textit{fundamental matrix}),
\begin{equation}
    \mathbf{F}
    =
    \mathbf{K}_R^{-T}
    \mathbf{E}
    \mathbf{K}_L^{-1},
    \label{eq:fundamental_matrix}
\end{equation}
dla której
\begin{equation}
    \bar{\mathbf{u}}_R^T
    \mathbf{F}
    \bar{\mathbf{u}}_L
    =
    0.
    \label{eq:fundamental_constraint}
\end{equation}

Zależności te są zgodne z~modelem skalibrowanej pary kamer przedstawionym w~\cite{nister2004fivepoint}.

Macierze $\mathbf{E}$ i~$\mathbf{F}$ opisują zatem to samo ograniczenie geometryczne, odpowiednio dla współrzędnych znormalizowanych i~pikselowych~\cite{nister2004fivepoint}.

\subsection{Dysparycja i triangulacja}
\label{subsec:dysparycja_triangulacja}

\textbf{Triangulacja} polega na wyznaczeniu położenia punktu w~przestrzeni na podstawie jego projekcji w~co najmniej dwóch obrazach oraz znanej geometrii kamer. Dla pary stereo punkt $\bar{\mathbf{X}}$ powinien spełniać zależności
\begin{equation}
\bar{\mathbf{u}}_L
\sim
\mathbf{P}_L\bar{\mathbf{X}},
\qquad
\bar{\mathbf{u}}_R
\sim
\mathbf{P}_R\bar{\mathbf{X}},
\label{eq:stereo_triangulation}
\end{equation}
gdzie $\mathbf{P}_L$ i~$\mathbf{P}_R$ są macierzami projekcji z~równania~\eqref{eq:stereo_projection_matrices}. W~obecności błędów pomiarowych promienie odpowiadające obu obserwacjom nie muszą przecinać się dokładnie, dlatego położenie punktu wyznacza się jako rozwiązanie najlepiej spełniające oba ograniczenia~\cite{kitani2017triangulation}.

Dla zrektyfikowanego układu stereo triangulacja przyjmuje szczególnie prostą postać. \textbf{Dysparycja} $d$ jest różnicą poziomych współrzędnych odpowiadających sobie punktów:
\begin{equation}
d=u'_L-u'_R.
\label{eq:disparity}
\end{equation}
Dla standardowego zrektyfikowanego układu, w~którym osie optyczne wirtualnych kamer są równoległe, a~baza $B$ jest skierowana wzdłuż osi poziomej, głębokość punktu $Z$ wynosi
\begin{equation}
Z=\frac{f'_x B}{d},
\label{eq:stereo_depth}
\end{equation}
gdzie $f'_x$ jest ogniskową poziomą kamery po rektyfikacji, wyrażoną w~pikselach. Współrzędne punktu w~układzie lewej zrektyfikowanej kamery można następnie wyznaczyć jako
\begin{equation}
X=\frac{(u'_L-c'_x)Z}{f'_x},
\qquad
Y=\frac{(v'_L-c'_y)Z}{f'_y},
\label{eq:stereo_point}
\end{equation}
gdzie $f'_x$, $f'_y$, $c'_x$ i~$c'_y$ są parametrami wewnętrznymi wirtualnej kamery po rektyfikacji. Nie muszą być one identyczne z~parametrami pierwotnej kamery z~równania~\eqref{eq:pinhole_projection}.

Ideę triangulacji przedstawiono na rysunku~\ref{fig:triangulation}.

\begin{figure}[!ht]
\centering
\includegraphics[width=0.76\linewidth]{theory_stereo_triangulation_geometry.png}
\caption{Idea triangulacji punktu na podstawie dwóch projekcji. W~obecności szumu promienie nie muszą przecinać się dokładnie, dlatego wybierane jest rozwiązanie najlepiej spełniające oba ograniczenia. Źródło:~\cite{kitani2017triangulation}.}
\label{fig:triangulation}
\end{figure}

Wrażliwość wyznaczonej głębokości na błąd dysparycji wynika z~równania~\eqref{eq:stereo_depth}:
\begin{equation}
\left|
\frac{\partial Z}{\partial d}
\right|
=
\frac{f'_x B}{d^2}.
\label{eq:depth_sensitivity}
\end{equation}
Ten sam błąd wyznaczenia dysparycji powoduje więc coraz większy błąd głębokości wraz ze wzrostem odległości punktu od kamer.


\section{Estymacja pozy metodą PnP}
\label{sec:estymacja_pnp}
\label{subsec:pnp}

\textbf{Problem PnP} (ang. \textit{Perspective-n-Point}) polega na wyznaczeniu pozy skalibrowanej kamery na podstawie zbioru korespondencji pomiędzy punktami $\mathbf{X}_i$ o~znanych współrzędnych 3D a~ich obserwacjami $\mathbf{u}_i$ na obrazie~\cite{lepetit2009epnp}. Przyjmując, że punkty 3D są zapisane w~układzie świata $W$, szukane parametry $\mathbf{R}$ i~$\mathbf{t}$ transformują je do układu kamery $C$ i~spełniają
\begin{equation}
s_i\bar{\mathbf{u}}_i
=
\mathbf{K}
\left(
\mathbf{R}\mathbf{X}_i+\mathbf{t}
\right),
\label{eq:pnp}
\end{equation}
gdzie $\bar{\mathbf{u}}_i$ jest punktem obrazu zapisanym we współrzędnych jednorodnych, $s_i$ współczynnikiem skali projekcji, a~$\mathbf{K}$ macierzą parametrów wewnętrznych kamery.

Wyznaczone parametry opisują transformację
\begin{equation}
\mathbf{T}_{CW}
=
\begin{bmatrix}
\mathbf{R} & \mathbf{t} \\
\mathbf{0}^T & 1
\end{bmatrix},
\end{equation}
czyli transformację punktów z~układu świata do układu kamery. Jeżeli poza kamery jest w~pracy reprezentowana jako $\mathbf{T}_{WC}$, należy zastosować transformację odwrotną $\mathbf{T}_{WC}=\mathbf{T}_{CW}^{-1}$.

W~praktyce część korespondencji 3D--2D może być błędna. Do odpornej estymacji pozy można zastosować \textbf{RANSAC} (ang. \textit{Random Sample Consensus}), który wielokrotnie losuje minimalne podzbiory danych, wyznacza na ich podstawie hipotezę modelu, a~następnie określa zbiór obserwacji zgodnych z~tą hipotezą w~zadanej tolerancji~\cite{fischler1981ransac}.

W~przypadku PnP zgodność korespondencji można oceniać na podstawie \textbf{błędu reprojekcji}, czyli odległości pomiędzy zmierzonym punktem obrazu $\mathbf{u}_i$ a~położeniem $\hat{\mathbf{u}}_i$ otrzymanym przez ponowną projekcję odpowiadającego mu punktu 3D:
\begin{equation}
e_i
=
\left\|
\mathbf{u}_i-\hat{\mathbf{u}}_i
\right\|_2.
\label{eq:reprojection_error}
\end{equation}
Korespondencje, dla których błąd jest mniejszy od przyjętego progu, tworzą zbiór obserwacji zgodnych z~daną hipotezą. Po wybraniu najlepszego zbioru model może zostać ponownie wyznaczony z~wykorzystaniem wszystkich należących do niego obserwacji~\cite{fischler1981ransac}.


\section{Graf pozy i domykanie pętli}
\label{sec:graf_pozy}

\textbf{Graf pozy} jest reprezentacją trajektorii, w~której węzły odpowiadają pozom robota, a~krawędzie opisują ograniczenia pomiędzy parami póz~\cite{grisetti2010graph}. Poza węzła $i$ jest zapisana jako $\mathbf{T}_i=\mathbf{T}_{W C_i}\in SE(3)$, czyli przekształca współrzędne z~układu kamery $C_i$ do układu świata $W$.

Krawędź $(i,j)$ zawiera pomiar względnej transformacji $\mathbf{Z}_{ij}=\mathbf{T}_{C_i C_j}$. Zgodnie z~konwencją przyjętą w~podsekcji~\ref{subsec:translacja_i_rotacja} transformacja ta przekształca współrzędne z~układu kamery $C_j$ do układu kamery $C_i$. Krawędź może być również związana z~macierzą informacji $\boldsymbol{\Omega}_{ij}$, określającą względne znaczenie poszczególnych składowych błędu~\cite{grisetti2010graph}.

Na podstawie aktualnych oszacowań obu póz można przewidzieć transformację względną:
\begin{equation}
    \hat{\mathbf{Z}}_{ij}
    =
    \mathbf{T}_i^{-1}\mathbf{T}_j.
    \label{eq:predicted_relative_pose}
\end{equation}

Zgodność pomiaru z~aktualnymi pozami można sprawdzić, obliczając transformację błędu:
\begin{equation}
    \mathbf{E}_{ij}
    =
    \mathbf{Z}_{ij}^{-1}\hat{\mathbf{Z}}_{ij}.
    \label{eq:pose_graph_error_transform}
\end{equation}
Jeżeli pomiar i~przewidywanie są identyczne, to $\mathbf{E}_{ij}=\mathbf{I}$. Macierz błędu zawiera jednak jednocześnie translację i~obrót, dlatego nie jest bezpośrednio używana jako reszta optymalizacji.

Mapa logarytmiczna grupy $SE(3)$ przekształca transformację błędu w~wektor sześciowymiarowy:
\begin{equation}
    \mathbf{e}_{ij}
    =
    \operatorname{Log}\!\left(\mathbf{E}_{ij}\right)
    =
    \operatorname{Log}\!\left(
        \mathbf{Z}_{ij}^{-1}
        \mathbf{T}_i^{-1}
        \mathbf{T}_j
    \right)
    \in\mathbb{R}^{6}.
    \label{eq:pose_graph_error}
\end{equation}
Operator $\operatorname{Log}$ w~równaniu~\eqref{eq:pose_graph_error} nie oznacza obliczania logarytmu każdego elementu macierzy. Odwzorowuje on transformację należącą do $SE(3)$ na sześć liczb: trzy składowe opisujące błąd translacji oraz trzy składowe wektora obrotu. Kierunek wektora obrotu wyznacza oś obrotu, a~jego długość odpowiada wartości kąta wyrażonej w~radianach. Dla $\mathbf{E}_{ij}=\mathbf{I}$ otrzymuje się $\mathbf{e}_{ij}=\mathbf{0}$~\cite{blanco2021se3}.

Optymalizacja grafu polega na znalezieniu takich póz, dla których ważone błędy wszystkich krawędzi są możliwie małe:
\begin{equation}
    \min_{\{\mathbf{T}_i\}}
    \sum_{(i,j)\in\mathcal{E}}
    \rho\!\left(
        \mathbf{e}_{ij}^{T}
        \boldsymbol{\Omega}_{ij}
        \mathbf{e}_{ij}
    \right),
    \label{eq:pose_graph}
\end{equation}
gdzie $\mathcal{E}$ jest zbiorem krawędzi, a~$\rho$ oznacza odporną funkcję kosztu ograniczającą wpływ dużych błędów. Ważenie poszczególnych składowych jest istotne między innymi dlatego, że błędy translacji i~rotacji mają różne jednostki.

Pierwsza poza grafu jest zwykle ustalana, ponieważ jednoczesne wykonanie tej samej sztywnej transformacji na wszystkich węzłach nie zmienia żadnej transformacji względnej. Usuwa to swobodę globalnej translacji i~rotacji rozwiązania~\cite{grisetti2010graph}.

\textbf{Domknięcie pętli} (ang. \textit{loop closure}) polega na wykryciu, że robot ponownie obserwuje wcześniej odwiedzony fragment otoczenia. Umożliwia to dodanie krawędzi pomiędzy węzłami odległymi w~czasie i~skorygowanie narastającego dryfu trajektorii.

Kandydaci na domknięcie pętli mogą być wybierani na podstawie podobieństwa bieżącej obserwacji do wcześniej zapisanych klatek. Samo podobieństwo wizualne nie potwierdza jednak zgodności geometrycznej, dlatego przed dodaniem krawędzi konieczne jest wyznaczenie transformacji względnej i~sprawdzenie jakości dopasowania~\cite{murartal2015orbslam}.

\section{Wyrównanie trajektorii metodą Umeyamy}
\label{subsec:wyrownanie_umeyamy}

Trajektoria estymowana i~trajektoria referencyjna mogą opisywać ten sam ruch, ale być zapisane w~różnych układach współrzędnych. Bezpośrednie porównanie odpowiadających sobie pozycji prowadziłoby wtedy do błędu wynikającego nie tylko z~niedokładności estymacji ruchu, lecz również z~różnicy pomiędzy przyjętymi układami odniesienia. Z~tego powodu przed obliczeniem błędu trajektorie mogą zostać wzajemnie wyrównane~\cite{sturm2012benchmark}.

Metoda Umeyamy wyznacza pojedynczą transformację, która w~sensie najmniejszych kwadratów najlepiej nakłada zbiór pozycji estymowanych na odpowiadające im pozycje referencyjne~\cite{umeyama1991least}. Dla par punktów $\mathbf{x}_i,\mathbf{y}_i\in\mathbb{R}^3$, oznaczających odpowiednio pozycję estymowaną i~referencyjną w~tej samej chwili, rozwiązywane jest zadanie
\begin{equation}
    \underset{s,\mathbf{R},\mathbf{t}}{\operatorname{min}}
    \frac{1}{N}\sum_{i=1}^{N}
    \left\|
        \mathbf{y}_i-
        \left(
            s\mathbf{R}\mathbf{x}_i+\mathbf{t}
        \right)
    \right\|^2,
    \qquad
    \mathbf{R}\in SO(3),
    \label{eq:umeyama_objective}
\end{equation}
gdzie $\mathbf{R}$ opisuje obrót pomiędzy układami współrzędnych, $\mathbf{t}$ ich względne przesunięcie, a~$s$ jest współczynnikiem skali.

Istotne jest, że ta sama transformacja jest stosowana do całej trajektorii. Wyrównanie nie koryguje więc lokalnych błędów estymacji ani narastającego dryfu, lecz usuwa globalną różnicę położenia i~orientacji pomiędzy układami odniesienia. Po wyrównaniu odległość pomiędzy odpowiadającymi sobie punktami trajektorii wynika zatem z~różnic w~oszacowanym przebiegu ruchu, a~nie z~dowolnego wyboru początku i~orientacji układu współrzędnych.

W~wariancie $SE(3)$ przyjmuje się $s=1$, dlatego wyrównaniu podlegają wyłącznie obrót i~translacja. Zachowana zostaje wówczas skala wyznaczona przez algorytm. W~wariancie $Sim(3)$ współczynnik $s$ jest dodatkowo estymowany, co pozwala usunąć również stałą globalną różnicę skali pomiędzy trajektoriami~\cite{umeyama1991least}. Zastosowanie $Sim(3)$ powoduje jednak, że globalny błąd skali nie jest uwzględniany w~późniejszej ocenie, dlatego wybór rodzaju wyrównania powinien odpowiadać temu, jaki aspekt działania systemu ma być oceniany.

    \chapter{Przegląd literatury}
\label{cha:przeglad_literatury}

W~rozdziale omówiono systemy VO i~SLAM, które dostarczyły istotnych rozwiązań architektonicznych lub algorytmicznych. Systemy PTAM, S-PTAM i~ORB-SLAM zostały opracowane dla obrazów i~nie przetwarzają bezpośrednio asynchronicznego strumienia zdarzeń. Ich uwzględnienie wynika z~możliwości zastosowania wybranych mechanizmów po wcześniejszym zgrupowaniu zdarzeń w~regularne reprezentacje obrazowe. Na takich reprezentacjach można wykrywać i~śledzić punkty, natomiast klatki kluczowe, rzadka mapa, relokalizacja i~domykanie pętli mogą działać według zasad znanych z~klasycznych systemów SLAM. Nie oznacza to wykorzystania ORB-SLAM jako gotowego algorytmu dla kamer zdarzeniowych, lecz adaptację wybranych elementów jego architektury. Możliwość łączenia klatek zdarzeniowych z~metodami opracowanymi dla obrazów pokazano między innymi w~systemach Rebecqa i~UltimateSLAM~\cite{rebecq2017eventvio,vidal2018ultimate}.

\section{Architektury odometrii wizyjnej i SLAM-u dla kamer klatkowych}
\label{sec:literatura_klasyczny_slam}

% PTAM i S-PTAM: rozdzielenie śledzenia od mapowania, klatki kluczowe, lokalna mapa i skala stereo.
% ORB-SLAM i ORB-SLAM2: cechy ORB, relokalizacja, rozpoznawanie miejsc i korekta pętli.
% Na końcu wskazać elementy ważne dla własnego systemu.

\subsection{PTAM i S-PTAM}
\label{subsec:literatura_ptam_sptam}

\textbf{PTAM} (ang. \textit{Parallel Tracking and Mapping}) jest jednokamerowym, opartym na punktach charakterystycznych (ang. \textit{feature-based}) systemem przeznaczonym do pracy w~niewielkich przestrzeniach rozszerzonej rzeczywistości~\cite{klein2007ptam}. Jego najważniejszym wkładem jest rozdzielenie śledzenia i~mapowania na dwa równolegle działające wątki. Pierwszy estymuje bieżącą pozę kamery względem mapy punktów 3D, natomiast drugi dodaje klatki kluczowe, trianguluje nowe punkty i~wykonuje optymalizację wiązki. Pozwala to prowadzić śledzenie z~częstotliwością klatek obrazu, podczas gdy bardziej kosztowna obliczeniowo aktualizacja mapy może być wykonywana niezależnie. Konfiguracja jednokamerowa nie zapewnia bezwzględnej skali, a~system jest przeznaczony przede wszystkim do niewielkich, względnie statycznych i~dostatecznie teksturowanych scen~\cite{klein2007ptam}.

\textbf{S-PTAM} (ang. \textit{Stereo Parallel Tracking and Mapping}) rozwija koncepcję równoległego śledzenia i~mapowania znaną z~PTAM, wykorzystując zsynchronizowaną parę kamer stereo~\cite{pire2017sptam}. Ogólną architekturę systemu oraz przepływ informacji pomiędzy jego głównymi modułami przedstawiono na rys.~\ref{fig:sptam_arch}.

\begin{figure}[!ht]
\centering
\includegraphics[width=0.75\linewidth]{literature_s_ptam_architecture.png}
\caption{Schemat architektury systemu S-PTAM. Źródło: \cite{pire2017sptam}.}
\label{fig:sptam_arch}
\end{figure}

Korespondencje pomiędzy obrazami lewej i~prawej kamery umożliwiają triangulację punktów już podczas inicjalizacji mapy oraz odtworzenie rzeczywistej skali sceny. Wątek śledzenia estymuje bieżącą pozę kamery na podstawie dopasowań pomiędzy punktami mapy 3D a~cechami obserwowanymi w~aktualnej parze obrazów, natomiast działający równolegle wątek mapowania przetwarza klatki kluczowe, dodaje nowe punkty mapy oraz wykonuje lokalną optymalizację wiązki. Dzięki rozdzieleniu obu zadań śledzenie może działać z~częstotliwością napływających obrazów, podczas gdy aktualizacja mapy wykonywana jest z~mniejszą częstotliwością~\cite{pire2017sptam}.

System obejmuje również moduł wykrywania i~domykania pętli. Po rozpoznaniu i~weryfikacji ponownej obserwacji wcześniej odwiedzonego miejsca do grafu pozy dodawane jest nowe ograniczenie, a~jego optymalizacja pozwala zmniejszyć dryf narastający podczas estymacji trajektorii~\cite{pire2017sptam}. Autorzy analizują ponadto wpływ zastosowanych detektorów i~deskryptorów cech na dokładność oraz koszt obliczeniowy systemu.

W~benchmarku KITTI S-PTAM uzyskał średni względny błąd translacji (ang. \textit{average relative translation error}, $t_{\mathrm{rel}}$) równy $1{,}19\%$ oraz średni względny błąd rotacji (ang. \textit{average relative rotation error}, $r_{\mathrm{rel}}$) równy $0{,}0025^\circ/\mathrm{m}$. Dla obrazów stereo o~rozdzielczości $1344\times391$ pikseli, rejestrowanych z~częstotliwością $10\,\mathrm{Hz}$, średni czas śledzenia wyniósł $55{,}22\,\mathrm{ms}$ na laptopie z~procesorem Intel Core i7 $2{,}8\,\mathrm{GHz}$~\cite{geiger2012kitti,pire2017sptam}.


Z~obu prac przyjęto ogólny podział na szybkie śledzenie oraz wolniejsze
zarządzanie mapą zbudowaną z~klatek kluczowych i~punktów 3D. S-PTAM jest
szczególnie bliski projektowi ze względu na przepływ
\textit{stereo -- triangulacja -- lokalna mapa -- estymacja pozy -- graf}.

\subsection{ORB-SLAM}
\label{subsec:literatura_orbslam}

\textbf{ORB-SLAM} wykorzystuje cechy pochodzące z algorytmu ORB we wszystkich podstawowych zadaniach systemu: śledzeniu, mapowaniu, relokalizacji i~domykaniu pętli~\cite{murartal2015orbslam}. Jego trzy główne wątki odpowiadają za estymację bieżącej pozy kamery, lokalne rozwijanie i~optymalizację mapy oraz wykrywanie i~korektę pętli. Mapa zawiera klatki kluczowe, punkty 3D oraz graf współwidoczności, który pozwala ograniczać śledzenie i~mapowanie do lokalnego fragmentu mapy. Kandydaci na domknięcie pętli są wybierani na podstawie podobieństwa wyglądu, a~następnie poddawani weryfikacji geometrycznej przed zaakceptowaniem pętli i~korektą grafu. Ograniczeniem konfiguracji jednokamerowej jest brak bezwzględnej skali oraz możliwość jej dryfu podczas estymacji trajektorii~\cite{murartal2015orbslam}.

\textbf{ORB-SLAM2} rozwija tę architekturę, dodając obsługę kamer stereo i~RGB-D~\cite{murartal2017orbslam2}. W~przypadku stereo system rozróżnia punkty bliskie i~dalekie. Punkty bliskie, dla których głębokość może być wiarygodnie wyznaczona z~pojedynczej pary obrazów, są bezpośrednio triangulowane i~dostarczają informacji o~metrycznej skali, translacji i~rotacji. Punkty dalekie dostarczają słabszej informacji o~translacji i~skali, dlatego są triangulowane dopiero po uzyskaniu obserwacji z~wielu widoków~\cite{murartal2017orbslam2}.

Na 11 sekwencjach KITTI, zarejestrowanych z~częstotliwością $10\,\mathrm{Hz}$ i~po rektyfikacji mających rozdzielczość $1240\times376$ pikseli, średni względny błąd translacji (ang. \textit{average relative translation error}, $t_{\mathrm{rel}}$) ORB-SLAM2 był na większości przebiegów mniejszy niż $1\%$. Dla sekwencji 00 wyniósł $0{,}70\%$, a~bezwzględny pierwiastek błędu średniokwadratowego translacji (ang. \textit{absolute translation Root Mean Square Error}, $t_{\mathrm{abs}}$) był równy $1{,}3\,\mathrm{m}$~\cite{geiger2012kitti,murartal2017orbslam2}.

Istotnym wzorcem wspólnym dla obu systemów jest wykorzystanie lokalnej mapy podczas śledzenia, selekcja klatek kluczowych oraz rozdzielenie wyszukiwania kandydatów na podstawie podobieństwa wizualnego od ich późniejszej weryfikacji geometrycznej. Taka organizacja stanowi również użyteczny punkt odniesienia dla systemów wykorzystujących obrazy tworzone ze zdarzeń, pod warunkiem możliwości niezawodnego wykrywania i~dopasowywania na nich cech.


\section{Odometria wizyjna na podstawie danych z kamer zdarzeniowych w układzie stereo}
\label{sec:literatura_event_vo}

% ESVO i ESVO2: powierzchnie czasu, śledzenie, mapowanie, skala stereo i rozszerzenie o IMU.
% Krótko podać wymagania i różnice względem własnego systemu.

\subsection{ESVO}
\label{subsec:literatura_esvo}

\textbf{ESVO} (ang. \textit{Event-based Stereo Visual Odometry}) jest bezpośrednią metodą odometrii wykorzystującą stereoskopowy układ skalibrowanych kamer zdarzeniowych~\cite{zhou2021esvo}. Zdarzenia reprezentowane są za pomocą powierzchni czasu (ang. \textit{time surface}), której wartości zależą od chwili wystąpienia ostatniego zdarzenia w~danym pikselu. Metoda nie opiera się na klasycznych deskryptorach i~ich dopasowywaniu, lecz wykorzystuje bezpośrednio czasowo-przestrzenną strukturę zdarzeń.

Moduł mapowania estymuje odwrotną głębokość przez maksymalizację zgodności czasowo-przestrzennej pomiędzy obserwacjami obu kamer, a~kolejne estymaty są łączone w~lokalną, półgęstą mapę. Moduł śledzenia wyznacza natomiast pozę kamery przez rejestrację punktów tej mapy względem bieżącej powierzchni czasu~\cite{zhou2021esvo}. Znana geometria układu stereo umożliwia odtworzenie głębokości w~skali metrycznej.

Na sekwencjach \texttt{rpg\_bin}~\cite{zhou2018stereoreconstruction} o~rozdzielczości $240\times180$ pikseli i~\texttt{upenn\_flying3}~\cite{zhu2018mvsec} o~rozdzielczości $346\times260$ pikseli system uzyskał wartość RMS bezwzględnego błędu trajektorii (ang. \textit{Absolute Trajectory Error}, ATE) równą odpowiednio $2{,}8\,\mathrm{cm}$ i~$11{,}1\,\mathrm{cm}$. Implementacja działała w~czasie rzeczywistym na laptopie z~procesorem Intel Core i7-8750H, przy czym moduł mapowania pracował zwykle z~częstotliwością $20\,\mathrm{Hz}$~\cite{zhou2021esvo}.

\textbf{ESVO2} rozwija tę architekturę do stereoskopowej odometrii wizyjno-inercyjnej~\cite{niu2025esvo2}. W~module mapowania wprowadzono efektywniejsze wybieranie punktów konturów na podstawie lokalnej dynamiki zdarzeń oraz połączono estymaty głębokości uzyskiwane ze stereo czasowego i~stereo statycznego. Preintegracja pomiarów IMU dostarcza wstępnego oszacowania ruchu dla modułu śledzenia, natomiast backend aktualizuje oszacowania prędkości liniowej oraz biasów IMU. Informacja inercyjna ogranicza wskazaną przez autorów degenerację śledzenia dotyczącą składowych obrotu \textit{pitch} i~\textit{yaw}~\cite{niu2025esvo2}.

ESVO2 oceniono na pięciu publicznych zbiorach z~kamerami o~rozdzielczościach od $240\times180$ do $1280\times720$ pikseli. Autorzy wykazali działanie w~czasie rzeczywistym na standardowym procesorze dla danych o~rozdzielczości do $640\times480$ pikseli~\cite{niu2025esvo2}.

ESVO i~ESVO2 wykorzystują bezpośrednio strukturę czasową zdarzeń i~nie wymagają klasycznego dopasowywania deskryptorów. Oba systemy realizują przede wszystkim lokalną estymację ruchu i~mapowanie, bez modułów globalnej relokalizacji i~domykania pętli.


\section{Rozpoznawanie miejsc w danych zdarzeniowych}
\label{sec:literatura_event_place}

% Odróżnić rozpoznanie miejsca, relokalizację i domknięcie pętli.
% Opisać metodę Milforda i zaznaczyć, że kandydat wymaga weryfikacji geometrycznej.

\textbf{Rozpoznawanie miejsca} (ang. \textit{visual place recognition}, VPR) polega na odnalezieniu wcześniejszej obserwacji odpowiadającej aktualnie obserwowanemu miejscu. Sam wynik podobieństwa nie określa jednak względnej transformacji przestrzennej, dlatego w~systemach SLAM rozpoznane miejsce stanowi zwykle jedynie kandydata do relokalizacji lub domknięcia pętli, wymagającego dodatkowej weryfikacji geometrycznej~\cite{murartal2015orbslam}.

Jedną z~pierwszych prób wykorzystania kamer zdarzeniowych do tego zadania przedstawili Milford i~współautorzy~\cite{milford2015towards}. Zdarzenia grupowano w~stałych oknach czasowych, tworząc niskorozdzielcze obrazy wejściowe dla algorytmu SeqSLAM, który porównuje sekwencje kolejnych obserwacji. Eksperymenty wykazały spadek skuteczności przy dużych różnicach prędkości przejścia. Wstępne testy wskazały również, że zwiększenie długości okna akumulacji i~pominięcie polaryzacji zdarzeń może poprawić wyniki, choć wariant ten nie został kompleksowo oceniony~\cite{milford2015towards}.

Metodę połączono z~uproszczoną odometrią i~grafową reprezentacją mapy, uzyskując półmetryczną, topologicznie poprawną mapę trasy o~długości \SI{2.7}{\kilo\metre}. Odometria zakładała jednak ruch na płaszczyźnie, a~translację przybliżano na podstawie zliczania kroków operatora, dlatego nie był to ogólny system SLAM dla swobodnego ruchu kamery~\cite{milford2015towards}.

Późniejsze badania potwierdziły możliwość rozpoznawania miejsc bezpośrednio na podstawie strumienia zdarzeń. Fischer i~Milford pokazali, że wystarczająca może być nawet liczba zdarzeń w~niewielkim zbiorze wybranych pikseli, a~reprezentacje zawierające stałą liczbę zdarzeń zwiększają odporność na zmianę prędkości ruchu~\cite{fischer2022eventvpr}. Wyniki te pokazują, że sposób grupowania i~reprezentacji zdarzeń ma istotny wpływ na skuteczność rozpoznawania miejsc.

\section{Systemy wizyjno-inercyjne i wielosensorowe}
\label{sec:literatura_vio}

% Krótko zdefiniować pełne VIO.
% Opisać zdarzeniowe VIO Rebecqa i UltimateSLAM: kompensację ruchu, śledzenie cech oraz połączenie zdarzeń, klatek i IMU.

\textbf{Odometrię wizyjno-inercyjną} (ang. \textit{Visual-Inertial Odometry}, VIO) można traktować jako wspólną estymację stanu na podstawie obserwacji wizyjnych i~pomiarów IMU. Estymowany stan obejmuje zwykle pozę, prędkość oraz obciążenia pomiarowe czujników inercyjnych. Przykładem rozwiniętego systemu tego typu jest \textbf{ORB-SLAM3}, który rozszerza cechowy SLAM o~ściśle sprzężoną estymację wizyjno-inercyjną~\cite{campos2021orbslam3}.

\subsection{Zdarzeniowe VIO Rebecqa}
\label{subsec:event_vio}

Rebecq i~współautorzy zaproponowali system VIO wykorzystujący jedną kamerę zdarzeniową i~pełną jednostkę IMU, bez korzystania z~klasycznych klatek intensywności z DVS (ang. \textit{Dynamic Vision Sensor}, czyli kamery zdarzeniowej)~\cite{rebecq2017eventvio}. Lokalizacja zdarzeń jest kompensowana ze względu na  ruch własny kamery i~przenoszone są one do wspólnego czasu odniesienia z~użyciem bieżącej estymaty ruchu oraz informacji o~głębokości. Następnie tworzone są klatki zdarzeniowe. Na ich podstawie wykrywane są cechy FAST, śledzone metodą Lucas--Kanade i~triangulowane do lokalnej mapy.

Logika oparta na klatkach kluczowych łączy obserwacje wizyjne z~pomiarami inercyjnymi i~wspólnie estymuje pozy, prędkość oraz obciążenia pomiarowe IMU~\cite{rebecq2017eventvio}.

W~eksperymentach na danych o~rozdzielczości $240\times180$ pikseli średni błąd położenia (ang. \textit{Mean Position Error}) stanowił $0{,}69\%$ przebytej drogi dla sekwencji \texttt{boxes\_6dof} i~$0{,}92\%$ dla \texttt{hdr\_boxes}. Całkowity czas przetwarzania jednego okna zdarzeń wyniósł $8{,}23\,\mathrm{ms}$ na laptopie z~procesorem Intel Core i5-4278U i~$20\,\mathrm{ms}$ na platformie Odroid XU4~\cite{rebecq2017eventvio}.

\subsection{UltimateSLAM}
\label{subsec:ultimate_slam}

\textbf{UltimateSLAM} rozszerza zdarzeniowe VIO Rebecqa o~równoległe wykorzystanie zwykłych obrazów intensywności~\cite{vidal2018ultimate}. Zarówno na obrazach zdarzeniowych skompensowanych ruchem, jak i~na zwykłych obrazach wykrywane są narożniki FAST, a~ich położenie jest następnie śledzone metodą Lucas--Kanade. Po wiarygodnej triangulacji punkty 3D oraz pomiary IMU są wykorzystywane we wspólnej nieliniowej optymalizacji stanu opartej na klatkach kluczowych. Przykładowe wyniki śledzenia dla obu rodzajów obrazów pokazano na rysunku~\ref{fig:ultimate_slam_tracks}~\cite{vidal2018ultimate}.

\begin{figure}[!ht]
    \centering
    \includegraphics[width=0.88\linewidth,trim=0 18 0 0,clip]{literature_ultimate_slam_feature_tracks.png}
    \caption{Ślady cech w~systemie UltimateSLAM na zwykłych obrazach (górny rząd) i~obrazach zdarzeniowych skompensowanych ruchem (dolny rząd), dla różnych warunków ruchu i~oświetlenia. Rysunek z~publikacji~\cite{vidal2018ultimate}.}
    \label{fig:ultimate_slam_tracks}
\end{figure}

Oba źródła obrazu wzajemnie się uzupełniają. Zdarzenia zachowują użyteczność przy szybkim ruchu i~trudnym oświetleniu, gdy zwykłe obrazy mogą być rozmyte lub niedoświetlone. Z~kolei przy bardzo małym ruchu liczba zdarzeń maleje, dlatego zwykłe obrazy umożliwiają stabilniejsze śledzenie cech~\cite{vidal2018ultimate}. System wymaga wspólnej kalibracji kamery zdarzeniowej, kamery klatkowej i~IMU oraz ich zgodności czasowej.

Na zbiorze Event Camera Dataset~\cite{mueggler2017eventdataset} o~rozdzielczości $240\times180$ pikseli, na podstawie średniego błędu położenia (ang. \textit{Mean Position Error}), autorzy odnotowali średni wzrost dokładności położenia o~$130\%$ względem wariantu wykorzystującego wyłącznie zdarzenia i~IMU oraz o~$85\%$ względem wariantu wykorzystującego obrazy klatkowe i~IMU~\cite{vidal2018ultimate}.

Publikacja koncentruje się na lokalnej estymacji stanu wizyjno-inercyjnego i~nie opisuje mechanizmu globalnego rozpoznawania miejsc ani domykania pętli.


\section{Wnioski z przeglądu literatury i wybór inspiracji projektowych}
\label{sec:wnioski_literatura}

Na podstawie analizowanych systemów przyjęto rozwiązanie, wykorzystujące punkty charakterystyczne. Położenie tych punktów w\~przestrzeni jest wyznaczane przez triangulację odpowiadających sobie obserwacji z\~kamer w układzie stereo. Znana geometria układu pozwala  uzyskać metryczną skalę translacji.

Dane zdarzeniowe nie są przetwarzane bezpośrednio jako pojedyncze zdarzenia, lecz grupowane w~reprezentacje obrazowe. Pozwala to zastosować klasyczne metody wykrywania i~śledzenia punktów, podobnie jak w~zdarzeniowych systemach Rebecqa i~UltimateSLAM~\cite{rebecq2017eventvio,vidal2018ultimate}. Z~tych prac zaczerpnięto również ideę kompensowania wpływu ruchu przy tworzeniu reprezentacji zdarzeń, przy czym wykorzystano do tego informację o~obrocie z~żyroskopu.

W~części odpowiedzialnej za mapowanie przyjęto klatki kluczowe, punkty 3D oraz lokalne wykorzystanie mapy, zgodnie z~rozwiązaniami stosowanymi w~S-PTAM i~systemach z~rodziny ORB-SLAM~\cite{pire2017sptam,murartal2015orbslam,murartal2017orbslam2}.

W~efekcie przyjęto zestaw rozwiązań łączący stereoskopową rekonstrukcję, śledzenie punktów na reprezentacjach zdarzeniowych, wspomaganie informacją z~żyroskopu oraz geometrycznie weryfikowane domykanie pętli. Szczegóły ich realizacji przedstawiono w~kolejnym rozdziale.

    \chapter{Projekt systemu}
\label{cha:projekt_systemu}

Projekt systemu zaprojektowano z~myślą o~zadaniach konkursowych EvSLAM i~M3ED. Zbiory danych różnią się sposobem organizacji informacji, strukturą pomiarów oraz zakresem dostępnej kalibracji, dlatego wymagają odrębnych etapów wczytania i~wstępnego przygotowania. Po ujednoliceniu danych wejściowych dalsze przetwarzanie realizuje wspólny algorytm estymacji ruchu i~budowy mapy.

W~rozdziale przedstawiono wymagania wynikające z~obu konkursów, przyjęty zestaw czujników oraz założenia projektowe systemu. Następnie opisano strukturę rozwiązania, przepływ danych pomiędzy modułami oraz najważniejsze elementy odpowiedzialne za estymację ruchu, mapowanie i~korektę trajektorii.


\section{Założenia wynikające z konkursów}
\label{sec:zalozenia_systemu}

\subsection{EvSLAM Challenge}
\label{subsec:projekt_evslam}

EvSLAM Challenge dotyczy estymacji stanu podczas szybkich manewrów. Udostępnione dane obejmują dwa strumienie z kamer zdarzeniowych DVXplorer, pomiary IMU oraz obrazy RGB, monochromatyczne obrazy z kamer podczerwieni i mapy głębi rejestrowane przez kamerę Intel RealSense D435. Kamery zdarzeniowe mają rozdzielczość $640\times480$ i~bazę $10\mathrm{cm}$, natomiast IMU pracuje z~częstotliwością $800\mathrm{Hz}$~\cite{evslam2025data}. Rozmieszczenie czujników na platformie przedstawiono na rysunku~\ref{fig:evslam_sensors}. Konkurs przewiduje wariant wykorzystujący wyłącznie kamery zdarzeniowe jako źródło informacji wizyjnej oraz wariant łączący zdarzenia z~monochromatycznymi obrazami z kamery Intel RealSense D435~\cite{evslam2025challenge}.

\begin{figure}[!ht]
\centering
\includegraphics[width=0.9\linewidth]{evslam_sensors.png}
\caption{Zestaw czujników wykorzystywany do rejestracji danych EvSLAM na platformie czterowirnikowego drona. Źródło: \cite{evslam2025data}.}
\label{fig:evslam_sensors}
\end{figure}

Dla każdej sekwencji należy wyznaczyć w~zadanych chwilach pozycję i~orientację lewej kamery oraz jej prędkość liniową wyrażoną w~układzie kamery. Ocena uwzględnia bezwzględny błąd trajektorii (ang. \textit{Absolute Trajectory Error}, ATE) oraz pole pod krzywą współczynnika sukcesu dla kolejnych progów względnego błędu prędkości (ang. \textit{Area Under the Curve}, AUC). Pierwsza z tych metryk jest brana pod uwagę wyłącznie w przypadku, w którym druga (AUC), jest równa dal dwóch różnych algorytmów.

\subsection{M3ED SLAM Challenge}
\label{subsec: project_m3ed}

Zbiór M3ED został zarejestrowany z~wykorzystaniem rozbudowanego zestawu czujników obejmującego stereoskopową parę kamer zdarzeniowych Prophesee EVK4 o~rozdzielczości $1280\times720$, stereoskopową parę kamer rejestrujących obrazy w~skali szarości, kamerę RGB, jednostkę IMU oraz 64-kanałowy skaner laserowy~\cite{m3ed2025data}. Dane pochodzą z~różnych platform, obejmujących samochód, bezzałogowy statek powietrzny oraz robota kroczącego, dzięki czemu zawierają ruch o~znacznie różniącej się dynamice i~charakterze.

W~M3ED SLAM Challenge wykorzystanie czujników jest ograniczone do dwóch wariantów. W~wariancie \textit{Event (+ IMU)} pozycja może być wyznaczana na podstawie jednej lub dwóch kamer zdarzeniowych, opcjonalnie z~wykorzystaniem IMU. Wariant \textit{Event + Mono (+ IMU)} dopuszcza dodatkowo obrazy z~jednej kamery globalnej migawki (ang. \textit{global shutter camera})~\cite{m3ed2025challenge}.

Konkurs obejmuje trzy sekwencje testowe: jazdę samochodem w~środowisku miejskim, szybki lot UAV oraz ruch robota kroczącego. Dla wskazanych znaczników czasu należy wyznaczyć pozycję i~orientację referencyjnej kamery zdarzeniowej. Dokładność jest oceniana za pomocą bezwzględnego błędu pozy (ang. \textit{Absolute Pose Error}, APE) względem trajektorii referencyjnej~\cite{m3ed2025challenge}. W~przeciwieństwie do EvSLAM wynik nie obejmuje prędkości liniowej -- wymaganym rezultatem jest wyłącznie trajektoria położenia robota.


\subsection{Wybór konfiguracji czujników}
\label{subsec:wybor_czujnikow}

W~pracy przyjęto wspólny dla obu konkursów wariant wykorzystujący
\textbf{dwie kamery zdarzeniowe i~IMU}. Informacja wizyjna pochodzi wyłącznie
ze zdarzeń, dzięki czemu system odpowiada wariantom niewymagającym obrazów z kamer klatkowych. Konfiguracja stereo została wybrana z~dwóch powodów:

\begin{itemize}
    \item znana baza kamer pozwala od razu wyznaczać głębokość i~metryczną skalę translacji;
    \item dwie obserwacje tej samej chwili ograniczają niejednoznaczność geometrii charakterystyczną dla pojedynczej kamery.
\end{itemize}

Z~IMU wykorzystywana jest prędkość kątowa. Żyroskop pomaga ograniczyć wpływ
szybkich obrotów na reprezentację zdarzeń oraz może dostarczyć pomocniczą
informację o~zmianie orientacji, podobnie jak w~systemach łączących zdarzenia
z~pomiarami inercyjnymi~\cite{rebecq2017eventvio,vidal2018ultimate}.
Akcelerometr nie jest uwzględniany, dlatego opracowane rozwiązanie nie jest
pełnym systemem VIO: nie estymuje obciążeń pomiarowych IMU
ani wspólnego stanu wizyjno-inercyjnego. Jest to świadome ograniczenie zakresu,
pozwalające skupić projekt na geometrii zdarzeniowego stereo VO i~SLAM.

\section{Cele i organizacja systemu}
\label{sec:architektura_systemu}

System zaprojektowano tak, aby po przygotowaniu danych wejściowych zbiory EvSLAM i~M3ED mogły być przetwarzane przez wspólną część algorytmu. Najważniejsze cele projektowe obejmują:

\begin{itemize}
\item estymację trajektorii referencyjnej kamery zdarzeniowej w~skali metrycznej;
\item wyznaczenie prędkości liniowej z~trajektorii na potrzeby zadania EvSLAM;
\item możliwość pracy wyłącznie jako odometria wizyjna lub z~dodatkowym wykorzystaniem mapy;
\item relokalizację po utracie śledzenia oraz korektę dryfu po wykryciu ponownie odwiedzonego miejsca.
\end{itemize}

Ogólną organizację systemu i~przepływ danych przedstawiono na rysunku~\ref{fig:system_architecture}. Główny tor przetwarzania prowadzi od zsynchronizowanych strumieni zdarzeń stereo do trajektorii lewej kamery i~wyznaczanej z~niej prędkości liniowej. Opcjonalne pomiary żyroskopu mogą wspierać przygotowanie obserwacji oraz estymację pozy, natomiast wariant SLAM rozszerza system o~rzadką mapę, rozpoznawanie miejsc, relokalizację i~korektę grafu pozy.

\begin{figure}[!ht]
    \centering
    \includegraphics[width=1\linewidth]{system_architecture.pdf}
    \caption{Architektura i~przepływ danych w~opracowanym systemie. Zielone pola oznaczają dane wejściowe, pomarańczowe wyniki, a~niebieskie elementy wariantu SLAM. Przerywana strzałka przedstawia warunkowe odzyskanie pozy przez relokalizację.}
    \label{fig:system_architecture}
\end{figure}

Przepływ danych w~systemie przedstawiono na rysunku~\ref{fig:system_architecture}. Zdarzenia z~obu kamer są synchronizowane i~dzielone na wspólne okna czasowe, a~następnie poddawane filtracji oraz korekcji zniekształceń wywołanych obrotem kamery. Wykorzystywane są przy tym zsynchronizowane z~kamerami pomiary żyroskopu, natomiast pomiary akcelerometru nie uczestniczą w~obliczeniach. Przygotowane zdarzenia są przekształcane w~reprezentacje obrazowe i~rektyfikowane. System wykrywa na nich punkty, śledzi je w~czasie i~wyznacza ich głębokość na podstawie obserwacji z~obu kamer. Ruch kamery jest określany z~wykorzystaniem punktów 3D z~poprzedniego okna lub mapy oraz ich bieżących obserwacji; w~konfiguracji końcowej translacja pochodzi z~wizji, a~zmiana orientacji z~żyroskopu. Wybrane obserwacje rozwijają rzadką mapę, która wspiera dalszą lokalizację, rozpoznawanie wcześniej odwiedzonych miejsc i~korektę trajektorii. Wynikiem jest trajektoria lewej kamery, a~dla EvSLAM również wyznaczona na jej podstawie prędkość liniowa. Dane referencyjne nie uczestniczą w~estymacji i~są wykorzystywane wyłącznie podczas oceny wyników.

    \chapter{Implementacja systemu}
\label{cha:implementacja}

System zaimplementowano w~języku Python z~wykorzystaniem bibliotek \textit{NumPy}, \textit{OpenCV} i~\textit{SciPy}. W~rozdziale opisano obliczenia wykonywane na kolejnych etapach oraz wizualizacje wyjść z poszczególnych etapów.

\section{Odczyt i przygotowanie danych}
\label{sec:implementacja_wejscie}

Dane wejściowe pochodzą z~gotowych zbiorów EvSLAM i~M3ED. Ze względu na różne sposoby ich zapisu (pliki \textit{.bag} dla EvSLAM i \textit{.h5} dla M3ED) zastosowano osobne mechanizmy odczytu, które przekazują do dalszego przetwarzania:

\begin{itemize}
    \item dwa strumienie zdarzeń zawierających współrzędne piksela, znacznik czasu i~polaryzację;
    \item parametry wewnętrzne, modele dystorsji i~rozdzielczości obu kamer;
    \item transformację $\mathbf{T}_{C_R C_L}$ opisującą wzajemne położenie kamer;
    \item opcjonalne pomiary prędkości kątowej oraz transformacje $\mathbf{T}_{C_L I}$ i~$\mathbf{T}_{C_R I}$ z~układu IMU do układów kamer.
\end{itemize}

System odczytuje wartości kalibracji, dostarczone wraz ze zbiorem, i weryfikuje ich poprawność. Pomiary żyroskopu są porządkowane według czasu, a~próbki o~powtórzonych znacznikach są usuwane.

Zdarzenia z~obu kamer są dzielone na wspólne okna $\mathcal{W}_k=[t_{s,k},t_{e,k})$. Każdemu oknu przypisywany jest znacznik odpowiadający środkowi przedziału.

\section{Przetwarzanie zdarzeń}
\label{sec:implementacja_zdarzenia}

\subsection{Filtr aktywności tła}
\label{subsec:okna_baf}

Filtr aktywności tła (ang. \textit{background activity filter}, BAF) odrzuca zdarzenia bez dostatecznego wsparcia w~bliskim otoczeniu przestrzennym i~czasowym. Ideę takiego testu pokazano na rysunku~\ref{fig:baf_neighborhood_source}~\cite{guo2023denoising}. W~opracowanym wariancie zamiast wyszukiwania najbliższych sąsiadów zliczane są aktywne piksele w~stałym, kwadratowym otoczeniu.

\begin{figure}[!ht]
    \centering
    \includegraphics[width=0.9\textwidth]{implementation_baf_neighborhood_guo2023.png}
    \caption{Zasada stosowana w~filtrach aktywności tła. Źródło:~\cite{guo2023denoising}.}
    \label{fig:baf_neighborhood_source}
\end{figure}

Dla każdej kamery filtr przechowuje czas ostatniego zdarzenia w~każdym pikselu. Nowe zdarzenie jest zachowywane, jeżeli w~kwadratowym sąsiedztwie o~rozmiarze $n \times n$ (lub $(2r+1)\times(2r+1)$ stosując promień, który w algorytmie jest parametrem filtru) wystąpiło co najmniej $k$ zdarzeń w~czasie nie dłuższym niż $t$. W~przeciwnym przypadku zostaje uznane za odosobnione zdarzenie tła i~odrzucone. Polaryzacja nie wpływa na decyzję filtra.

W~końcowej konfiguracji EvSLAM, na podstawie prób i testów rozwojowych przyjęto $r=2$, czyli otoczenie $5\times5$ pikseli, czas $t=10\,\mathrm{ms}$ i~wymaganie co najmniej $k=5$ aktywnych sąsiadów. Wpływ filtra o takich parametrach na klatkę zdarzeniową pokazano na rysunku~\ref{fig:implementation_baf_comparison}. Dla M3ED filtr wyłączono.

\begin{figure}[!ht]
    \centering
    \begin{subfigure}[t]{0.48\textwidth}
        \centering
        \includegraphics[width=\textwidth]{implementation_evslam_seq007_baf_before.png}
        \caption{Reprezentacja bez filtracji.}
    \end{subfigure}
    \hfill
    \begin{subfigure}[t]{0.48\textwidth}
        \centering
        \includegraphics[width=\textwidth]{implementation_evslam_seq007_baf_after.png}
        \caption{Reprezentacja po filtracji BAF.}
    \end{subfigure}
    \caption{Wpływ filtra aktywności tła na reprezentację zdarzeń lewej kamery dla tego samego okna sekwencji \texttt{seq007\_test} ze zbioru EvSLAM. Niebieskie i~czerwone piksele reprezentacji wejściowych oznaczają odpowiednio zdarzenia dodatnie i~ujemne. Ponieważ szum jest stosunkowo niewielki, żeby zaobserwować efekty są widoczne dopiero po dokładniejszej obserwacji.}
    \label{fig:implementation_baf_comparison}
\end{figure}

\subsection{Kompensacja ruchu obrotowego}
\label{subsec:kompensacja_ruchu}
Zastosowane rozwiązanie oparto na metodzie tworzenia obrazów zdarzeniowych skompensowanych ruchem, przedstawionej w pracy \cite{rebecq2017eventvio} oraz wykorzystanej później w~systemie UltimateSLAM~\cite{vidal2018ultimate}. W~metodach tych zdarzenia są przenoszone do wspólnej chwili na podstawie ruchu wyznaczonego z~pomiarów inercyjnych. W~opracowanym systemie zastosowano uproszczenie ograniczające kompensację do obrotu zmierzonego przez żyroskop. Dzięki temu nie jest wymagana znajomość głębokości poszczególnych zdarzeń.

Kompensacja ruchu przenosi zdarzenia z~różnych chwil do wspólnego czasu odniesienia, którym jest środek okna. Zdarzenia są dzielone na stałą liczbę przedziałów czasowych (przyjęto $32$). Dla każdego przedziału przyjmowany jest średni czas należących do niego zdarzeń, dzięki czemu nie trzeba wyznaczać osobnego obrotu dla każdego zdarzenia.

Obrót pomiędzy średnim czasem przedziału a~czasem odniesienia jest obliczany przez całkowanie pomiarów żyroskopu. Próbki na granicach tego zakresu są wyznaczane przez interpolację. Dla dwóch kolejnych chwil $t_i$ i~$t_{i+1}$ przyrost obrotu wynosi
\begin{equation}
    \Delta\boldsymbol{\theta}_i=
    \frac{\boldsymbol{\omega}_i+\boldsymbol{\omega}_{i+1}}{2}\Delta t_i,
    \qquad
    \Delta t_i=t_{i+1}-t_i,
    \label{eq:implementation_gyro_integration}
\end{equation}
gdzie $\boldsymbol{\omega}_i$ i~$\boldsymbol{\omega}_{i+1}$ są pomiarami prędkości kątowej, a~$\Delta\boldsymbol{\theta}_i$ jest wektorem obrotu. Kierunek tego wektora wyznacza oś obrotu, natomiast jego długość odpowiada kątowi wyrażonemu w~radianach. Wektor jest przekształcany w~macierz obrotu, a~kolejne przyrosty są składane zgodnie z~kolejnością pomiarów~\cite{forster2015preintegration}. System nie estymuje obciążenia żyroskopu. Otrzymany obrót jest przenoszony z~układu IMU do układu odpowiedniej kamery z~wykorzystaniem kalibracji między czujnikami.

Wektor zapisuje obrót: kąt dla każdej osi. Kąt obrotu jest równy długości tego wektora:
\begin{equation}
    \alpha_i=\left\|\Delta\boldsymbol{\theta}_i\right\|,
    \qquad
    \mathbf{a}_i=\frac{\Delta\boldsymbol{\theta}_i}{\alpha_i},
    \label{eq:implementation_axis_angle}
\end{equation}
gdzie $\mathbf{a}_i=(a_x,a_y,a_z)^T$ jest jednostkowym wektorem osi obrotu. Na jego podstawie tworzona jest macierz
\begin{equation}
    \mathbf{S}_i=
    \begin{bmatrix}
        0 & -a_z & a_y\\
        a_z & 0 & -a_x\\
        -a_y & a_x & 0
    \end{bmatrix},
    \qquad
    \Delta\mathbf{R}_i=
    \mathbf{I}+\sin(\alpha_i)\mathbf{S}_i+
    \left(1-\cos(\alpha_i)\right)\mathbf{S}_i^2,
    \label{eq:implementation_rodrigues}
\end{equation}
gdzie $\mathbf{I}$ oznacza macierz jednostkową. Jest to wzór Rodriguesa przekształcający obrót opisany osią i~kątem w~macierz obrotu~\cite{blanco2021se3}. Dla zerowego kąta przyjmowana jest macierz jednostkowa. Macierze kolejnych przyrostów są następnie składane zgodnie z~kolejnością pomiarów.

Dla zdarzenia zarejestrowanego w~pikselu $\bar{\mathbf{u}}_k=(x_k,y_k,1)^T$ nowe położenie jest wyznaczane jako
\begin{equation}
    \bar{\mathbf{u}}'_k\sim
    \mathbf{K}\mathbf{R}_{C_r C_k}\mathbf{K}^{-1}\bar{\mathbf{u}}_k,
    \label{eq:implementation_motion_compensation}
\end{equation}
gdzie $\mathbf{K}$ jest macierzą parametrów wewnętrznych kamery, a~$\mathbf{R}_{C_r C_k}$ przekształca kierunek obserwacji z~chwili zdarzenia $t_k$ do czasu odniesienia $t_r$. Najpierw współrzędne piksela są więc zamieniane na kierunek obserwacji za pomocą $\mathbf{K}^{-1}$, następnie ten kierunek jest obracany, a~macierz $\mathbf{K}$ ponownie przekształca go na współrzędne obrazu. Symbol $\sim$ oznacza równość z~dokładnością do skali; współrzędne piksela otrzymuje się po podzieleniu dwóch pierwszych elementów wyniku przez element trzeci.
Otrzymane współrzędne są zaokrąglane do najbliższego piksela, a~zdarzenia przeniesione poza obraz zostają odrzucone. Kompensowany jest wyłącznie obrót kamery. Uwzględnienie translacji wymagałoby znajomości głębokości każdego zdarzenia, która na tym etapie nie jest jeszcze dostępna.
Po kompensacji widocznej na rysunku~\ref{fig:implementation_motion_compensation} krawędzie są węższe, a~przeciwne polaryzacje lepiej się pokrywają. Jeżeli próbki IMU nie obejmują całego przedziału, okno jest przetwarzane bez kompensacji. Można zaobserwować, że obrót własny robota, który powodował część rozmycia ruchu, został skutecznie skompensowany. Rozmycie ruchu wynikające z translacji robota pozostaje.

\begin{figure}[!ht]
    \centering
    \includegraphics[width=0.9\textwidth]{implementation_evslam_seq007_rotation_compensation.png}
    \caption{Reprezentacje zdarzeń lewej i~prawej kamery przed kompensacją obrotu oraz po jej wykonaniu dla fragmentu sekwencji \texttt{seq007\_test} ze zbioru EvSLAM \cite{evslam2025data}.}
    \label{fig:implementation_motion_compensation}
\end{figure}

\subsection{Agregacja zdarzeń}
\label{subsec:implementacja_agregacja}

Po filtracji i~opcjonalnej kompensacji każde okno jest przekształcane w~ośmiobitową reprezentację obrazową. Zaimplementowano standardową klatkę zdarzeniową opisaną równaniem~\eqref{eq:standard_event_frame} oraz wykładniczo zanikającą klatkę zdarzeniową z~równania~\eqref{eq:exponential_event_frame}. W~obu wariantach wykorzystywane jest ostatnie zdarzenie zarejestrowane w~każdym pikselu.

Wartość 127 oznacza piksel bez zdarzenia. W~wariancie standardowym zdarzenie dodatnie otrzymuje wartość 255, a~ujemne 0. W~wariancie wykładniczym odchylenie od wartości neutralnej maleje wraz z~wiekiem zdarzenia. W konfiguracjach końcowych stosowany jest wariant wykładniczy ze stałą zaniku $\tau$, przyjmowaną jako równa czasowi agregacji ramek. Porównanie obu zaimplementowanych reprezentacji z~odpowiadającą im klatką intensywności przedstawiono na rysunku~\ref{fig:implementation_event_representations}.

\begin{figure}[!ht]
    \centering
    \begin{subfigure}[t]{0.31\textwidth}
        \centering
        \includegraphics[width=\textwidth]{implementation_evslam_seq007_intensity_frame.png}
        \caption{Klatka RGB z tradycyjnej kamery.}
    \end{subfigure}
    \hfill
    \begin{subfigure}[t]{0.31\textwidth}
        \centering
        \includegraphics[width=\textwidth]{implementation_evslam_seq007_standard_event_frame.png}
        \caption{Standardowa klatka zdarzeniowa. $t=12ms$.}
    \end{subfigure}
    \hfill
    \begin{subfigure}[t]{0.31\textwidth}
        \centering
        \includegraphics[width=\textwidth]{implementation_evslam_seq007_exponential_event_frame.png}
        \caption{Wykładniczo zanikająca klatka zdarzeniowa. $t=12ms$; $\tau=8ms$.}
    \end{subfigure}
    \caption{Porównanie klatki kamery konwencjonalnej i~dwóch zaimplementowanych reprezentacji tego samego fragmentu danych EvSLAM. W~wizualizacji kolor niebieski oznacza zdarzenia dodatnie, a~czerwony zdarzenia ujemne.}
    \label{fig:implementation_event_representations}
\end{figure}

\subsection{Rektyfikacja}
\label{subsec:implementacja_rektyfikacja}

Mapy rektyfikacji są wyznaczane raz, z~parametrów obu kamer i~transformacji $\mathbf{T}_{C_R C_L}$. Obrazy wynikowe mają wspólne położenie punktów głównych i~obejmują wspólne pole widzenia. Reprezentacje są przemapowywane z~interpolacją najbliższego sąsiada, aby nie tworzyć pośrednich wartości między rzadkimi strukturami zdarzeniowymi. Na rysunku~\ref{fig:implementation_rectification} odpowiadające sobie fragmenty przecinają te same poziome linie.

\begin{figure}[!ht]
    \centering
    \includegraphics[width=\textwidth]{implementation_evslam_seq007_stereo_rectification.png}
    \caption{Zrektyfikowana para reprezentacji zdarzeń z~zaznaczonymi poziomymi liniami pomocniczymi.}
    \label{fig:implementation_rectification}
\end{figure}

\section{Śledzenie cech i wyznaczanie głębokości}
\label{sec:implementacja_sledzenie_stereo}

Na lewej zrektyfikowanej reprezentacji wykrywane jest do 1000 punktów GFTT~\cite{shi1994good}. Pomiędzy oknami są one śledzone piramidalną metodą Lucas--Kanade~\cite{lucas1981iterative}, z~oknem $21\times21$ pikseli i~czterema poziomami piramidy. Ślady niepoprawne lub kończące się poza obrazem są odrzucane, a~w~miejscach oddalonych o~co najmniej siedem pikseli od aktywnych cech wykrywane są nowe punkty. Wynik pokazano na rysunku~\ref{fig:implementation_tracking}.

Dopasowanie stereo również wykorzystuje metodę Lucas--Kanade. Punkt w~prawej reprezentacji musi leżeć wewnątrz obrazu, mieć dodatnią dysparycję z~zadanego zakresu i~różnić się pionowo o~nie więcej niż dwa piksele. Dla EvSLAM stosowana jest dodatkowo kontrola w~obu kierunkach: po przeniesieniu punktu z~lewego obrazu do prawego i~z~powrotem błąd nie może przekroczyć jednego piksela.

Zaakceptowane pary są triangulowane zgodnie z~równaniami~\eqref{eq:disparity} i~\eqref{eq:stereo_point}. Odrzucane są punkty o~niepoprawnych współrzędnych oraz głębokości spoza skonfigurowanego zakresu. Punkty dodawane do mapy są ograniczone do $10\,\mathrm{m}$, ponieważ głębokość odległych punktów, odpowiadających małym wartościom dysparycji, jest szczególnie wrażliwa na błąd dopasowania. Przykład zaakceptowanych korespondencji przedstawiono na rysunku~\ref{fig:implementation_stereo_depth}.

\begin{figure}[!ht]
    \centering
    \begin{subfigure}[t]{0.72\textwidth}
        \centering
        \includegraphics[width=\textwidth]{implementation_evslam_seq007_feature_tracking.png}
        \caption{Śledzenie punktów na lewej reprezentacji. Zielone odcinki łączą położenia punktów uznanych za te same w~poprzednim i~bieżącym oknie, a~żółte kropki oznaczają ich bieżące położenia. W centrum obrazu można zaobserwować błędnie śledzone punkty -- zmiana ich położenia nie jest zgodna z ruchem kamery. Wartości \texttt{tracks}, \texttt{active} i~\texttt{redetect} oznaczają odpowiednio liczbę poprawnie prześledzonych punktów, wszystkich aktywnych punktów oraz nowych punktów wykrytych po zakończeniu śledzenia.}
        \label{fig:implementation_tracking}
    \end{subfigure}
    \par\medskip
    \begin{subfigure}[t]{\textwidth}
        \centering
        \includegraphics[width=\textwidth]{implementation_evslam_seq007_stereo_depth.png}
        \caption{Dopasowania pomiędzy lewą i~prawą reprezentacją. Kropki wskazują położenia odpowiadających sobie punktów, a~linie łączą pary korespondencji. Zmiana ich koloru od żółtego do czerwonego oznacza wzrost wyznaczonej głębokości. Wartości \texttt{in}, \texttt{matched} i~\texttt{3d} oznaczają liczbę punktów przekazanych do dopasowania, zaakceptowanych korespondencji i~poprawnie wyznaczonych punktów 3D, natomiast \texttt{z\_med} jest medianą ich głębokości wyrażoną w~metrach.}
        \label{fig:implementation_stereo_depth}
    \end{subfigure}
    \caption{Wyniki pośrednie śledzenia cech i~wyznaczania rzadkiej głębokości stereo.}
    \label{fig:implementation_tracking_stereo}
\end{figure}

\section{Estymacja pozy kamery}
\label{sec:implementacja_pozy}

Układ świata $W$ jest inicjalizowany zgodnie z~położeniem lewej kamery w~pierwszym oknie. Jej początkowej pozie przypisywana jest macierz jednostkowa $\mathbf{T}_{W C_0}=\mathbf{I}$, a~dla obserwowanych punktów wyznaczane są początkowe współrzędne 3D. W~kolejnych oknach ruch kamery jest estymowany na podstawie korespondencji pomiędzy punktami 3D a~ich bieżącymi obserwacjami 2D. Najpierw wykorzystywane są punkty mapy lokalnej. Jeżeli estymacja względem mapy nie powiedzie się, system korzysta z~punktów wyznaczonych dla poprzedniej pary stereo.

Rotacja i~translacja są wyznaczane wspólnie za pomocą iteracyjnej metody PnP dostępnej w~bibliotece OpenCV~\cite{opencv_calib3d}. Obliczenia odpowiadają modelowi projekcji opisanemu równaniem~\eqref{eq:pnp}. Algorytm RANSAC służy do odrzucania błędnych korespondencji na podstawie błędu reprojekcji~\cite{fischler1981ransac}. Metryczna skala translacji wynika z~głębokości punktów otrzymanej z~układu stereo~\cite{hartley2004multiple}.

Dodano również możliwość przyjęcia rotacji wyznaczonej przez całkowanie pomiarów żyroskopu. Estymacja nadal korzysta z~tych samych korespondencji 3D--2D i~modelu projekcji, jednak rotacja zostaje ustalona, dlatego wyznaczana jest tylko translacja. Wariant ten pozwala porównać dryf orientacji wynikający z~estymacji wizyjnej z~dryfem otrzymanym z~żyroskopu. Jego wadą jest bezpośrednie przenoszenie i~akumulowanie błędów pomiaru prędkości kątowej, w~szczególności błędu systematycznego żyroskopu~\cite{forster2015preintegration}.

Niezależnie od źródła rotacji, do uruchomienia estymacji potrzeba co najmniej 15 korespondencji. Rozwiązanie musi zawierać co najmniej 15 inlierów, stanowiących nie mniej niż $25\%$ wszystkich korespondencji, a~mediana błędu reprojekcji nie może przekroczyć $2{,}5$ piksela. RANSAC wykonuje do 100 prób, z~progiem reprojekcji 3 piksele i~zadaną pewnością $0{,}999$. Odrzucane są również kroki większe niż $0{,}2\,\mathrm{m}$ lub $15^\circ$. W~wariancie pełnego PnP wynik wizyjny może być dodatkowo odrzucony, gdy różnica względem przyrostu z~żyroskopu przekracza $5^\circ$; kontrola ta nie zmienia konfiguracji końcowej, w~której rotacja jest bezpośrednio ustalana z~żyroskopu. Gdy dane inercyjne nie obejmują wymaganego przedziału czasu, system wyznacza pełną pozę za pomocą PnP.

Zaakceptowane oszacowanie jest sprowadzane do transformacji $\mathbf{T}_{C_k C_{k-1}}$. Zgodnie z~konwencją określoną w~równaniu~\eqref{eq:frame_transform} aktualizacja pozy ma postać
\begin{equation}
    \mathbf{T}_{W C_k}
    =
    \mathbf{T}_{W C_{k-1}}
    \mathbf{T}_{C_k C_{k-1}}^{-1}.
    \label{eq:implementation_pose_update}
\end{equation}
Transformacja odwrotna jest potrzebna, ponieważ $\mathbf{T}_{C_k C_{k-1}}$ przekształca współrzędne z~poprzedniego układu kamery do bieżącego. Jeżeli nie uda się uzyskać poprawnego rozwiązania, zachowywana jest ostatnia pozycja, natomiast śledzenie punktów i~triangulacja są kontynuowane.

\subsection{Klatki kluczowe i punkty mapy}\
\label{subsec}

Mapa ma strukturę opartą na klatkach kluczowych i~punktach 3D. Pierwsze przetworzone okno tworzy początkową klatkę kluczową, a~następna może zostać dodana po odstępie co najmniej pięciu okien. Nowa klatka powstaje, gdy obrót względem poprzedniej osiągnie $10^\circ$, udział aktywnych śladów powiązanych z~mapą spadnie poniżej $0{,}6$ albo stosunek przesunięcia kamery do mediany głębokości przekroczy $0{,}15$. Jeżeli żaden z~tych warunków nie zostanie wcześniej spełniony, klatka jest dodawana po 30 oknach. Podane wartości są parametrami przyjętymi w~opracowanym systemie. Dobrano je podczas testów rozwojowych, tak aby reagować na wyraźną zmianę widoku lub utratę punktów mapy, a~jednocześnie nie tworzyć nadmiarowych klatek kluczowych.

Punkt mapy jest tworzony, gdy cecha zapisana w~klatce kluczowej ma poprawną głębokość stereo oraz deskryptor ORB~\cite{rublee2011orb}. Jego położenie zostaje przekształcone do układu świata zgodnie z~równaniem~\eqref{eq:frame_transform}. Jeżeli ten sam ślad wystąpi w~kolejnej klatce kluczowej, położenie punktu jest uśredniane z~nową obserwacją. Punkt zostaje potwierdzony po drugiej obserwacji, natomiast niepotwierdzone punkty są usuwane, gdy ich początkowa klatka nie należy już do pięciu ostatnich klatek kluczowych.

Mapa lokalna zawiera punkty obserwowane w~pięciu ostatnich klatkach kluczowych. W~pierwszej kolejności wykorzystywane są zachowane powiązania śladów. Pozostałe cechy są dopasowywane na podstawie odległości Hamminga pomiędzy deskryptorami ORB. Dla odległości $d_1\leq d_2$ dwóch najlepszych kandydatów stosowany jest test ilorazowy Lowe'a~\cite{lowe2004distinctive}:
\begin{equation}
    d_1<\eta d_2,
    \label{eq:implementation_descriptor_ratio}
\end{equation}
gdzie w~opracowanym systemie $\eta=0{,}75$. Dopasowanie jest przyjmowane tylko wtedy, gdy najlepszy kandydat jest lepszy od drugiego i~znajduje się dostatecznie blisko przewidywanej reprojekcji punktu. Każda cecha obrazu i~każdy punkt mapy mogą zostać wykorzystane tylko w~jednym dopasowaniu.

Rysunek~\ref{fig:implementation_sparse_map} przedstawia mapę uzyskaną dla pełnej sekwencji \texttt{seq007\_test}. Czerwona linia łączy pozycje klatek kluczowych, natomiast kolor punktu oznacza liczbę jego obserwacji. Pełna mapa zawiera również pojedyncze punkty położone wyraźnie poza zasadniczą strukturą otoczenia, powstałe wskutek błędnych dopasowań lub triangulacji. Nie wpływają one bezpośrednio na estymowaną pozę: punkt jest wykorzystywany dopiero po dopasowaniu deskryptora, sprawdzeniu jego reprojekcji i~weryfikacji korespondencji przez RANSAC~\cite{fischler1981ransac}.

\begin{figure}[!ht]
    \centering
    \begin{subfigure}[t]{0.49\textwidth}
        \centering
        \includegraphics[width=\textwidth]{implementation_evslam_seq007_sparse_map_side_view.png}
        \caption{Widok przestrzenny.}
        \label{fig:implementation_sparse_map_seq007_3d}
    \end{subfigure}
    \hfill
    \begin{subfigure}[t]{0.49\textwidth}
        \centering
        \includegraphics[width=\textwidth]{implementation_evslam_seq007_sparse_map_top_view.png}
        \caption{Rzut z~góry na płaszczyznę $X$-$Z$.}
        \label{fig:implementation_sparse_map_seq007_topdown}
    \end{subfigure}
    \caption{Rzadka mapa uzyskana dla pełnej sekwencji \texttt{seq007\_test} -- robota poruszającego się po okręgach w zamkniętym pokoju. Czerwona linia przedstawia pozycje klatek kluczowych, a~kolor punktów oznacza liczbę ich obserwacji. Pokazano punkty utworzone w~odległości nie większej niż $5\,\mathrm{m}$ od kamery.}
    \label{fig:implementation_sparse_map}
\end{figure}


\subsection{Rozpoznawanie miejsc i relokalizacja}
\label{subsec:relocalizacja_petle}

Obserwacja wykorzystywana do rozpoznawania miejsca jest tworzona niezależnie od punktów śledzonych podczas estymacji pozy. Z~reprezentacji zdarzeń powstaje obraz aktywności bez rozróżniania polaryzacji: bezwzględna różnica względem wartości neutralnej 127 jest podwajana i~ograniczana do 255, a~następnie poddawana filtracji medianowej $3\times3$. Na tak przygotowanym obrazie wyznaczane są cechy ORB oraz odpowiadające im głębokości stereo~\cite{rublee2011orb}. Równolegle obraz jest zmniejszany do $40\times24$ pikseli, centrowany i~normalizowany, tworząc uproszczoną sygnaturę całej obserwacji. Podział na szybkie wyszukiwanie podobnych klatek i~późniejszą weryfikację geometryczną jest wzorowany na ORB-SLAM2~\cite{murartal2017orbslam2}, wykorzystano jednak własną, prostszą sygnaturę obrazu.

Iloczyn skalarny znormalizowanych sygnatur służy do wstępnego uporządkowania klatek kluczowych. Ich deskryptory ORB są następnie porównywane za pomocą odległości Hamminga i~testu ilorazowego z~równania~\eqref{eq:implementation_descriptor_ratio}. W~tym teście przyjęto $\eta=0{,}75$ dla EvSLAM i~$\eta=0{,}85$ dla M3ED. Do dalszej weryfikacji wymagane jest co najmniej 50 dopasowań, a~najwyżej trzy najlepsze klatki są sprawdzane metodą PnP z~algorytmem RANSAC~\cite{fischler1981ransac}. Transformacja jest sprawdzana w~obu kierunkach: z~punktów 3D kandydata do bieżących obserwacji 2D oraz z~bieżących punktów 3D do obserwacji zapisanych w~kandydacie.

Dopasowanie zostaje zaakceptowane, gdy co najmniej jeden kierunek spełnia progi liczby i~udziału inlierów oraz ma medianę błędu reprojekcji nie większą niż 2 piksele. Dla EvSLAM wymagane jest co najmniej 25 inlierów i~udział równy przynajmniej $0{,}4$, natomiast dla M3ED odpowiednio 20 i~$0{,}3$. Łagodniejsze wartości przyjęto jednakowo dla wszystkich sekwencji M3ED, ponieważ wcześniejsze progi odrzucały poprawnych kandydatów przed utworzeniem krawędzi pętli. Pozostałe kryteria są wspólne dla obu zbiorów. Przebieg procesu przedstawiono na rysunku~\ref{fig:place_recognition_flow}.

\begin{figure}[!ht]
    \centering
    \includegraphics[height=0.35\textheight]{implementation_place_recognition_flow.png}
    \caption{Wyszukiwanie i~weryfikacja wcześniej odwiedzonego miejsca. Zależnie od sytuacji poprawne dopasowanie służy do relokalizacji albo utworzenia krawędzi pętli.}
    \label{fig:place_recognition_flow}
\end{figure}

Relokalizacja jest podejmowana po trzech kolejnych nieudanych próbach estymacji pozy i~obejmuje całą historię klatek kluczowych. Zaakceptowana transformacja pozwala wyznaczyć bieżącą pozę względem odnalezionej klatki. Podczas domykania pętli rozpoznawanie miejsca jest wykonywane po utworzeniu klatki kluczowej, z~pominięciem kandydatów zbyt bliskich w~kolejności przetwarzania. Potwierdzone dopasowanie tworzy krawędź pętli przekazywaną do optymalizacji grafu pozy.

Przykład działania warstwy SLAM przedstawiono na rysunku~\ref{fig:implementation_slam_diagnostics_seq007}. Górny wykres pokazuje liczbę śledzonych i~aktywnych cech, korespondencji z~mapą oraz inlierów. Dolny wskazuje źródło pozy w~kolejnych oknach. Czerwone krzyżyki oznaczają okna, w~których powiązania aktywnych cech z~mapą lokalną zostały uzupełnione na podstawie deskryptorów ORB; nie oznaczają one relokalizacji. Gwiazdy wskazują zaakceptowane domknięcia pętli. W~przedstawionym przebiegu większość póz wyznaczono względem mapy lokalnej, zaakceptowano 15 pętli i~wykonano 4 relokalizacje. przez większość sekwencji pozycja była skutecznie wyznaczana względem mapy lokalnej, natomiast relokalizacja była potrzebna jedynie w~nielicznych momentach utraty śledzenia

\begin{figure}[!ht]
    \centering
    \includegraphics[width=\textwidth]{implementation_slam_diagnostics_seq007.png}
    \caption{Diagnostyka śledzenia i~mechanizmów SLAM dla pełnej sekwencji \texttt{seq007}. Górny wykres przedstawia liczby cech, korespondencji z~mapą i~inlierów, natomiast dolny pokazuje źródło estymowanej pozy, dopasowania ORB do mapy lokalnej oraz zaakceptowane domknięcia pętli.}
    \label{fig:implementation_slam_diagnostics_seq007}
\end{figure}

\subsection{Optymalizacja trajektorii}
\label{subsec:optymalizacja_grafu_implementacja}

Trajektoria składa się z póz robota w chwilach czasowych. Można je reprezentować jako graf. W celu ograniczenia rozmiaru grafu trajektorii, jego węzłami są pozy klatek kluczowych. Pomiędzy kolejnymi klatkami kluczowymi dodawane są krawędzie zawierające względną transformację wyznaczoną na podstawie bieżącej trajektorii. Krawędź pętli jest dodawana dopiero po wykryciu wcześniej odwiedzonego miejsca i~pozytywnym wyniku jego weryfikacji geometrycznej.

Błąd każdej krawędzi jest wyznaczany zgodnie z~równaniem~\eqref{eq:pose_graph_error}. W~implementacji nie jest przechowywana osobna pełna macierz informacji $\boldsymbol{\Omega}_{ij}$ dla każdej krawędzi. Zamiast tego składowe błędu są skalowane za pomocą wspólnych wag translacji i~rotacji:
\begin{equation}
    \mathbf{r}_{ij}
    =
    \mathbf{w}\odot\mathbf{e}_{ij},
    \qquad
    \mathbf{w}
    =
    (w_t,w_t,w_t,w_R,w_R,w_R)^T,
    \label{eq:implementation_pose_graph_residual}
\end{equation}
gdzie $w_t$ i~$w_R$ określają odpowiednio wagi składowych translacyjnych i~rotacyjnych, a~$\odot$ oznacza mnożenie odpowiadających sobie składowych.

Optymalizacja jest realizowana za pomocą funkcji \texttt{scipy.optimize.least\_squares}. Do każdej składowej wektora $\mathbf{r}_{ij}$ stosowana jest funkcja Hubera~\cite{huber1964robust}, ograniczająca wpływ dużych reszt. Jej próg określa parametr \texttt{huber\_scale}, którego wartość domyślna wynosi $1$. Pierwsza poza grafu pozostaje stała.

Po zaakceptowaniu pętli jej krawędź jest dodawana do grafu, a optymalizacja jest uruchamiana. Po jej poprawnym zakończeniu aktualizowane są pozy klatek kluczowych i~położenia punktów mapy, a~zapisane pozy klatek pośrednich są ponownie wyznaczane względem skorygowanych klatek kluczowych.

Zmiennymi optymalizacji są wyłącznie pozy klatek kluczowych. Położenia punktów 3D i~ich obserwacje 2D nie są optymalizowane wspólnie z~pozami, dlatego zastosowany mechanizm nie jest optymalizacją wiązki.

\section{Wyznaczanie wyników końcowych}
\label{sec:implementacja_wyjscia}

Podstawowym wynikiem działania systemu jest trajektoria lewej kamery zapisana jako pozy $\mathbf{T}_{W C_k}$. W~przypadku EvSLAM wymagane jest również wyznaczenie prędkości liniowej kamery, natomiast format wynikowy M3ED zawiera wyłącznie pozycję i~orientację.
Prędkość dla EvSLAM jest obliczana na podstawie wyznaczonej trajektorii. W~otoczeniu każdej chwili $t_i$ do kolejnych współrzędnych pozycji kamery w~układzie świata dopasowywany jest lokalny wielomian, a~jego pochodna w~punkcie $t_i$ stanowi oszacowanie prędkości~\cite{savitzky1964smoothing}. W~konfiguracji końcowej stosowany jest wielomian stopnia 2 i~okno 41 próbek.

Dla wielomianu drugiego stopnia zależność ta przyjmuje postać:
\begin{equation}
    \mathbf{p}_W(t_i+\Delta t)
    \approx
    \mathbf{a}_0
    +
    \mathbf{a}_1\Delta t
    +
    \mathbf{a}_2\Delta t^2,
    \qquad
    \mathbf{v}_W(t_i)=\mathbf{a}_1,
    \qquad
    \mathbf{v}_{C_k}
    =
    \mathbf{R}_{W C_k}^{T}\mathbf{v}_W,
    \label{eq:implementation_velocity}
\end{equation}
gdzie $\mathbf{p}_W$ jest pozycją kamery, $\Delta t$ przesunięciem czasu względem $t_i$, a~$\mathbf{a}_0$, $\mathbf{a}_1$ i~$\mathbf{a}_2$ współczynnikami dopasowanego wielomianu. Ostatnia zależność przekształca wyznaczoną prędkość z~układu świata do układu bieżącej kamery.

Przed utworzeniem pliku wynikowego trajektoria jest interpolowana do znaczników czasu wymaganych przez dany konkurs. Następnie prędkość jest wyznaczana na podstawie tak otrzymanej trajektorii i~zapisywana w~układzie kamery.
Zastosowanie takiego lokalnego dopasowania wielomianowego pozwala ograniczyć wzmacnianie niewielkich zmian i szumu trajektorii podczas numerycznego różniczkowania, bez rozszerzania estymatora SLAM o dodatkowy stan prędkości.

    \chapter{Metodyka badań}
\label{cha:metodyka_badan}

Celem badań była ocena dokładności opracowanego systemu oraz sprawdzenie wpływu jego najważniejszych parametrów na wynik. Ich zadaniem jest weryfikacja:
\begin{itemize}
    \item jaki efekt daje wykorzystanie mapy, relokalizacji i~domykania pętli względem odometrii wizyjnej, opartej wyłącznie na estymacji zmiany położenia między kolejnymi klatkami zdarzeniowymi. Pozwala to przedstawić różnicę między VO a~SLAM i~ocenić korzyść z~zastosowania mechanizmów SLAM;
    \item czy zastosowanie filtracji BAF poprawia wynik dla sekwencji badanych zbiorów i~jak zwiększa koszt obliczeniowy algorytmu;
    \item jak zmienia się estymacja po zastąpieniu rotacji wyznaczanej metodą PnP rotacją otrzymaną bezpośrednio z~żyroskopu;
    \item jaki wpływ na wynik ma czas akumulacji zdarzeń;
    \item jaką jakość osiąga konfiguracja końcowa na sekwencjach EvSLAM i~M3ED.
\end{itemize}
Każde porównanie kontrolowane dotyczy jednego zdefiniowanego wariantu systemu, a~pozostałe ustawienia pozostają niezmienione.

\section{Zbiory danych i zakres badań}
\label{sec:zbiory_danych}

W~badaniach wykorzystano zbiory EvSLAM i~M3ED opisane w~rozdziale~\ref{cha:projekt_systemu}.
Sekwencje EvSLAM obejmują ruch drona, platformy liniowej, ramienia robota i~robota kołowego.
Dla M3ED wybrano jedną sekwencję testową przeznaczoną do ewaluacji konkursu \cite{m3ed2025challenge} -- \path{spot_outdoor_day_penno_building_loop} oraz dwie sekwencje ze zbioru \cite{m3ed2025data}, nie brane pod uwagę przy ocenie wyników konkursu, lecz zawierające dane referencyjne (ang. \textit{ground truth}), umożliwiające ocenę jakości. Pierwszą z nich jest \path{falcon_outdoor_day_fast_flight_2} -- specjalnie wybrana jako trudny test możliwości algorytmu, mogącego wykraczać poza jego możliwości. Jest to szybki lot drona nad terenem zawierającym niewiele punktów odniesienia, znajdujących się daleko od kamer, co utrudnia poprawne oszacowanie głębi. Dodatkowym wyzwaniem jest to, że dron w~tej sekwencji wraca w~to samo miejsce tylko raz -- w~momencie lądowania. Kolejna sekwencja jest potencjalnie łatwiejsza -- jest to lot drona w~zamkniętym pomieszczeniu, pełnym charakterystycznych obiektów -- \path{falcon_indoor_flight_1}. W~tabeli użyto krótszych nazw odpowiadających sekwencjom.

Zestawienie użytych danych przedstawiono w~tabeli~\ref{tab:badane_sekwencje}.
\begin{table}[!ht]
    \centering
    \small
    \caption{Sekwencje wykorzystane w~badaniach.}
    \label{tab:badane_sekwencje}
    \begin{tabularx}{\textwidth}{
        >{\raggedright\arraybackslash}p{0.09\textwidth}
        >{\raggedright\arraybackslash}p{0.20\textwidth}
        >{\raggedright\arraybackslash}p{0.11\textwidth}
        >{\raggedright\arraybackslash}X
        >{\raggedright\arraybackslash}p{0.11\textwidth}
    }
        \toprule
        Zbiór & Sekwencja & Czas trwania & Platforma i~scenariusz & Referencja \\
        \midrule
        EvSLAM & \texttt{seq001} & 41,0~s & dron, spokojny ruch po owalnym torze & dostępna \\
        EvSLAM & \texttt{seq002} & 38,9~s & dron, szybszy ruch po owalnym torze & dostępna \\
        EvSLAM & \texttt{seq003} & 41,2~s & dron, spokojny ruch z~dominującymi odcinkami prostymi & dostępna \\
        EvSLAM & \texttt{seq004} & 33,4~s & dron, szybszy ruch z~dominującymi odcinkami prostymi & dostępna \\
        EvSLAM & \texttt{seq005} & 12,8~s & platforma liniowa, ruch po prowadnicy & dostępna \\
        EvSLAM & \texttt{seq006} & 26,8~s & ramię robota, krótkie i~gwałtowne ruchy kamery & dostępna \\
        EvSLAM & \texttt{seq007} & 35,9~s & robot kołowy, jazda po okręgu & dostępna \\
        EvSLAM & \texttt{seq007\_test} & 32,5~s & robot kołowy, jazda po okręgu, sekwencja niebrana pod uwagę w~wyniku końcowym, przeznaczona do~testów algorytmu & dostępna \\
        \midrule
        M3ED & \texttt{fast\_flight\_2} & 163~s & dron, szybki lot na zewnątrz w~dzień & dostępna \\
        M3ED & \texttt{indoor\_flight\_1} & 56,9~s & dron, lot między przeszkodami wewnątrz budynku & dostępna \\
        M3ED & \texttt{spot\_loop} & 295,9~s & robot kroczący, zewnętrzna trasa z~powrotem w~okolice początku & niedostępna \\
        \bottomrule
    \end{tabularx}
\end{table}

\section{Konfiguracja eksperymentów}
\label{sec:konfiguracje_badan}

Dla sekwencji zbioru EvSLAM stosuje się okna agregacji zdarzeń o~długości $t=12\,\mathrm{ms}$, natomiast dla M3ED $t=24\,\mathrm{ms}$. W~obu zbiorach wykorzystuje się wykładniczo zanikające klatki zdarzeniowe ze stałą zaniku $\tau=t$, kompensację obrotu do środka okna, detektor cech GFTT i~rotację z~żyroskopu. Stosowane są również te same kryteria PnP, wyboru klatek kluczowych oraz optymalizacji grafu. Geometria stereo jest przetwarzana w~ten sam sposób, z~kalibracją właściwą dla danego zestawu kamer. BAF jest włączony dla EvSLAM, a~wyłączony dla M3ED. Maksymalna głębokość stereo wynosi odpowiednio $10\,\mathrm{m}$ i~$100\,\mathrm{m}$.

W~rozpoznawaniu pętli M3ED zastosowano łagodniejsze progi odpowiedzialne za wykrywanie ponownie odwiedzonego miejsca.
Różnice w~parametrach wynikają z~odmiennych charakterystyk zbiorów. M3ED zawiera znacznie dłuższe sekwencje, często rejestrowane na zewnątrz, w~których obiekty znajdują się dalej od kamer, a~robot przebywa większe dystanse i~rzadziej wraca w~to samo miejsce. Z~tego powodu zwiększono czułość wykrywania pętli dla M3ED w~porównaniu z~EvSLAM.
Parametry zostały dobrane eksperymentalnie na podstawie wyników opisanych w~sekcji~\ref{sec:wyniki_eksperymenty} oraz testów rozwojowych.

\section{Ocena jakości estymowanej trajektorii}
\label{sec:protokol_oceny}

\subsection{Wyrównanie trajektorii}
\label{subsec:wyrownanie}

Próbki estymowane są wyrównywane czasowo z~referencyjnymi. W~oficjalnym protokole EvSLAM pozycje są wyrównywane metodą Umeyamy bez estymacji skali, czyli za pomocą transformacji $SE(3)$. Wyznaczane są globalny obrót i~translacja, natomiast przyjmuje się skalę równą $1$. Tak wyrównane pozycje służą do obliczenia ATE. Wartość AUC jest wyznaczana bezpośrednio z~dopasowanych w~czasie prędkości zapisanych w~układzie kamery~\cite{evslam2026evaluation,zhong2026evslam}.

M3ED w~ewaluacji wyników korzysta z~narzędzia \textit{evo}~\cite{evo2025alignment}. Stosowane jest wyrównanie metodą Umeyamy (podrozdział ~\ref{subsec:wyrownanie_umeyamy}) ze skalą, czyli transformacja $Sim(3)$, w~której wyznaczane są obrót, translacja i~współczynnik skali.

Dla porównania wyników pomiędzy zbiorami do obliczenia metryk dla obu zbiorów zastosowano transformację $SE(3)$.
Do wizualizacji trajektorie wyjściowe przesunięto jedynie do punktu początkowego trajektorii referencyjnej. Pozwala to ocenić rzeczywisty dryf oraz dokładność zachowania skali systemu.

\subsection{Metryki}
\label{subsec:metryki_trajektorii}

Niech $\mathbf{p}_i^{\mathrm{est}}$ i~$\mathbf{p}_i^{\mathrm{gt}}$ oznaczają odpowiednio dopasowane w~czasie pozycję estymowaną i~referencyjną. Po transformacji $SE(3)$ błąd próbki wynosi
\begin{equation}
    e_i=\left\|\mathbf{p}_i^{\mathrm{gt}}-\bigl(\mathbf{R}\mathbf{p}_i^{\mathrm{est}}+\mathbf{t}\bigr)\right\|_2,
    \label{eq:methodology_position_error}
\end{equation}
gdzie $\mathbf{R}$ i~$\mathbf{t}$ są obrotem i~translacją wyznaczonymi metodą Umeyamy.

\textbf{ATE} (ang. \textit{Absolute Trajectory Error}) określa globalny błąd translacji:
\begin{equation}
    \mathrm{ATE_{RMSE}}=\sqrt{\frac{1}{N}\sum_{i=1}^{N}e_i^2}.
    \label{eq:methodology_ate}
\end{equation}
Analogicznie raportowany \textbf{RMSE} (ang. \textit{Root Mean Square Error}) rotacji jest pierwiastkiem ze średniego kwadratu kąta różnicy orientacji i~jest podawany w~stopniach.

\textbf{RPE} (ang. \textit{Relative Pose Error}) opisuje błąd ruchu względnego w~odstępie 1~s~\cite{sturm2012benchmark}. Dla pary próbek $i, j$ definiuje się
\begin{equation}
    \Delta\mathbf{T}_{ij}=\mathbf{T}_i^{-1}\mathbf{T}_j,
    \qquad
    \mathbf{E}_{ij}=\left(\Delta\mathbf{T}_{ij}^{\mathrm{gt}}\right)^{-1}\Delta\mathbf{T}_{ij}^{\mathrm{est}},
    \label{eq:methodology_rpe}
\end{equation}
gdzie $j$ jest próbką położoną najbliżej chwili $t_i+1$~s, a~$\mathbf{T}_i$ oznacza pozę w~chwili $t_i$ (estymowaną lub referencyjną). Osobno podaje się normę translacyjnej części $\mathbf{E}_{ij}$ w~metrach oraz kąt jej części obrotowej w~stopniach.

Miara \textbf{AUC} jest polem pod ważoną krzywą sukcesu prędkości~\cite{evslam2026evaluation}. Dla znormalizowanego błędu
\begin{equation}
    \varepsilon_i=\frac{\left\|\mathbf{v}_i^{\mathrm{gt}}-\mathbf{v}_i^{\mathrm{est}}\right\|_2}{\left\|\mathbf{v}_i^{\mathrm{gt}}\right\|_2}
    \label{eq:methodology_velocity_error}
\end{equation}
próbce przypisuje się wagę $w_i=\|\mathbf{v}_i^{\mathrm{gt}}\|_2/\sum_j\|\mathbf{v}_j^{\mathrm{gt}}\|_2$. Większą wagę otrzymują zatem próbki odpowiadające szybszemu ruchowi. Krzywa sukcesu ma postać
\begin{equation}
    S(\xi)=\sum_i w_i\,\mathbb{I}(\varepsilon_i<\xi),
    \qquad
    \mathrm{AUC}=\int_0^1 S(\xi)\,\mathrm{d}\xi,
    \label{eq:methodology_auc}
\end{equation}
gdzie $\mathbb{I}$ jest funkcją wskaźnikową. Próbki o~zerowej prędkości referencyjnej otrzymują zerową wagę; w~implementacji ich błąd względny jest przyjmowany jako równy zeru, co zapobiega dzieleniu przez zero. Pole jest przybliżane dla 500 równomiernie rozmieszczonych progów $\xi$ od 0 do 1.

Zestawienie raportowanych wielkości zawiera tabela~\ref{tab:metryki_trajektorii}.

\begin{table}[!ht]
    \centering
    \small
    \caption{Metryki stosowane do oceny trajektorii.}
    \label{tab:metryki_trajektorii}
    \begin{tabularx}{\textwidth}{
        >{\raggedright\arraybackslash}p{0.22\textwidth}
        >{\raggedright\arraybackslash}X
        >{\raggedright\arraybackslash}p{0.16\textwidth}
    }
        \toprule
        Metryka & Znaczenie & Jednostka \\
        \midrule
        ATE RMSE & globalny błąd położenia & m \\
        RMSE rotacji & globalna różnica orientacji & stopnie \\
        RPE translacji & lokalny błąd ruchu & m \\
        RPE rotacji & lokalny błąd obrotu & stopnie \\
        AUC & pole pod ważoną krzywą sukcesu prędkości & brak \\
        \bottomrule
    \end{tabularx}
\end{table}

\section{Kontrolowane porównania}
\label{sec:badania_wariantow}

Zakres eksperymentów zestawiono w~tabeli~\ref{tab:macierz_eksperymentow}. Porównywane warianty muszą obejmować ten sam fragment sekwencji i~korzystać z~tego samego protokołu synchronizacji oraz metryk.

\begin{table}[!ht]
    \centering
    \footnotesize
    \caption{Macierz kontrolowanych porównań.}
    \label{tab:macierz_eksperymentow}
    \begin{tabularx}{\textwidth}{p{0.22\textwidth} p{0.34\textwidth} X}
        \toprule
        Badany element & Sekwencja & Warianty \\
        \midrule
        mechanizmy SLAM & EvSLAM \texttt{seq007\_test} & pełny SLAM / VO \\
        filtr BAF & EvSLAM \texttt{seq001}, 961,7--976,7~s; M3ED \texttt{indoor\_flight\_1}, 5--20~s & włączony / wyłączony \\
        źródło rotacji & M3ED \texttt{indoor\_flight\_1}, całość & PnP / żyroskop \\
        czas akumulacji & EvSLAM \texttt{seq006}, 5--20~s; M3ED \texttt{indoor\_flight\_1}, 5--20~s & 6 / 12 / 24 / 48~ms \\
        \bottomrule
    \end{tabularx}
\end{table}

    \chapter{Wyniki}
\label{cha:wyniki}

W~rozdziale przedstawiono wyniki eksperymentów porównawczych opisanych w~rozdziale~\ref{cha:metodyka_badan}, które pomagają dobrać najważniejsze parametry algorytmu i~ocenić ich wpływ na jakość wyniku. Przedstawiono również pełne wyniki dla wszystkich sekwencji ze zbioru EvSLAM oraz wybranych sekwencji ze zbioru M3ED w~celu końcowej ewaluacji.

\section{Eksperymenty porównawcze}
\label{sec:wyniki_eksperymenty}

\subsection{SLAM i odometria wizyjna}
\label{subsec:wyniki_slam_vo}

Uruchomiono sekwencję \texttt{seq007\_test} z~wyłączonymi wybranymi mechanizmami algorytmu. Były to:
\begin{itemize}
    \item budowanie rzadkiej mapy,
    \item estymacja pozy względem mapy,
    \item zamknięcie pętli,
    \item optymalizacja grafu trajektorii.
\end{itemize}

Wariant odometrii wizyjnej uzyskano przez wyłączenie mechanizmów związanych z~utrzymywaniem i~wykorzystywaniem mapy. Ruch kamery był w~nim wyznaczany wyłącznie na podstawie punktów 3D otrzymanych z~poprzedniej pary stereo. Wyłączono:
\begin{itemize}
    \item tworzenie klatek kluczowych i~rzadkiej mapy punktów 3D,
    \item estymację pozy na podstawie punktów mapy lokalnej,
    \item rozpoznawanie miejsc i~zamknięcie pętli,
    \item optymalizację grafu trajektorii.
\end{itemize}

Wyniki przedstawiono w~tabeli~\ref{tab:wyniki_slam_vo}.

\begin{table}[!ht]
    \centering
    \small
    \caption{Porównanie SLAM i~VO na wspólnym fragmencie \texttt{seq007\_test}; wyrównanie $SE(3)$, RPE dla około 1~s.}
    \label{tab:wyniki_slam_vo}
    \resizebox{\textwidth}{!}{%
        \begin{tabular}{lrrrrrr}
            \toprule
            Wariant & ATE RMSE [m] & RMSE rotacji [$^\circ$] & AUC & RPE translacji [m] & RPE rotacji [$^\circ$] & Próbki \\
            \midrule
            SLAM & 0,151 & 3,179 & 0,756 & 0,187 & 4,227 & 275 \\
            VO & 0,544 & 16,256 & 0,772 & 0,174 & 5,712 & 275 \\
            \bottomrule
        \end{tabular}%
    }
\end{table}

\begin{figure}[!ht]
    \centering
    \begin{subfigure}[t]{0.49\textwidth}
        \centering
        \includegraphics[width=0.49\linewidth]{ablation_evslam_seq007_slam_first_pose_trajectory_x.png}
        \includegraphics[width=0.49\linewidth]{ablation_evslam_seq007_slam_first_pose_trajectory_y.png}
        \includegraphics[width=0.49\linewidth]{ablation_evslam_seq007_slam_first_pose_trajectory_z.png}
        \includegraphics[width=0.49\linewidth]{ablation_evslam_seq007_slam_first_pose_trajectory_3d.png}
        \caption{SLAM.}
    \end{subfigure}
    \hfill
    \begin{subfigure}[t]{0.49\textwidth}
        \centering
        \includegraphics[width=0.49\linewidth]{ablation_evslam_seq007_vo_first_pose_trajectory_x.png}
        \includegraphics[width=0.49\linewidth]{ablation_evslam_seq007_vo_first_pose_trajectory_y.png}
        \includegraphics[width=0.49\linewidth]{ablation_evslam_seq007_vo_first_pose_trajectory_z.png}
        \includegraphics[width=0.49\linewidth]{ablation_evslam_seq007_vo_first_pose_trajectory_3d.png}
        \caption{Odometria wizyjna.}
    \end{subfigure}
    \caption{Trajektorie w~historycznym porównaniu SLAM i~VO. Kolejne małe wykresy przedstawiają $x$, $y$, $z$ i~widok 3D.}
    \label{fig:wyniki_slam_vo}
\end{figure}

Warstwa SLAM wyraźnie zmniejszyła globalny błąd położenia i~orientacji oraz rotacyjny RPE. Nie poprawiła jednak AUC ani translacyjnego RPE. Na wykresach trajektorii przedstawionych na rysunku~\ref{fig:wyniki_slam_vo} można zaobserwować spodziewaną różnicę między SLAM a~VO. W~systemie bez korekty pozy na podstawie wykrycia wcześniej odwiedzonego miejsca występuje znaczny dryf estymowanego położenia.

\subsection{Filtr BAF}
\label{subsec:wyniki_baf}

W celu weryfikacji efektów zastosowania filtracji uruchomiono system z włączonym filtrem BAF oraz bez niego. Wybrano sekwencje EvSLAM \texttt{seq001} oraz M3ED \texttt{falcon\_indoor\_flight\_1}. Dla obu sekwencji przetworzono fragmenty z jednakowego przedziału czasowego. Aby zmniejszyć czas trwania eksperymentu zastosowano krótki 15-sekundowy przedział (od 5 do 20 s), który jednak pozwala na ocenę efektów filtracji. Rozpoczęcie eksperymentu od 5 sekundy pozwala na pominięcie początkowego etapu, w którym ruch nie występuje lub jest powolny.

\begin{table}[!ht]
    \centering
    \small
    \caption{Porównanie z~BAF i~bez BAF na fragmentach EvSLAM \texttt{seq001} i~M3ED \texttt{falcon\_indoor\_flight\_1};}
    \label{tab:wyniki_baf}
    \resizebox{\textwidth}{!}{%
        \begin{tabular}{llrrrrrr}
            \toprule
            Zbiór & BAF & Zmiana czasu przetwarzania po zastosowaniu filtracji & ATE RMSE $SE(3)$ [m] & RMSE rotacji [$^\circ$] & AUC & RPE translacji [m] & RPE rotacji [$^\circ$] \\
            \midrule
            EvSLAM & wyłączony & -- & 0,675 & 5,197 & 0,852 & 0,411 & 2,540 \\
            EvSLAM & włączony & +7,2\% & 0,574 & 6,784 & 0,849 & 0,396 & 2,544 \\
            M3ED & wyłączony & -- & 0,103 & 1,540 & 0,866 & 0,068 & 0,482 \\
            M3ED & włączony & +207,5\% & 0,065 & 2,332 & 0,903 & 0,050 & 0,435 \\
            \bottomrule
        \end{tabular}%
    }
\end{table}

Na fragmencie EvSLAM filtr zmniejszył ATE i~translacyjny RPE, lecz nieznacznie pogorszył pozostałe miary. Czas wykonania wzrósł przy tym z~781~s do 837~s, czyli o~7,2\%. Na fragmencie M3ED filtr poprawił wszystkie raportowane błędy poza globalnym RMSE rotacji, ale wydłużył czas przetwarzania z~255~s do 784~s, czyli o~207,5\%. Podane wyniki pozwalają ocenić koszt BAF wyłącznie dla badanych sekwencji i~konfiguracji. Ze względu na większą liczbę zdarzeń występujących w~sekwencji M3ED filtr znacznie wydłużył obliczenia. Ponieważ cecha ta występuje również w~innych sekwencjach tego zbioru, zdecydowano o~wyłączeniu filtracji dla M3ED i~pozostawieniu jej dla EvSLAM.
Odpowiadające tym wynikom przebiegi trajektorii przedstawiono na rysunku~\ref{fig:wyniki_baf}.

\begin{figure}[!ht]
    \centering
    \begin{subfigure}[t]{0.49\textwidth}
        \centering
        \includegraphics[width=0.49\linewidth]{ablation_evslam_seq001_baf_disabled_first_pose_trajectory_x.png}
        \includegraphics[width=0.49\linewidth]{ablation_evslam_seq001_baf_disabled_first_pose_trajectory_y.png}
        \includegraphics[width=0.49\linewidth]{ablation_evslam_seq001_baf_disabled_first_pose_trajectory_z.png}
        \includegraphics[width=0.49\linewidth]{ablation_evslam_seq001_baf_disabled_first_pose_trajectory_3d.png}
        \caption{EvSLAM, BAF wyłączony.}
    \end{subfigure}
    \hfill
    \begin{subfigure}[t]{0.49\textwidth}
        \centering
        \includegraphics[width=0.49\linewidth]{ablation_evslam_seq001_baf_enabled_first_pose_trajectory_x.png}
        \includegraphics[width=0.49\linewidth]{ablation_evslam_seq001_baf_enabled_first_pose_trajectory_y.png}
        \includegraphics[width=0.49\linewidth]{ablation_evslam_seq001_baf_enabled_first_pose_trajectory_z.png}
        \includegraphics[width=0.49\linewidth]{ablation_evslam_seq001_baf_enabled_first_pose_trajectory_3d.png}
        \caption{EvSLAM, BAF włączony.}
    \end{subfigure}

    \begin{subfigure}[t]{0.49\textwidth}
        \centering
        \includegraphics[width=0.49\linewidth]{ablation_m3ed_falcon_indoor_flight_1_baf_disabled_first_pose_trajectory_x.png}
        \includegraphics[width=0.49\linewidth]{ablation_m3ed_falcon_indoor_flight_1_baf_disabled_first_pose_trajectory_y.png}
        \includegraphics[width=0.49\linewidth]{ablation_m3ed_falcon_indoor_flight_1_baf_disabled_first_pose_trajectory_z.png}
        \includegraphics[width=0.49\linewidth]{ablation_m3ed_falcon_indoor_flight_1_baf_disabled_first_pose_trajectory_3d.png}
        \caption{M3ED, BAF wyłączony.}
    \end{subfigure}
    \hfill
    \begin{subfigure}[t]{0.49\textwidth}
        \centering
        \includegraphics[width=0.49\linewidth]{ablation_m3ed_falcon_indoor_flight_1_baf_enabled_first_pose_trajectory_x.png}
        \includegraphics[width=0.49\linewidth]{ablation_m3ed_falcon_indoor_flight_1_baf_enabled_first_pose_trajectory_y.png}
        \includegraphics[width=0.49\linewidth]{ablation_m3ed_falcon_indoor_flight_1_baf_enabled_first_pose_trajectory_z.png}
        \includegraphics[width=0.49\linewidth]{ablation_m3ed_falcon_indoor_flight_1_baf_enabled_first_pose_trajectory_3d.png}
        \caption{M3ED, BAF włączony.}
    \end{subfigure}
    \caption{Trajektorie 15-sekundowych fragmentów w~porównaniu BAF. Kolejne małe wykresy przedstawiają $x$, $y$, $z$ i~widok 3D.}
    \label{fig:wyniki_baf}
\end{figure}

\clearpage

\subsection{Źródło rotacji}
\label{subsec:wyniki_rotation_source}

System uruchomiono w~dwóch wariantach na pełnej sekwencji M3ED \texttt{falcon\_indoor\_flight\_1}. W~pierwszym PnP wyznacza rotację i~translację, a~w~drugim rotacja pochodzi z~żyroskopu i~z~wizji wyznaczana jest tylko translacja. Wyniki zestawiono w~tabeli~\ref{tab:wyniki_rotation_source}.

\begin{table}[!ht]
    \centering
    \small
    \caption{Porównanie źródła rotacji na pełnej sekwencji M3ED \texttt{falcon\_indoor\_flight\_1}; wyrównanie $SE(3)$, RPE dla około 1~s.}
    \label{tab:wyniki_rotation_source}
    \resizebox{\textwidth}{!}{%
        \begin{tabular}{lrrrrr}
            \toprule
            Wariant & ATE RMSE [m] & RMSE rot. [$^\circ$] & AUC & RPE trans. [m] & RPE rot. [$^\circ$] \\
            \midrule
            pełne PnP & 0,240 & 9,885 & 0,748 & 0,152 & 1,861 \\
            rotacja z~żyroskopu & 0,162 & 4,434 & 0,830 & 0,092 & 0,514 \\
            \bottomrule
        \end{tabular}%
    }
\end{table}

\begin{figure}[!ht]
    \centering
    \begin{subfigure}[t]{0.96\textwidth}
        \centering
        \includegraphics[width=0.24\linewidth]{ablation_m3ed_falcon_indoor_flight_1_rotation_from_pnp_first_pose_trajectory_x.png}
        \includegraphics[width=0.24\linewidth]{ablation_m3ed_falcon_indoor_flight_1_rotation_from_pnp_first_pose_trajectory_y.png}
        \includegraphics[width=0.24\linewidth]{ablation_m3ed_falcon_indoor_flight_1_rotation_from_pnp_first_pose_trajectory_z.png}
        \includegraphics[width=0.24\linewidth]{ablation_m3ed_falcon_indoor_flight_1_rotation_from_pnp_first_pose_trajectory_3d.png}
        \caption{Rotacja i~translacja z~PnP.}
    \end{subfigure}

    \vspace{2mm}

    \begin{subfigure}[t]{0.96\textwidth}
        \centering
        \includegraphics[width=0.24\linewidth]{ablation_m3ed_falcon_indoor_flight_1_rotation_from_imu_first_pose_trajectory_x.png}
        \includegraphics[width=0.24\linewidth]{ablation_m3ed_falcon_indoor_flight_1_rotation_from_imu_first_pose_trajectory_y.png}
        \includegraphics[width=0.24\linewidth]{ablation_m3ed_falcon_indoor_flight_1_rotation_from_imu_first_pose_trajectory_z.png}
        \includegraphics[width=0.24\linewidth]{ablation_m3ed_falcon_indoor_flight_1_rotation_from_imu_first_pose_trajectory_3d.png}
        \caption{Rotacja z~żyroskopu, translacja z~wizji.}
    \end{subfigure}
    \caption{Trajektorie pełnej sekwencji M3ED \texttt{falcon\_indoor\_flight\_1} w~porównaniu źródła rotacji. W~każdym wierszu pokazano kolejno $x$, $y$, $z$ i~widok 3D.}
    \label{fig:wyniki_rotation_source}
\end{figure}

Rotacja z~żyroskopu zmniejszyła ATE, RMSE rotacji oraz obie składowe RPE, a~także zwiększyła AUC. Dla tej sekwencji jej bezpośrednie użycie było więc korzystniejsze od wyznaczania pełnej pozy za pomocą PnP. Trajektorie przedstawiono na rysunku~\ref{fig:wyniki_rotation_source}. Ponieważ w~czasie rozwoju systemu podobne efekty zauważono także dla innych testowanych sekwencji, w~konfiguracji końcowej przyjęto rotację wyznaczaną za pomocą żyroskopu.

\subsection{Czas akumulacji zdarzeń}
\label{subsec:wyniki_accumulation}

Rysunek~\ref{fig:wyniki_window_previews} pokazuje bezpośredni wpływ długości okna na reprezentację tego samego momentu sekwencji. Klatki utworzoen przy użyciu dłuższego okna zawierają więcej aktywności.

Wpływ na estymację oceniono na fragmentach z jednakowego przedziału 5--20~s sekwencji EvSLAM \texttt{seq006} i~M3ED \texttt{falcon\_indoor\_flight\_1}. W~każdym wariancie stała zaniku $\tau$ była równa długości okna. Przedział czasowy ograniczony do $15 s$ pozwala skutecznie ocenić wpływ rozmiaru okna, pozwalając na przeprowadzenie eksperymentów dla większego zbioru badanych wartości, dzięki ograniczeniu czasu przetwarzania sekwencji. Rozpoczęcie eksperymentu od 5 sekundy pozwala na pominięcie początkowego etapu, w którym ruch nie występuje lub jest powolny.

\begin{figure}[!ht]
    \centering
    \begin{subfigure}[t]{0.235\textwidth}
        \centering
        \includegraphics[width=\textwidth]{ablation_evslam_seq006_accumulation_006ms_left_event_frame.png}
        \caption{6~ms.}
    \end{subfigure}
    \hfill
    \begin{subfigure}[t]{0.235\textwidth}
        \centering
        \includegraphics[width=\textwidth]{ablation_evslam_seq006_accumulation_012ms_left_event_frame.png}
        \caption{12~ms.}
    \end{subfigure}
    \hfill
    \begin{subfigure}[t]{0.235\textwidth}
        \centering
        \includegraphics[width=\textwidth]{ablation_evslam_seq006_accumulation_024ms_left_event_frame.png}
        \caption{24~ms.}
    \end{subfigure}
    \hfill
    \begin{subfigure}[t]{0.235\textwidth}
        \centering
        \includegraphics[width=\textwidth]{ablation_evslam_seq006_accumulation_048ms_left_event_frame.png}
        \caption{48~ms.}
    \end{subfigure}
    \caption{Reprezentacje zdarzeń z~lewej kamery utworzone dla tego samego fragmentu \texttt{seq006} przy czterech czasach akumulacji.}
    \label{fig:wyniki_window_previews}
\end{figure}

\begin{table}[!ht]
    \centering
    \small
    \caption{Porównanie czasu akumulacji; metryki po wyrównaniu $SE(3)$.}
    \label{tab:wyniki_accumulation}
    \resizebox{\textwidth}{!}{%
        \begin{tabular}{llrrrrr}
            \toprule
            Zbiór & Okno [ms] & ATE RMSE [m] & RMSE rot. [$^\circ$] & AUC & RPE trans. [m] & RPE rot. [$^\circ$] \\
            \midrule
            EvSLAM & 6 & 0,102 & 11,900 & 0,498 & 0,186 & 19,800 \\
            EvSLAM & 12 & 0,111 & 12,740 & 0,502 & 0,196 & 20,704 \\
            EvSLAM & 24 & 0,161 & 21,058 & 0,398 & 0,228 & 24,788 \\
            EvSLAM & 48 & 0,296 & 47,304 & 0,093 & 0,452 & 63,585 \\
            M3ED & 6 & 0,431 & 8,208 & 0,515 & 0,295 & 0,671 \\
            M3ED & 12 & 0,224 & 5,102 & 0,717 & 0,131 & 0,367 \\
            M3ED & 24 & 0,103 & 1,540 & 0,866 & 0,068 & 0,481 \\
            M3ED & 48 & 0,055 & 2,856 & 0,927 & 0,049 & 0,493 \\
            \bottomrule
        \end{tabular}%
    }
\end{table}

\clearpage

Rysunki~\ref{fig:wyniki_accumulation_evslam} i~\ref{fig:wyniki_accumulation_m3ed} przedstawiają odpowiednio przebiegi EvSLAM i~M3ED.

\begin{figure}[!ht]
    \centering
    \begin{subfigure}[t]{0.96\textwidth}
        \centering
        \includegraphics[width=0.24\linewidth]{ablation_evslam_seq006_accumulation_006ms_first_pose_trajectory_x.png}
        \includegraphics[width=0.24\linewidth]{ablation_evslam_seq006_accumulation_006ms_first_pose_trajectory_y.png}
        \includegraphics[width=0.24\linewidth]{ablation_evslam_seq006_accumulation_006ms_first_pose_trajectory_z.png}
        \includegraphics[width=0.24\linewidth]{ablation_evslam_seq006_accumulation_006ms_first_pose_trajectory_3d.png}
        \caption{Okno 6~ms.}
    \end{subfigure}

    \begin{subfigure}[t]{0.96\textwidth}
        \centering
        \includegraphics[width=0.24\linewidth]{ablation_evslam_seq006_accumulation_012ms_first_pose_trajectory_x.png}
        \includegraphics[width=0.24\linewidth]{ablation_evslam_seq006_accumulation_012ms_first_pose_trajectory_y.png}
        \includegraphics[width=0.24\linewidth]{ablation_evslam_seq006_accumulation_012ms_first_pose_trajectory_z.png}
        \includegraphics[width=0.24\linewidth]{ablation_evslam_seq006_accumulation_012ms_first_pose_trajectory_3d.png}
        \caption{Okno 12~ms.}
    \end{subfigure}

    \begin{subfigure}[t]{0.96\textwidth}
        \centering
        \includegraphics[width=0.24\linewidth]{ablation_evslam_seq006_accumulation_024ms_first_pose_trajectory_x.png}
        \includegraphics[width=0.24\linewidth]{ablation_evslam_seq006_accumulation_024ms_first_pose_trajectory_y.png}
        \includegraphics[width=0.24\linewidth]{ablation_evslam_seq006_accumulation_024ms_first_pose_trajectory_z.png}
        \includegraphics[width=0.24\linewidth]{ablation_evslam_seq006_accumulation_024ms_first_pose_trajectory_3d.png}
        \caption{Okno 24~ms.}
    \end{subfigure}

    \begin{subfigure}[t]{0.96\textwidth}
        \centering
        \includegraphics[width=0.24\linewidth]{ablation_evslam_seq006_accumulation_048ms_first_pose_trajectory_x.png}
        \includegraphics[width=0.24\linewidth]{ablation_evslam_seq006_accumulation_048ms_first_pose_trajectory_y.png}
        \includegraphics[width=0.24\linewidth]{ablation_evslam_seq006_accumulation_048ms_first_pose_trajectory_z.png}
        \includegraphics[width=0.24\linewidth]{ablation_evslam_seq006_accumulation_048ms_first_pose_trajectory_3d.png}
        \caption{Okno 48~ms.}
    \end{subfigure}
    \caption{Trajektorie fragmentu 5--20~s sekwencji EvSLAM \texttt{seq006} dla czterech czasów akumulacji. W~każdym wierszu pokazano kolejno $x$, $y$, $z$ i~widok 3D.}
    \label{fig:wyniki_accumulation_evslam}
\end{figure}

\begin{figure}[!ht]
    \centering
    \begin{subfigure}[t]{0.96\textwidth}
        \centering
        \includegraphics[width=0.24\linewidth]{ablation_m3ed_falcon_indoor_flight_1_accumulation_006ms_first_pose_trajectory_x.png}
        \includegraphics[width=0.24\linewidth]{ablation_m3ed_falcon_indoor_flight_1_accumulation_006ms_first_pose_trajectory_y.png}
        \includegraphics[width=0.24\linewidth]{ablation_m3ed_falcon_indoor_flight_1_accumulation_006ms_first_pose_trajectory_z.png}
        \includegraphics[width=0.24\linewidth]{ablation_m3ed_falcon_indoor_flight_1_accumulation_006ms_first_pose_trajectory_3d.png}
        \caption{Okno 6~ms.}
    \end{subfigure}

    \begin{subfigure}[t]{0.96\textwidth}
        \centering
        \includegraphics[width=0.24\linewidth]{ablation_m3ed_falcon_indoor_flight_1_accumulation_012ms_first_pose_trajectory_x.png}
        \includegraphics[width=0.24\linewidth]{ablation_m3ed_falcon_indoor_flight_1_accumulation_012ms_first_pose_trajectory_y.png}
        \includegraphics[width=0.24\linewidth]{ablation_m3ed_falcon_indoor_flight_1_accumulation_012ms_first_pose_trajectory_z.png}
        \includegraphics[width=0.24\linewidth]{ablation_m3ed_falcon_indoor_flight_1_accumulation_012ms_first_pose_trajectory_3d.png}
        \caption{Okno 12~ms.}
    \end{subfigure}

    \begin{subfigure}[t]{0.96\textwidth}
        \centering
        \includegraphics[width=0.24\linewidth]{ablation_m3ed_falcon_indoor_flight_1_accumulation_024ms_first_pose_trajectory_x.png}
        \includegraphics[width=0.24\linewidth]{ablation_m3ed_falcon_indoor_flight_1_accumulation_024ms_first_pose_trajectory_y.png}
        \includegraphics[width=0.24\linewidth]{ablation_m3ed_falcon_indoor_flight_1_accumulation_024ms_first_pose_trajectory_z.png}
        \includegraphics[width=0.24\linewidth]{ablation_m3ed_falcon_indoor_flight_1_accumulation_024ms_first_pose_trajectory_3d.png}
        \caption{Okno 24~ms.}
    \end{subfigure}

    \begin{subfigure}[t]{0.96\textwidth}
        \centering
        \includegraphics[width=0.24\linewidth]{ablation_m3ed_falcon_indoor_flight_1_accumulation_048ms_first_pose_trajectory_x.png}
        \includegraphics[width=0.24\linewidth]{ablation_m3ed_falcon_indoor_flight_1_accumulation_048ms_first_pose_trajectory_y.png}
        \includegraphics[width=0.24\linewidth]{ablation_m3ed_falcon_indoor_flight_1_accumulation_048ms_first_pose_trajectory_z.png}
        \includegraphics[width=0.24\linewidth]{ablation_m3ed_falcon_indoor_flight_1_accumulation_048ms_first_pose_trajectory_3d.png}
        \caption{Okno 48~ms.}
    \end{subfigure}
    \caption{Trajektorie fragmentu 5--20~s sekwencji M3ED \texttt{falcon\_indoor\_flight\_1} dla czterech czasów akumulacji. W~każdym wierszu pokazano kolejno $x$, $y$, $z$ i~widok 3D.}
    \label{fig:wyniki_accumulation_m3ed}
\end{figure}

\clearpage

Wyniki zestawione w~tabeli~\ref{tab:wyniki_accumulation} pokazują, że na fragmencie sekwencji ze zbioru EvSLAM najniższe ATE i~błędy rotacji uzyskano dla 6~ms, a~najwyższe AUC dla 12~ms. Dłuższe okna pogarszały wynik szybkiego ruchu ramienia. Na fragmencie sekwencji ze zbioru M3ED okno 48~ms dało najmniejsze ATE i~translacyjny RPE oraz najwyższe AUC, natomiast najmniejszy globalny błąd rotacji uzyskano dla 24~ms. Porównanie potwierdza, że korzystny czas akumulacji zależy od charakteru ruchu i~zbioru danych. W~konfiguracji końcowej przyjęto:
\begin{itemize}
    \item 12~ms dla EvSLAM, ze względu na najwyższą wartość najważniejszej w~konkursie metryki AUC;
    \item 24~ms dla M3ED jako wartość bezpieczniejszą dla często występujących w~tym zbiorze sekwencji bardziej dynamicznych niż testowana, dla której najlepszy wynik uzyskano przy 48~ms.
\end{itemize}

\section{Wyniki dla zbioru EvSLAM}
\label{sec:wyniki_evslam}

\subsection{Dokładność trajektorii i prędkości}
\label{subsec:wyniki_evslam_metryki}

Dla zbioru EvSLAM przetworzono wszystkie siedem dostępnych sekwencji testowych.

\begin{table}[!ht]
    \centering
    \small
    \caption{Wyniki końcowe na zbiorze EvSLAM. ATE RMSE i~AUC obliczono zgodnie z~oficjalnym protokołem~\cite{zhong2026evslam,evslam2026evaluation}.}
    \label{tab:wyniki_evslam_final}
    \begin{tabularx}{\textwidth}{l X r r r}
        \toprule
        Sekwencja & Scenariusz & Pokrycie GT [\%] & ATE RMSE [m] & AUC \\
        \midrule
        \texttt{seq001} & dron, ruch po okręgu & 99,8 & 0,901 & 0,821 \\
        \texttt{seq002} & dron, szybszy ruch po okręgu & 99,7 & 0,780 & 0,828 \\
        \texttt{seq003} & dron, ruch z~odcinkami prostymi & 99,8 & 0,333 & 0,852 \\
        \texttt{seq004} & dron, szybszy ruch z~odcinkami prostymi & 99,7 & 0,649 & 0,757 \\
        \texttt{seq005} & platforma liniowa & 98,4 & 0,077 & 0,800 \\
        \texttt{seq006} & ramię robota & 99,6 & 0,111 & 0,506 \\
        \texttt{seq007} & robot kołowy, jazda po okręgu & 100,0 & 0,801 & 0,433 \\
        \midrule
        Suma konkursowa & & & 3,654 & 4,996 \\
        Średnia & & 99,6 & 0,522 & 0,714 \\
        \bottomrule
    \end{tabularx}
\end{table}

Najmniejszy błąd położenia uzyskano dla ruchu platformy liniowej, a~największy dla \texttt{seq001}. Najwyższe AUC osiągnięto dla \texttt{seq003}, natomiast najniższe dla \texttt{seq007}. Różna kolejność sekwencji według ATE i~AUC pokazuje, że miary oceniają inne własności wyniku: zgodność globalnej trajektorii oraz chwilową estymację prędkości. W~rankingu konkursowym sumowane są wartości siedmiu sekwencji; średnie podano jedynie dla łatwiejszej interpretacji.

Na wszystkich rysunkach porównujących trajektorię estymowaną z~GT wyrównano ją wyłącznie do pierwszej dopasowanej pozy referencyjnej. Dalszy przebieg ani skala nie są korygowane, dlatego na wykresach pozostają widoczne narastające błędy orientacji i~skali, jeśli występują. Wyrównanie Umeyamy jest używane wyłącznie do obliczania metryk.

Przebiegi dla \texttt{seq001}--\texttt{seq007} przedstawiono odpowiednio na rysunkach~\ref{fig:evslam_seq001_trajectory}, \ref{fig:evslam_seq002_trajectory}, \ref{fig:evslam_seq003_trajectory}, \ref{fig:evslam_seq004_trajectory}, \ref{fig:evslam_seq005_trajectory}, \ref{fig:evslam_seq006_trajectory} i~\ref{fig:evslam_seq007_trajectory}. Każdy rysunek przedstawia wykresy estymowanych składowych położenia, widok trajektorii w~trójwymiarowej przestrzeni oraz wykresy estymowanych składowych prędkości liniowej. Na rysunkach przedstawiono również dane z trajektorii referencyjnych (ang. \textit{groundtruth}), co umożliwia wizualną ocenę jakości wyników.

\begin{figure}[!ht]
    \centering
    \includegraphics[width=0.40\textwidth]{results_evslam_seq001_trajectory_x.png}
    \includegraphics[width=0.40\textwidth]{results_evslam_seq001_trajectory_y.png}
    \includegraphics[width=0.40\textwidth]{results_evslam_seq001_trajectory_z.png}
    \includegraphics[width=0.40\textwidth]{results_evslam_seq001_trajectory_3d.png}
    \includegraphics[width=0.31\textwidth]{results_evslam_seq001_velocity_x.png}
    \includegraphics[width=0.31\textwidth]{results_evslam_seq001_velocity_y.png}
    \includegraphics[width=0.31\textwidth]{results_evslam_seq001_velocity_z.png}
    \caption{\texttt{seq001}, całość, konfiguracja końcowa: współrzędne położenia $x$, $y$, $z$, widok 3D oraz składowe prędkości $v_x$, $v_y$ i~$v_z$.}
    \label{fig:evslam_seq001_trajectory}
\end{figure}

\begin{figure}[!ht]
    \centering
    \includegraphics[width=0.40\textwidth]{results_evslam_seq002_trajectory_x.png}
    \includegraphics[width=0.40\textwidth]{results_evslam_seq002_trajectory_y.png}
    \includegraphics[width=0.40\textwidth]{results_evslam_seq002_trajectory_z.png}
    \includegraphics[width=0.40\textwidth]{results_evslam_seq002_trajectory_3d.png}
    \includegraphics[width=0.31\textwidth]{results_evslam_seq002_velocity_x.png}
    \includegraphics[width=0.31\textwidth]{results_evslam_seq002_velocity_y.png}
    \includegraphics[width=0.31\textwidth]{results_evslam_seq002_velocity_z.png}
    \caption{\texttt{seq002}, całość, konfiguracja końcowa: współrzędne położenia $x$, $y$, $z$, widok 3D oraz składowe prędkości $v_x$, $v_y$ i~$v_z$.}
    \label{fig:evslam_seq002_trajectory}
\end{figure}

\begin{figure}[!ht]
    \centering
    \includegraphics[width=0.40\textwidth]{results_evslam_seq003_trajectory_x.png}
    \includegraphics[width=0.40\textwidth]{results_evslam_seq003_trajectory_y.png}
    \includegraphics[width=0.40\textwidth]{results_evslam_seq003_trajectory_z.png}
    \includegraphics[width=0.40\textwidth]{results_evslam_seq003_trajectory_3d.png}
    \includegraphics[width=0.31\textwidth]{results_evslam_seq003_velocity_x.png}
    \includegraphics[width=0.31\textwidth]{results_evslam_seq003_velocity_y.png}
    \includegraphics[width=0.31\textwidth]{results_evslam_seq003_velocity_z.png}
    \caption{\texttt{seq003}, całość, konfiguracja końcowa: współrzędne położenia $x$, $y$, $z$, widok 3D oraz składowe prędkości $v_x$, $v_y$ i~$v_z$.}
    \label{fig:evslam_seq003_trajectory}
\end{figure}

\begin{figure}[!ht]
    \centering
    \includegraphics[width=0.40\textwidth]{results_evslam_seq004_trajectory_x.png}
    \includegraphics[width=0.40\textwidth]{results_evslam_seq004_trajectory_y.png}
    \includegraphics[width=0.40\textwidth]{results_evslam_seq004_trajectory_z.png}
    \includegraphics[width=0.40\textwidth]{results_evslam_seq004_trajectory_3d.png}
    \includegraphics[width=0.31\textwidth]{results_evslam_seq004_velocity_x.png}
    \includegraphics[width=0.31\textwidth]{results_evslam_seq004_velocity_y.png}
    \includegraphics[width=0.31\textwidth]{results_evslam_seq004_velocity_z.png}
    \caption{\texttt{seq004}, całość, konfiguracja końcowa: współrzędne położenia $x$, $y$, $z$, widok 3D oraz składowe prędkości $v_x$, $v_y$ i~$v_z$.}
    \label{fig:evslam_seq004_trajectory}
\end{figure}

\begin{figure}[!ht]
    \centering
    \includegraphics[width=0.40\textwidth]{results_evslam_seq005_trajectory_x.png}
    \includegraphics[width=0.40\textwidth]{results_evslam_seq005_trajectory_y.png}
    \includegraphics[width=0.40\textwidth]{results_evslam_seq005_trajectory_z.png}
    \includegraphics[width=0.40\textwidth]{results_evslam_seq005_trajectory_3d.png}
    \includegraphics[width=0.31\textwidth]{results_evslam_seq005_velocity_x.png}
    \includegraphics[width=0.31\textwidth]{results_evslam_seq005_velocity_y.png}
    \includegraphics[width=0.31\textwidth]{results_evslam_seq005_velocity_z.png}
    \caption{\texttt{seq005}, całość, konfiguracja końcowa: współrzędne położenia $x$, $y$, $z$, widok 3D oraz składowe prędkości $v_x$, $v_y$ i~$v_z$.}
    \label{fig:evslam_seq005_trajectory}
\end{figure}

\begin{figure}[!ht]
    \centering
    \includegraphics[width=0.40\textwidth]{results_evslam_seq006_trajectory_x.png}
    \includegraphics[width=0.40\textwidth]{results_evslam_seq006_trajectory_y.png}
    \includegraphics[width=0.40\textwidth]{results_evslam_seq006_trajectory_z.png}
    \includegraphics[width=0.40\textwidth]{results_evslam_seq006_trajectory_3d.png}
    \includegraphics[width=0.31\textwidth]{results_evslam_seq006_velocity_x.png}
    \includegraphics[width=0.31\textwidth]{results_evslam_seq006_velocity_y.png}
    \includegraphics[width=0.31\textwidth]{results_evslam_seq006_velocity_z.png}
    \caption{\texttt{seq006}, całość, konfiguracja końcowa: współrzędne położenia $x$, $y$, $z$, widok 3D oraz składowe prędkości $v_x$, $v_y$ i~$v_z$.}
    \label{fig:evslam_seq006_trajectory}
\end{figure}

\begin{figure}[!ht]
    \centering
    \includegraphics[width=0.40\textwidth]{results_evslam_seq007_trajectory_x.png}
    \includegraphics[width=0.40\textwidth]{results_evslam_seq007_trajectory_y.png}
    \includegraphics[width=0.40\textwidth]{results_evslam_seq007_trajectory_z.png}
    \includegraphics[width=0.40\textwidth]{results_evslam_seq007_trajectory_3d.png}
    \includegraphics[width=0.31\textwidth]{results_evslam_seq007_velocity_x.png}
    \includegraphics[width=0.31\textwidth]{results_evslam_seq007_velocity_y.png}
    \includegraphics[width=0.31\textwidth]{results_evslam_seq007_velocity_z.png}
    \caption{\texttt{seq007}, całość, konfiguracja końcowa: współrzędne położenia $x$, $y$, $z$, widok 3D oraz składowe prędkości $v_x$, $v_y$ i~$v_z$.}
    \label{fig:evslam_seq007_trajectory}
\end{figure}

Rysunek~\ref{fig:evslam_success_curve} przedstawia krzywe sukcesu, których pola tworzą wartości AUC. Niższy przebieg krzywych \texttt{seq006} i~\texttt{seq007} odpowiada ich niższym wynikom w~tabeli~\ref{tab:wyniki_evslam_final}.

\begin{figure}[!ht]
    \centering
    \includegraphics[width=0.82\textwidth]{results_evslam_official_velocity_success_curve.pdf}
    \caption{Ważone krzywe sukcesu prędkości dla siedmiu sekwencji EvSLAM.}
    \label{fig:evslam_success_curve}
\end{figure}

\clearpage

\subsection{Porównanie z opublikowanymi wynikami}
\label{subsec:porownanie_evslam_literatura}

Tabela~\ref{tab:evslam_sota_comparison} zestawia wyniki systemu z~wybranymi rezultatami opublikowanymi dla tych samych sekwencji. FETAP~\cite{liu2025fetap} korzysta z~jednej kamery zdarzeniowej, kamery klatkowej i~IMU; SuperEvent~\cite{burkhardt2025superevent} oraz EVIV~\cite{lu2025eviv} z~pary kamer zdarzeniowych i~IMU, a~SDEVO~\cite{zhong2025sdevo} wyłącznie z~pary kamer zdarzeniowych. Opracowany system wykorzystuje stereo zdarzeniowe i~żyroskop, dlatego zestawienie pokazuje poziom wyników na wspólnym protokole, lecz nie stanowi porównania metod o~identycznych wejściach~\cite{zhong2026evslam}.

\begin{table}[!ht]
    \centering
    \scriptsize
    \caption{Porównanie ATE RMSE [m] / AUC na podstawie tabeli VII publikacji~\cite{zhong2026evslam}.}
    \label{tab:evslam_sota_comparison}
    \resizebox{\textwidth}{!}{%
        \begin{tabular}{lccccc}
            \toprule
            Sekwencja & Opracowany system & FETAP~\cite{liu2025fetap} & SuperEvent~\cite{burkhardt2025superevent} & SDEVO~\cite{zhong2025sdevo} & EVIV~\cite{lu2025eviv} \\
            \midrule
            \texttt{seq001} & 0,901 / 0,821 & 0,292 / 0,890 & 0,189 / 0,879 & 0,373 / 0,860 & 2,068 / 0,872 \\
            \texttt{seq002} & 0,780 / 0,828 & 0,113 / 0,899 & 0,282 / 0,883 & 0,388 / 0,871 & 1,659 / 0,872 \\
            \texttt{seq003} & 0,333 / 0,852 & 0,116 / 0,930 & 0,080 / 0,904 & 0,418 / 0,905 & 0,885 / 0,879 \\
            \texttt{seq004} & 0,649 / 0,757 & 0,080 / 0,911 & 0,055 / 0,901 & 0,344 / 0,875 & 2,283 / 0,808 \\
            \texttt{seq005} & 0,077 / 0,800 & 0,019 / 0,921 & 0,045 / 0,869 & 0,038 / 0,914 & 0,227 / 0,861 \\
            \texttt{seq006} & 0,111 / 0,506 & 0,048 / 0,831 & 0,033 / 0,857 & 0,069 / 0,855 & 1,966 / 0,805 \\
            \texttt{seq007} & 0,801 / 0,433 & 0,149 / 0,904 & 0,179 / 0,844 & 0,263 / 0,854 & 1,505 / 0,784 \\
            \midrule
            Średnia & 0,522 / 0,714 & 0,117 / 0,898 & 0,123 / 0,877 & 0,271 / 0,876 & 1,513 / 0,840 \\
            \bottomrule
        \end{tabular}%
    }
\end{table}

Średnia wartość ATE opracowanego systemu jest mniejsza niż dla EVIV, ale większa niż dla FETAP, SuperEvent i~SDEVO. Średnia wartość AUC jest niższa od wszystkich zestawionych metod. Największa różnica AUC występuje dla \texttt{seq007}; bez dodatkowej diagnostyki nie przesądza to jednak o~jej przyczynie ani o~uogólnieniu wyniku na cały rodzaj ruchu.

W~tabeli~\ref{tab:evslam_competition_results} zestawiono wynik opracowanego systemu z~końcowym rankingiem konkursu EvSLAM. Wszystkie wartości dotyczą tych samych siedmiu sekwencji i~zostały obliczone według wspólnego protokołu ewaluacyjnego. Suma AUC ocenia dokładność wyznaczania prędkości, przy czym większa wartość oznacza lepszy wynik, natomiast suma ATE opisuje łączny błąd położenia i~powinna być jak najmniejsza. Miejsca w tabeli w pierwszej kolejności wyznacza wartość ATE. AUC jest porównywane dopiero w przypadku remisu.

\newpage

Rdzeń opracowanego systemu łączy przetwarzanie klatek zdarzeniowych i~śledzenie punktów inspirowane zdarzeniowym VIO Rebecqa~\cite{rebecq2017eventvio} oraz UltimateSLAM~\cite{vidal2018ultimate}, strukturę stereo i~mapy lokalnej wzorowaną na S-PTAM~\cite{pire2017sptam}, a~także rozpoznawanie miejsc i~domykanie pętli inspirowane ORB-SLAM2~\cite{murartal2017orbslam2}. Wynik przedstawiony w~tabeli pozwala ocenić takie połączenie klasycznych metod geometrycznych i~reprezentacji
danych zdarzeniowych na tle pozostałych rozwiązań konkursowych.

\begin{table}[!ht]
    \centering
    \small
    \caption{Końcowe wyniki konkursu EvSLAM i~wynik opracowanego systemu~\cite{evslam2025challenge}.}
    \label{tab:evslam_competition_results}
    \begin{tabularx}{0.82\textwidth}{c X r r}
        \toprule
        Miejsce & Metoda & Suma AUC $\uparrow$ & Suma ATE [m] $\downarrow$ \\
        \midrule
        1 & SuperPoint + FETAP + VINS & 6,286 & 0,817 \\
        2 & OKVIS2 + SuperEvent & 6,271 & 0,448 \\
        3 & OKVIS2 + SuperEvent bez obrazów & 6,137 & 0,863 \\
        4 & SDEVO & 6,133 & 1,894 \\
        5 & REFIO & 6,083 & 1,344 \\
        6 & DEIO + LC & 6,048 & 0,464 \\
        7 & EVIV & 5,881 & 10,594 \\
        -- & Opracowany system & 4,996 & 3,654 \\
        8 & \texttt{test-0829} & 2,829 & 243,636 \\
        \bottomrule
    \end{tabularx}
\end{table}

Największą sumę AUC uzyskało rozwiązanie łączące SuperPoint, FETAP i~VINS, natomiast najmniejszą sumę ATE osiągnęło połączenie OKVIS2 z~SuperEvent. Wyniki wariantów SuperEvent pokazują, że wykorzystanie obrazów z~kamer klatkowych zwiększyło sumę AUC z~$6{,}137$ do $6{,}271$, a~sumę ATE zmniejszyło z~$0{,}863,\mathrm{m}$ do $0{,}448,\mathrm{m}$. Kolejność rozwiązań zależy zatem od rozpatrywanej miary: metoda najlepsza pod względem estymacji prędkości nie uzyskała najmniejszego błędu położenia.

Opracowany system uzyskał sumę AUC równą $4{,}996$, czyli wynik niższy od siedmiu najwyżej sklasyfikowanych metod, ale wyższy od ostatniego zgłoszenia w~rankingu. W~porównaniu ze zwycięskim rozwiązaniem suma AUC była niższa o~$1{,}290$, czyli o~$20{,}5\%$, natomiast suma ATE była większa o~$2{,}837,\mathrm{m}$, a więc ponad trzykrotnie. Jednocześnie suma ATE opracowanego systemu była mniejsza niż dla EVIV oraz ostatniego zgłoszenia. W~porównaniu z~EVIV system dokładniej odtworzył położenie, ale uzyskał słabszy wynik estymacji prędkości. Potwierdza to, że wyznaczanie prędkości jest elementem wymagającym dalszego rozwoju.

Wartości dla opracowanego systemu obliczono za pomocą oficjalnego narzędzia ewaluacyjnego. Jego wiersz w tabeli \ref{tab:evslam_competition_results} nie jest więc oficjalnym zgłoszeniem konkursowym, a~położenia względem wyników rankingu nie należy traktować jako przyznanego miejsca.

\clearpage

\section{Wyniki dla zbioru M3ED}
\label{sec:wyniki_m3ed}

\subsection{Sekwencje z dostępną trajektorią referencyjną}
\label{subsec:wyniki_m3ed_gt}

Estymowaną trajektorię i~prędkość w zestawieniu z trajektorią referencyjną dla sekwencji \texttt{falcon\_indoor\_flight\_1} przedstawiono na rysunku~\ref{fig:wyniki_m3ed_indoor_trajectory}.

\begin{table}[!ht]
    \centering
    \small
    \caption{Wyniki pełnych przebiegów dwóch sekwencji M3ED z~dostępną trajektorią referencyjną, w~protokołach określonych w~metodyce.}
    \label{tab:wyniki_m3ed_gt}
    \resizebox{\textwidth}{!}{%
        \begin{tabular}{lrlrrrrr}
            \toprule
            Sekwencja & Okno [ms] & Protokół & ATE/APE RMSE [m] & RMSE rotacji [$^\circ$] & AUC & RPE translacji [m] & RPE rotacji [$^\circ$] \\
            \midrule
            \texttt{falcon\_indoor\_flight\_1} & 24 & wspólny $SE(3)$ & 0,162 & 4,434 & 0,830 & 0,092 & 0,514 \\
            \texttt{falcon\_indoor\_flight\_1} & 24 & oficjalny $Sim(3)$ & 0,161 & -- & -- & -- & -- \\
            \texttt{falcon\_outdoor\_day\_fast\_flight\_2} & 24 & wspólny $SE(3)$ & 25,305 & 47,762 & 0,605 & 2,406 & 4,072 \\
            \texttt{falcon\_outdoor\_day\_fast\_flight\_2} & 24 & oficjalny $Sim(3)$ & 12,403 & -- & -- & -- & -- \\
            \bottomrule
        \end{tabular}%
    }
\end{table}

\begin{figure}[!ht]
    \centering
    \includegraphics[width=0.40\textwidth]{results_m3ed_falcon_indoor_flight_1_first_pose_trajectory_x.png}
    \includegraphics[width=0.40\textwidth]{results_m3ed_falcon_indoor_flight_1_first_pose_trajectory_y.png}
    \includegraphics[width=0.40\textwidth]{results_m3ed_falcon_indoor_flight_1_first_pose_trajectory_z.png}
    \includegraphics[width=0.40\textwidth]{results_m3ed_falcon_indoor_flight_1_first_pose_trajectory_3d.png}
    \includegraphics[width=0.31\textwidth]{results_m3ed_falcon_indoor_flight_1_first_pose_velocity_x.png}
    \includegraphics[width=0.31\textwidth]{results_m3ed_falcon_indoor_flight_1_first_pose_velocity_y.png}
    \includegraphics[width=0.31\textwidth]{results_m3ed_falcon_indoor_flight_1_first_pose_velocity_z.png}
    \caption{Pełna sekwencja \texttt{falcon\_indoor\_flight\_1}: współrzędne położenia $x$, $y$, $z$, widok 3D oraz składowe prędkości $v_x$, $v_y$ i~$v_z$.}
    \label{fig:wyniki_m3ed_indoor_trajectory}
\end{figure}

Graf przebiegu zawiera jedną zaakceptowaną krawędź pętli. Otrzymano pełne pokrycie trajektorii referencyjnej, a~zbliżone wartości ATE po wyrównaniu $SE(3)$ i~APE po $Sim(3)$ wskazują, że dodatkowa korekta skali ma niewielki wpływ na ten wynik.

\begin{figure}[!ht]
    \centering
    \includegraphics[width=0.40\textwidth]{results_m3ed_falcon_outdoor_day_fast_flight_2_first_pose_trajectory_x.png}
    \includegraphics[width=0.40\textwidth]{results_m3ed_falcon_outdoor_day_fast_flight_2_first_pose_trajectory_y.png}
    \includegraphics[width=0.40\textwidth]{results_m3ed_falcon_outdoor_day_fast_flight_2_first_pose_trajectory_z.png}
    \includegraphics[width=0.40\textwidth]{results_m3ed_falcon_outdoor_day_fast_flight_2_first_pose_trajectory_3d.png}
    \includegraphics[width=0.31\textwidth]{results_m3ed_falcon_outdoor_day_fast_flight_2_first_pose_velocity_x.png}
    \includegraphics[width=0.31\textwidth]{results_m3ed_falcon_outdoor_day_fast_flight_2_first_pose_velocity_y.png}
    \includegraphics[width=0.31\textwidth]{results_m3ed_falcon_outdoor_day_fast_flight_2_first_pose_velocity_z.png}
    \caption{Pełna sekwencja \texttt{falcon\_outdoor\_day\_fast\_flight\_2} z~oknem 24~ms: współrzędne położenia $x$, $y$, $z$, widok 3D oraz składowe prędkości $v_x$, $v_y$ i~$v_z$.}
    \label{fig:wyniki_m3ed_fast_flight_2_trajectory}
\end{figure}

W~sekwencji \texttt{falcon\_outdoor\_day\_fast\_flight\_2} system przetworzył cały przebieg, jednak trajektoria wykazuje znaczny dryf. Potwierdzają to wartości ATE i~APE oraz rozbieżności przebiegów widoczne na rysunku~\ref{fig:wyniki_m3ed_fast_flight_2_trajectory}. Tak słaby wynik pokazuje ograniczenia systemu w~sekwencji zawierającej niewiele charakterystycznych punktów odniesienia, położonych w~dużej odległości od kamer, co znacznie utrudnia poprawne oszacowanie głębi. Dodatkowym utrudnieniem jest brak warunków do wykrycia pętli na trasie drona. Wraca on w~mało charakterystyczne miejsce początkowe dopiero podczas lądowania i~porusza się wtedy bardzo wolno, co ogranicza liczbę generowanych zdarzeń oraz uniemożliwia wykrycie pętli i~korektę trajektorii.

\begin{table}[!ht]
    \centering
    \small
    \caption{Końcowe wyniki konkursu M3ED oraz wyniki opracowanego systemu według tej samej metryki APE RMSE po wyrównaniu $Sim(3)$~\cite{m3ed2025challenge,m3ed2025leaderboard}. W~oficjalnej tabeli wyników konkursu znajdują się tylko dwa zgłoszenia.}
    \label{tab:m3ed_competition_results}
    \begin{tabularx}{\textwidth}{X X r}
        \toprule
        Metoda & Zakres wyniku & APE RMSE [m] $\downarrow$ \\
        \midrule
        NAIL\_SLAM -- Stereo-DEVO & suma dla trzech ukrytych sekwencji konkursowych & 4,84 \\
        lziupenn -- reconstruction & suma dla trzech ukrytych sekwencji konkursowych & 47,49 \\
        Opracowany system & \path{falcon_outdoor_day_fast_flight_2}, jedna jawna sekwencja & 12,403 \\
        Opracowany system & \path{falcon_indoor_flight_1}, jedna jawna sekwencja & 0,161 \\
        \bottomrule
    \end{tabularx}
\end{table}

Wyniki konkursowe przedstawione w~tabeli~\ref{tab:m3ed_competition_results} są sumą błędów trzech innych ukrytych sekwencji, natomiast wartości opracowanego systemu dotyczą pojedynczych sekwencji jawnych. Na ich podstawie trudno oszacować, jaki wynik opracowany system uzyskałby w~konkursie M3ED. Wyniki dla dwóch różnych sekwencji referencyjnych pokazują, że jakość estymacji jest silnie uzależniona od charakterystyki konkretnej sekwencji testowej.

\subsection{Sekwencje bez dostępnej trajektorii referencyjnej}
\label{subsec:wyniki_m3ed_no_gt}

Do prezentacji jakościowej wybrano pełną sekwencję \texttt{spot\_loop}. Jej trajektoria referencyjna jest ukryta, dlatego tabela~\ref{tab:wyniki_m3ed_no_gt} zawiera wyłącznie informacje o~kompletności oraz diagnostykę zamkniętej trasy. Nie są to miary dokładności.

\begin{table}[!ht]
    \centering
    \small
    \caption{Diagnostyka pełnego przebiegu M3ED \texttt{spot\_loop}.}
    \label{tab:wyniki_m3ed_no_gt}
    \begin{tabularx}{0.78\textwidth}{X r}
        \toprule
        Wielkość & Wartość \\
        \midrule
        status przebiegu & ukończony \\
        czas przetwarzania & 6741~s (1~h 52 min 21~s) \\
        liczba zapisanych póz & 2960 \\
        czas trwania trajektorii & 295,893~s \\
        oszacowana długość trasy & 292,859~m \\
        odległość start--koniec & 22,829~m \\
        odległość start--koniec / długość trasy & 7,80\% \\
        zaakceptowane domknięcia pętli & 0 \\
        \bottomrule
    \end{tabularx}
\end{table}

\begin{figure}[!ht]
    \centering
    \includegraphics[width=0.47\textwidth]{results_m3ed_spot_outdoor_day_penno_building_loop_trajectory_x.png}
    \includegraphics[width=0.47\textwidth]{results_m3ed_spot_outdoor_day_penno_building_loop_trajectory_y.png}
    \includegraphics[width=0.47\textwidth]{results_m3ed_spot_outdoor_day_penno_building_loop_trajectory_z.png}
    \includegraphics[width=0.47\textwidth]{results_m3ed_spot_outdoor_day_penno_building_loop_trajectory_3d.png}
    \caption{Estymowana trajektoria pełnej sekwencji M3ED \texttt{spot\_loop}: współrzędne $x$, $y$, $z$ i~widok 3D. Przedstawion wyłącznie trajektorię obliczoną przez opracowany system, ponieważ trajektoria referencyjna dla tej sekwencji nie jest dostępna.}
    \label{fig:wyniki_m3ed_no_gt_trajectory}
\end{figure}

Przebieg objął całe 295,9~s danych i~utworzył 2\,212 klatek kluczowych oraz 33\,554 punkty mapy. Trajektorię pokazano na rysunku~\ref{fig:wyniki_m3ed_no_gt_trajectory}. System nie zaakceptował domknięcia pętli, a~końcowa odległość od punktu startu odpowiada 7,80\% długości estymowanej trasy. Wynik potwierdza wykonanie pełnego przebiegu, ale bez GT nie pozwala określić błędu trajektorii.

\clearpage

\section{Ocena działania systemu}
\label{sec:wnioski}

\subsection{Wyniki końcowe}

Dla siedmiu sekwencji EvSLAM uzyskano średnie pokrycie trajektorii referencyjnej równe $99{,}6\%$, średnie ATE RMSE wynoszące $0{,}522\,\mathrm{m}$ i~średnie AUC równe $0{,}714$. Najmniejszy błąd położenia wystąpił dla ruchu liniowego w~\texttt{seq005} i~wyniósł $0{,}077\,\mathrm{m}$, natomiast największe błędy uzyskano w~\texttt{seq001} i~\texttt{seq007}. W~porównaniu z~metodami zestawionymi w~tabeli~\ref{tab:evslam_sota_comparison} średnie ATE było mniejsze niż dla EVIV, ale większe niż dla FETAP, SuperEvent* i~SDEVO, przy czym metody te korzystają z~różnych zestawów czujników. Niższe AUC wskazuje, że wyznaczanie prędkości pozostaje słabszym elementem rozwiązania. Pełne wyniki zawiera tabela~\ref{tab:wyniki_evslam_final}.

Wyniki M3ED zależały od scenariusza. Dla \texttt{falcon\_indoor\_flight\_1} uzyskano ATE RMSE równe $0{,}162\,\mathrm{m}$ po wyrównaniu $SE(3)$. Dla zarejestrowanej na zewnątrz sekwencji \texttt{falcon\_outdoor\_day\_fast\_flight\_2} wartości ATE RMSE i~APE RMSE wyniosły odpowiednio $25{,}305\,\mathrm{m}$ i~$12{,}403\,\mathrm{m}$, co potwierdza znaczny dryf podczas szybkiego lotu z~odległymi punktami odniesienia. Wyniki obu przebiegów przedstawiono w~tabeli~\ref{tab:wyniki_m3ed_gt}. System przetworzył także całe $295{,}9\,\mathrm{s}$ sekwencji \texttt{spot\_loop}, lecz nie wykrył w~niej pętli. Ze względu na brak trajektorii referencyjnej wynik ten potwierdza jedynie kompletność przebiegu, a~nie jego dokładność.

\subsection{Wpływ badanych parametrów na wynik}

Na wspólnym fragmencie \texttt{seq007\_test} zastosowanie mapy i~korekty grafu zmniejszyło ATE RMSE z~$0{,}544\,\mathrm{m}$ do $0{,}151\,\mathrm{m}$, a~RMSE rotacji z~$16{,}256^\circ$ do $3{,}179^\circ$. Nie poprawiły się AUC ani translacyjny RPE, co oznacza, że warstwa SLAM ograniczała przede wszystkim globalny dryf pozycji~(tabela~\ref{tab:wyniki_slam_vo}). Dla \texttt{falcon\_indoor\_flight\_1} przyjęcie rotacji z~żyroskopu zamiast pełnej estymacji PnP zmniejszyło ATE z~$0{,}240\,\mathrm{m}$ do $0{,}162\,\mathrm{m}$ i~RMSE rotacji z~$9{,}885^\circ$ do $4{,}434^\circ$, co uzasadnia wybór tego wariantu w~konfiguracji końcowej~(tabela~\ref{tab:wyniki_rotation_source}).

W~badanych fragmentach filtr BAF poprawiał część miar, ale jego koszt zależał od gęstości strumienia zdarzeń. Wydłużył czas badanego fragmentu EvSLAM o~$7{,}2\%$, a~M3ED o~$207{,}5\%$, dlatego dla M3ED został wyłączony~(tabela~\ref{tab:wyniki_baf}). Korzystny czas akumulacji również zależał od danych: dla szybkiego ruchu ramienia EvSLAM najlepsze ATE uzyskano przy $6\,\mathrm{ms}$, a~najwyższe AUC przy $12\,\mathrm{ms}$, natomiast na badanym fragmencie M3ED większość miar poprawiała się do $48\,\mathrm{ms}$. W~konfiguracji końcowej przyjęto odpowiednio $12\,\mathrm{ms}$ i~$24\,\mathrm{ms}$~(tabela~\ref{tab:wyniki_accumulation}).

\subsection{Ograniczenia}

Największy dryf wystąpił w~zewnętrznej sekwencji M3ED, w~której dominowały odległe elementy sceny. Mała dysparycja odpowiadająca takim punktom zwiększała niepewność głębokości, a~tym samym również niepewność translacji wyznaczanej na podstawie punktów 3D. W~tym przebiegu oraz w~sekwencji \texttt{spot\_loop} nie wykryto żadnej pętli, mimo powrotu robota w~okolice początku drugiej z~nich. Wskazuje to na ograniczoną odporność rozpoznawania miejsc na zmianę punktu widzenia i~niewielką aktywność zdarzeń.

Ograniczenia ujawniły się również w~sekwencjach EvSLAM. Podczas postoju lub wolnego ruchu mała liczba zdarzeń utrudniała wyznaczenie punktów 3D i~utworzenie wystarczającej liczby korespondencji PnP. Na początku kilku sekwencji opóźniało to rozpoczęcie poprawnej estymacji i~powodowało błąd początkowego fragmentu trajektorii, na przykład w~\texttt{seq006}.

W~\texttt{seq007} zaobserwowano ponadto wyraźne zaniżenie skali translacji, widoczne na rysunku~\ref{fig:evslam_seq007_trajectory}. Prawdopodobną przyczyną jest wrażliwość estymacji na błędy dopasowania stereo: zawyżona dysparycja prowadzi do zaniżenia głębokości punktów.

System nie estymuje błędów systematycznych IMU ani prędkości jako części stanu. Bezpośrednie wykorzystanie żyroskopu poprawiło wynik kontrolowanego porównania, ale może przenosić jego dryf do trajektorii. Ograniczeniem pozostaje również wydajność. Nie oceniano osobno czasu wykonania poszczególnych sekwencji, jednak system nie jest zoptymalizowany pod kątem wydajności obliczeniowej, a~przetwarzanie obszernych sekwencji trwa stosunkowo długo.

    \chapter{Podsumowanie}
\label{cha:podsumowanie}

% \section{Realizacja celu pracy}
% \label{sec:realizacja_celu}

Celem pracy było opracowanie systemu wyznaczającego ruch robota na podstawie danych z~dwóch kamer zdarzeniowych, przeznaczonego do przetwarzania zbiorów EvSLAM i~M3ED. Cel ten został zrealizowany. Powstał system zawierający moduły odpowiedzialne za synchronizację i~przygotowanie zdarzeń, śledzenie punktów, geometrię stereo oraz estymację trajektorii lewej kamery. Dla EvSLAM wyznaczana jest również wymagana przez konkurs prędkość liniowa.

Odometria wizyjna została rozszerzona o~rzadką mapę, lokalizację względem jej punktów, relokalizację, domykanie pętli i~optymalizację grafu pozy. System przetworzył wszystkie oceniane sekwencje EvSLAM oraz wybrane pełne przebiegi M3ED. Działa on na gotowych, zebranych wcześniej sekwencjach i~obecnie nie jest przystosowany do pracy w~czasie rzeczywistym na robocie mobilnym. Wykorzystuje IMU, lecz używa wyłącznie żyroskopu, dlatego nie stanowi pełnego estymatora wizyjno-inercyjnego. Prędkość oblicza z~gotowej trajektorii położenia, co jest prostym rozwiązaniem stosowanym na potrzeby uzyskania wyników wymaganych przez EvSLAM.

\section{Kierunki dalszych prac}
\label{sec:dalsze_prace}

Najważniejsze kierunki rozwoju, w~kolejności wynikającej z~wykazanych ograniczeń, obejmują:

\begin{enumerate}
    \item utworzenie pełnego estymatora wizyjno-inercyjnego, który wspólnie wyznaczałby pozę, prędkość oraz błędy systematyczne IMU z~wykorzystaniem preintegracji pomiarów~\cite{forster2015preintegration} -- mogłoby to ograniczyć dryf rotacji, będący jednym z~najważniejszych ograniczeń widocznych w~wynikach dla sekwencji \texttt{falcon\_outdoor\_day\_fast\_flight\_2};
    \item uwzględnienie niepewności estymacji głębi podczas wyboru punktów do estymacji pozy;
    \item wprowadzenie adaptacyjnego doboru czasu akumulacji zależnego od dynamiki sceny, który zwiększyłby odporność systemu na zmiany prędkości ruchu robota i~pozwoliłby skutecznie estymować pozę z~większej liczby klatek zdarzeniowych;
    \item dodanie lokalnej optymalizacji wiązki (ang. \textit{bundle adjustment}), stosowanej w~bardziej rozbudowanych systemach SLAM, między innymi w~ORB-SLAM2~\cite{murartal2017orbslam2}, oraz zastosowanie bardziej wydajnego optymalizatora grafowego dla większych map~\cite{grisetti2010graph};
    \item ocena złożoności obliczeniowej i~przyspieszenie najbardziej kosztownych etapów, a~następnie sprawdzenie systemu na bieżącym strumieniu danych.
\end{enumerate}

    
    % itd.
    % \appendix
    % \include{dodatekA}
    % \include{dodatekB}
    % itd.
    \printbibliography

    Kod projektu oraz inne pliki są dostępne w~repozytorium: 
    
    \url{https://github.com/romannowak9/dvs-slam-agh/event-slam}

\end{document}
